\documentclass{article} % For LaTeX2e
\usepackage{iclr2021_conference,times}

% Optional math commands from https://github.com/goodfeli/dlbook_notation.
%%%%% NEW MATH DEFINITIONS %%%%%

\usepackage{amsmath,amsfonts,bm}

% Mark sections of captions for referring to divisions of figures
\newcommand{\figleft}{{\em (Left)}}
\newcommand{\figcenter}{{\em (Center)}}
\newcommand{\figright}{{\em (Right)}}
\newcommand{\figtop}{{\em (Top)}}
\newcommand{\figbottom}{{\em (Bottom)}}
\newcommand{\captiona}{{\em (a)}}
\newcommand{\captionb}{{\em (b)}}
\newcommand{\captionc}{{\em (c)}}
\newcommand{\captiond}{{\em (d)}}

% Highlight a newly defined term
\newcommand{\newterm}[1]{{\bf #1}}


% Figure reference, lower-case.
\def\figref#1{figure~\ref{#1}}
% Figure reference, capital. For start of sentence
\def\Figref#1{Figure~\ref{#1}}
\def\twofigref#1#2{figures \ref{#1} and \ref{#2}}
\def\quadfigref#1#2#3#4{figures \ref{#1}, \ref{#2}, \ref{#3} and \ref{#4}}
% Section reference, lower-case.
\def\secref#1{section~\ref{#1}}
% Section reference, capital.
\def\Secref#1{Section~\ref{#1}}
% Reference to two sections.
\def\twosecrefs#1#2{sections \ref{#1} and \ref{#2}}
% Reference to three sections.
\def\secrefs#1#2#3{sections \ref{#1}, \ref{#2} and \ref{#3}}
% Reference to an equation, lower-case.
%\def\eqref#1{equation~\ref{#1}}
% Reference to an equation, upper case
%\def\Eqref#1{Equation~\ref{#1}}
% A raw reference to an equation---avoid using if possible
\def\plaineqref#1{\ref{#1}}
% Reference to a chapter, lower-case.
\def\chapref#1{chapter~\ref{#1}}
% Reference to an equation, upper case.
\def\Chapref#1{Chapter~\ref{#1}}
% Reference to a range of chapters
\def\rangechapref#1#2{chapters\ref{#1}--\ref{#2}}
% Reference to an algorithm, lower-case.
\def\algref#1{algorithm~\ref{#1}}
% Reference to an algorithm, upper case.
\def\Algref#1{Algorithm~\ref{#1}}
\def\twoalgref#1#2{algorithms \ref{#1} and \ref{#2}}
\def\Twoalgref#1#2{Algorithms \ref{#1} and \ref{#2}}
% Reference to a part, lower case
\def\partref#1{part~\ref{#1}}
% Reference to a part, upper case
\def\Partref#1{Part~\ref{#1}}
\def\twopartref#1#2{parts \ref{#1} and \ref{#2}}

\def\ceil#1{\lceil #1 \rceil}
\def\floor#1{\lfloor #1 \rfloor}
\def\1{\bm{1}}
\newcommand{\train}{\mathcal{D}}
\newcommand{\valid}{\mathcal{D_{\mathrm{valid}}}}
\newcommand{\test}{\mathcal{D_{\mathrm{test}}}}

\def\eps{{\epsilon}}


% Random variables
\def\reta{{\textnormal{$\eta$}}}
\def\ra{{\textnormal{a}}}
\def\rb{{\textnormal{b}}}
\def\rc{{\textnormal{c}}}
\def\rd{{\textnormal{d}}}
\def\re{{\textnormal{e}}}
\def\rf{{\textnormal{f}}}
\def\rg{{\textnormal{g}}}
\def\rh{{\textnormal{h}}}
\def\ri{{\textnormal{i}}}
\def\rj{{\textnormal{j}}}
\def\rk{{\textnormal{k}}}
\def\rl{{\textnormal{l}}}
% rm is already a command, just don't name any random variables m
\def\rn{{\textnormal{n}}}
\def\ro{{\textnormal{o}}}
\def\rp{{\textnormal{p}}}
\def\rq{{\textnormal{q}}}
\def\rr{{\textnormal{r}}}
\def\rs{{\textnormal{s}}}
\def\rt{{\textnormal{t}}}
\def\ru{{\textnormal{u}}}
\def\rv{{\textnormal{v}}}
\def\rw{{\textnormal{w}}}
\def\rx{{\textnormal{x}}}
\def\ry{{\textnormal{y}}}
\def\rz{{\textnormal{z}}}

% Random vectors
\def\rvepsilon{{\mathbf{\epsilon}}}
\def\rvtheta{{\mathbf{\theta}}}
\def\rva{{\mathbf{a}}}
\def\rvb{{\mathbf{b}}}
\def\rvc{{\mathbf{c}}}
\def\rvd{{\mathbf{d}}}
\def\rve{{\mathbf{e}}}
\def\rvf{{\mathbf{f}}}
\def\rvg{{\mathbf{g}}}
\def\rvh{{\mathbf{h}}}
\def\rvu{{\mathbf{i}}}
\def\rvj{{\mathbf{j}}}
\def\rvk{{\mathbf{k}}}
\def\rvl{{\mathbf{l}}}
\def\rvm{{\mathbf{m}}}
\def\rvn{{\mathbf{n}}}
\def\rvo{{\mathbf{o}}}
\def\rvp{{\mathbf{p}}}
\def\rvq{{\mathbf{q}}}
\def\rvr{{\mathbf{r}}}
\def\rvs{{\mathbf{s}}}
\def\rvt{{\mathbf{t}}}
\def\rvu{{\mathbf{u}}}
\def\rvv{{\mathbf{v}}}
\def\rvw{{\mathbf{w}}}
\def\rvx{{\mathbf{x}}}
\def\rvy{{\mathbf{y}}}
\def\rvz{{\mathbf{z}}}

% Elements of random vectors
\def\erva{{\textnormal{a}}}
\def\ervb{{\textnormal{b}}}
\def\ervc{{\textnormal{c}}}
\def\ervd{{\textnormal{d}}}
\def\erve{{\textnormal{e}}}
\def\ervf{{\textnormal{f}}}
\def\ervg{{\textnormal{g}}}
\def\ervh{{\textnormal{h}}}
\def\ervi{{\textnormal{i}}}
\def\ervj{{\textnormal{j}}}
\def\ervk{{\textnormal{k}}}
\def\ervl{{\textnormal{l}}}
\def\ervm{{\textnormal{m}}}
\def\ervn{{\textnormal{n}}}
\def\ervo{{\textnormal{o}}}
\def\ervp{{\textnormal{p}}}
\def\ervq{{\textnormal{q}}}
\def\ervr{{\textnormal{r}}}
\def\ervs{{\textnormal{s}}}
\def\ervt{{\textnormal{t}}}
\def\ervu{{\textnormal{u}}}
\def\ervv{{\textnormal{v}}}
\def\ervw{{\textnormal{w}}}
\def\ervx{{\textnormal{x}}}
\def\ervy{{\textnormal{y}}}
\def\ervz{{\textnormal{z}}}

% Random matrices
\def\rmA{{\mathbf{A}}}
\def\rmB{{\mathbf{B}}}
\def\rmC{{\mathbf{C}}}
\def\rmD{{\mathbf{D}}}
\def\rmE{{\mathbf{E}}}
\def\rmF{{\mathbf{F}}}
\def\rmG{{\mathbf{G}}}
\def\rmH{{\mathbf{H}}}
\def\rmI{{\mathbf{I}}}
\def\rmJ{{\mathbf{J}}}
\def\rmK{{\mathbf{K}}}
\def\rmL{{\mathbf{L}}}
\def\rmM{{\mathbf{M}}}
\def\rmN{{\mathbf{N}}}
\def\rmO{{\mathbf{O}}}
\def\rmP{{\mathbf{P}}}
\def\rmQ{{\mathbf{Q}}}
\def\rmR{{\mathbf{R}}}
\def\rmS{{\mathbf{S}}}
\def\rmT{{\mathbf{T}}}
\def\rmU{{\mathbf{U}}}
\def\rmV{{\mathbf{V}}}
\def\rmW{{\mathbf{W}}}
\def\rmX{{\mathbf{X}}}
\def\rmY{{\mathbf{Y}}}
\def\rmZ{{\mathbf{Z}}}

% Elements of random matrices
\def\ermA{{\textnormal{A}}}
\def\ermB{{\textnormal{B}}}
\def\ermC{{\textnormal{C}}}
\def\ermD{{\textnormal{D}}}
\def\ermE{{\textnormal{E}}}
\def\ermF{{\textnormal{F}}}
\def\ermG{{\textnormal{G}}}
\def\ermH{{\textnormal{H}}}
\def\ermI{{\textnormal{I}}}
\def\ermJ{{\textnormal{J}}}
\def\ermK{{\textnormal{K}}}
\def\ermL{{\textnormal{L}}}
\def\ermM{{\textnormal{M}}}
\def\ermN{{\textnormal{N}}}
\def\ermO{{\textnormal{O}}}
\def\ermP{{\textnormal{P}}}
\def\ermQ{{\textnormal{Q}}}
\def\ermR{{\textnormal{R}}}
\def\ermS{{\textnormal{S}}}
\def\ermT{{\textnormal{T}}}
\def\ermU{{\textnormal{U}}}
\def\ermV{{\textnormal{V}}}
\def\ermW{{\textnormal{W}}}
\def\ermX{{\textnormal{X}}}
\def\ermY{{\textnormal{Y}}}
\def\ermZ{{\textnormal{Z}}}

% Vectors
\def\vzero{{\bm{0}}}
\def\vone{{\bm{1}}}
\def\vmu{{\bm{\mu}}}
\def\vtheta{{\bm{\theta}}}
\def\va{{\bm{a}}}
\def\vb{{\bm{b}}}
\def\vc{{\bm{c}}}
\def\vd{{\bm{d}}}
\def\ve{{\bm{e}}}
\def\vf{{\bm{f}}}
\def\vg{{\bm{g}}}
\def\vh{{\bm{h}}}
\def\vi{{\bm{i}}}
\def\vj{{\bm{j}}}
\def\vk{{\bm{k}}}
\def\vl{{\bm{l}}}
\def\vm{{\bm{m}}}
\def\vn{{\bm{n}}}
\def\vo{{\bm{o}}}
\def\vp{{\bm{p}}}
\def\vq{{\bm{q}}}
\def\vr{{\bm{r}}}
\def\vs{{\bm{s}}}
\def\vt{{\bm{t}}}
\def\vu{{\bm{u}}}
\def\vv{{\bm{v}}}
\def\vw{{\bm{w}}}
\def\vx{{\bm{x}}}
\def\vy{{\bm{y}}}
\def\vz{{\bm{z}}}

% Elements of vectors
\def\evalpha{{\alpha}}
\def\evbeta{{\beta}}
\def\evepsilon{{\epsilon}}
\def\evlambda{{\lambda}}
\def\evomega{{\omega}}
\def\evmu{{\mu}}
\def\evpsi{{\psi}}
\def\evsigma{{\sigma}}
\def\evtheta{{\theta}}
\def\eva{{a}}
\def\evb{{b}}
\def\evc{{c}}
\def\evd{{d}}
\def\eve{{e}}
\def\evf{{f}}
\def\evg{{g}}
\def\evh{{h}}
\def\evi{{i}}
\def\evj{{j}}
\def\evk{{k}}
\def\evl{{l}}
\def\evm{{m}}
\def\evn{{n}}
\def\evo{{o}}
\def\evp{{p}}
\def\evq{{q}}
\def\evr{{r}}
\def\evs{{s}}
\def\evt{{t}}
\def\evu{{u}}
\def\evv{{v}}
\def\evw{{w}}
\def\evx{{x}}
\def\evy{{y}}
\def\evz{{z}}

% Matrix
\def\mA{{\bm{A}}}
\def\mB{{\bm{B}}}
\def\mC{{\bm{C}}}
\def\mD{{\bm{D}}}
\def\mE{{\bm{E}}}
\def\mF{{\bm{F}}}
\def\mG{{\bm{G}}}
\def\mH{{\bm{H}}}
\def\mI{{\bm{I}}}
\def\mJ{{\bm{J}}}
\def\mK{{\bm{K}}}
\def\mL{{\bm{L}}}
\def\mM{{\bm{M}}}
\def\mN{{\bm{N}}}
\def\mO{{\bm{O}}}
\def\mP{{\bm{P}}}
\def\mQ{{\bm{Q}}}
\def\mR{{\bm{R}}}
\def\mS{{\bm{S}}}
\def\mT{{\bm{T}}}
\def\mU{{\bm{U}}}
\def\mV{{\bm{V}}}
\def\mW{{\bm{W}}}
\def\mX{{\bm{X}}}
\def\mY{{\bm{Y}}}
\def\mZ{{\bm{Z}}}
\def\mBeta{{\bm{\beta}}}
\def\mPhi{{\bm{\Phi}}}
\def\mLambda{{\bm{\Lambda}}}
\def\mSigma{{\bm{\Sigma}}}

% Tensor
\DeclareMathAlphabet{\mathsfit}{\encodingdefault}{\sfdefault}{m}{sl}
\SetMathAlphabet{\mathsfit}{bold}{\encodingdefault}{\sfdefault}{bx}{n}
\newcommand{\tens}[1]{\bm{\mathsfit{#1}}}
\def\tA{{\tens{A}}}
\def\tB{{\tens{B}}}
\def\tC{{\tens{C}}}
\def\tD{{\tens{D}}}
\def\tE{{\tens{E}}}
\def\tF{{\tens{F}}}
\def\tG{{\tens{G}}}
\def\tH{{\tens{H}}}
\def\tI{{\tens{I}}}
\def\tJ{{\tens{J}}}
\def\tK{{\tens{K}}}
\def\tL{{\tens{L}}}
\def\tM{{\tens{M}}}
\def\tN{{\tens{N}}}
\def\tO{{\tens{O}}}
\def\tP{{\tens{P}}}
\def\tQ{{\tens{Q}}}
\def\tR{{\tens{R}}}
\def\tS{{\tens{S}}}
\def\tT{{\tens{T}}}
\def\tU{{\tens{U}}}
\def\tV{{\tens{V}}}
\def\tW{{\tens{W}}}
\def\tX{{\tens{X}}}
\def\tY{{\tens{Y}}}
\def\tZ{{\tens{Z}}}


% Graph
\def\gA{{\mathcal{A}}}
\def\gB{{\mathcal{B}}}
\def\gC{{\mathcal{C}}}
\def\gD{{\mathcal{D}}}
\def\gE{{\mathcal{E}}}
\def\gF{{\mathcal{F}}}
\def\gG{{\mathcal{G}}}
\def\gH{{\mathcal{H}}}
\def\gI{{\mathcal{I}}}
\def\gJ{{\mathcal{J}}}
\def\gK{{\mathcal{K}}}
\def\gL{{\mathcal{L}}}
\def\gM{{\mathcal{M}}}
\def\gN{{\mathcal{N}}}
\def\gO{{\mathcal{O}}}
\def\gP{{\mathcal{P}}}
\def\gQ{{\mathcal{Q}}}
\def\gR{{\mathcal{R}}}
\def\gS{{\mathcal{S}}}
\def\gT{{\mathcal{T}}}
\def\gU{{\mathcal{U}}}
\def\gV{{\mathcal{V}}}
\def\gW{{\mathcal{W}}}
\def\gX{{\mathcal{X}}}
\def\gY{{\mathcal{Y}}}
\def\gZ{{\mathcal{Z}}}

% Sets
\def\sA{{\mathbb{A}}}
\def\sB{{\mathbb{B}}}
\def\sC{{\mathbb{C}}}
\def\sD{{\mathbb{D}}}
% Don't use a set called E, because this would be the same as our symbol
% for expectation.
\def\sF{{\mathbb{F}}}
\def\sG{{\mathbb{G}}}
\def\sH{{\mathbb{H}}}
\def\sI{{\mathbb{I}}}
\def\sJ{{\mathbb{J}}}
\def\sK{{\mathbb{K}}}
\def\sL{{\mathbb{L}}}
\def\sM{{\mathbb{M}}}
\def\sN{{\mathbb{N}}}
\def\sO{{\mathbb{O}}}
\def\sP{{\mathbb{P}}}
\def\sQ{{\mathbb{Q}}}
\def\sR{{\mathbb{R}}}
\def\sS{{\mathbb{S}}}
\def\sT{{\mathbb{T}}}
\def\sU{{\mathbb{U}}}
\def\sV{{\mathbb{V}}}
\def\sW{{\mathbb{W}}}
\def\sX{{\mathbb{X}}}
\def\sY{{\mathbb{Y}}}
\def\sZ{{\mathbb{Z}}}

% Entries of a matrix
\def\emLambda{{\Lambda}}
\def\emA{{A}}
\def\emB{{B}}
\def\emC{{C}}
\def\emD{{D}}
\def\emE{{E}}
\def\emF{{F}}
\def\emG{{G}}
\def\emH{{H}}
\def\emI{{I}}
\def\emJ{{J}}
\def\emK{{K}}
\def\emL{{L}}
\def\emM{{M}}
\def\emN{{N}}
\def\emO{{O}}
\def\emP{{P}}
\def\emQ{{Q}}
\def\emR{{R}}
\def\emS{{S}}
\def\emT{{T}}
\def\emU{{U}}
\def\emV{{V}}
\def\emW{{W}}
\def\emX{{X}}
\def\emY{{Y}}
\def\emZ{{Z}}
\def\emSigma{{\Sigma}}

% entries of a tensor
% Same font as tensor, without \bm wrapper
\newcommand{\etens}[1]{\mathsfit{#1}}
\def\etLambda{{\etens{\Lambda}}}
\def\etA{{\etens{A}}}
\def\etB{{\etens{B}}}
\def\etC{{\etens{C}}}
\def\etD{{\etens{D}}}
\def\etE{{\etens{E}}}
\def\etF{{\etens{F}}}
\def\etG{{\etens{G}}}
\def\etH{{\etens{H}}}
\def\etI{{\etens{I}}}
\def\etJ{{\etens{J}}}
\def\etK{{\etens{K}}}
\def\etL{{\etens{L}}}
\def\etM{{\etens{M}}}
\def\etN{{\etens{N}}}
\def\etO{{\etens{O}}}
\def\etP{{\etens{P}}}
\def\etQ{{\etens{Q}}}
\def\etR{{\etens{R}}}
\def\etS{{\etens{S}}}
\def\etT{{\etens{T}}}
\def\etU{{\etens{U}}}
\def\etV{{\etens{V}}}
\def\etW{{\etens{W}}}
\def\etX{{\etens{X}}}
\def\etY{{\etens{Y}}}
\def\etZ{{\etens{Z}}}

% The true underlying data generating distribution
\newcommand{\pdata}{p_{\rm{data}}}
% The empirical distribution defined by the training set
\newcommand{\ptrain}{\hat{p}_{\rm{data}}}
\newcommand{\Ptrain}{\hat{P}_{\rm{data}}}
% The model distribution
\newcommand{\pmodel}{p_{\rm{model}}}
\newcommand{\Pmodel}{P_{\rm{model}}}
\newcommand{\ptildemodel}{\tilde{p}_{\rm{model}}}
% Stochastic autoencoder distributions
\newcommand{\pencode}{p_{\rm{encoder}}}
\newcommand{\pdecode}{p_{\rm{decoder}}}
\newcommand{\precons}{p_{\rm{reconstruct}}}

\newcommand{\laplace}{\mathrm{Laplace}} % Laplace distribution

\newcommand{\E}{\mathbb{E}}
\newcommand{\Ls}{\mathcal{L}}
\newcommand{\R}{\mathbb{R}}
\newcommand{\emp}{\tilde{p}}
\newcommand{\lr}{\alpha}
\newcommand{\reg}{\lambda}
\newcommand{\rect}{\mathrm{rectifier}}
\newcommand{\softmax}{\mathrm{softmax}}
\newcommand{\sigmoid}{\sigma}
\newcommand{\softplus}{\zeta}
\newcommand{\KL}{D_{\mathrm{KL}}}
\newcommand{\Var}{\mathrm{Var}}
\newcommand{\standarderror}{\mathrm{SE}}
\newcommand{\Cov}{\mathrm{Cov}}
% Wolfram Mathworld says $L^2$ is for function spaces and $\ell^2$ is for vectors
% But then they seem to use $L^2$ for vectors throughout the site, and so does
% wikipedia.
\newcommand{\normlzero}{L^0}
\newcommand{\normlone}{L^1}
\newcommand{\normltwo}{L^2}
\newcommand{\normlp}{L^p}
\newcommand{\normmax}{L^\infty}

\newcommand{\parents}{Pa} % See usage in notation.tex. Chosen to match Daphne's book.

\DeclareMathOperator*{\argmax}{arg\,max}
\DeclareMathOperator*{\argmin}{arg\,min}

\DeclareMathOperator{\sign}{sign}
\DeclareMathOperator{\Tr}{Tr}
\let\ab\allowbreak


\usepackage{cy}
\usepackage{nicefrac}
\usepackage{caption}
\usepackage{subcaption}

\newcommand{\valpha}{{\bm{\alpha}}}
\newcommand{\vbeta}{{\bm{\beta}}}
\newcommand{\vbetafc}{\vbeta_{\rm fc}}
\newcommand{\vbetadiag}{\vbeta_{\rm diag}}
\newcommand{\vbetaconv}{\vbeta_{\rm conv}}
\newcommand{\flip}{\mathop{\rm flip}} % flip operation for vectors
\newcommand{\bTheta}{\bm{\Theta}}
\newcommand{\bThetafc}{\bTheta_{\rm fc}}
\newcommand{\bThetadiag}{\bTheta_{\rm diag}}
\newcommand{\bThetaconv}{\bTheta_{\rm conv}}
\newcommand{\ffc}{f_{\rm fc}}
\newcommand{\fdiag}{f_{\rm diag}}
\newcommand{\fconv}{f_{\rm conv}}
\newcommand{\tMfc}{\tM_{\rm fc}}
\newcommand{\tMdiag}{\tM_{\rm diag}}
\newcommand{\tMconv}{\tM_{\rm conv}}
\newcommand{\vrho}{{\bm{\rho}}}
\newcommand{\veta}{{\bm{\eta}}}
\newcommand{\vnu}{{\bm{\nu}}}

\allowdisplaybreaks

%%%%% overlap two symbols
% \makeatletter
% \newcommand{\superimpose}[2]{%
%   {\ooalign{$#1\@firstoftwo#2$\cr\hfil$#1\@secondoftwo#2$\hfil\cr}}}
% \makeatother
% \newcommand{\prododot}[2]{\mathpalette\superimpose{{\prod_{#1}^{#2}}{\odot}}}

\renewcommand{\Re}{\mathop{\rm Re}}
\renewcommand{\Im}{\mathop{\rm Im}}
\newcommand{\Arg}{\mathop{\rm arg}}

\usepackage{enumitem}

\usepackage{hyperref}
\usepackage{url}

\usepackage{color,colortbl}
\definecolor{darkblue}{rgb}{0.0,0.0,0.65}
\definecolor{darkred}{rgb}{0.65,0.0,0.0}
\hypersetup{
  colorlinks = true,
  citecolor  = darkblue,
  linkcolor  = darkred,
  filecolor  = darkblue,
  urlcolor   = darkblue,
}
% Title: Unifying perspectives on implicit bias in linear neural network training

% Abstract: Overparametrized neural networks have infinitely many solutions that achieve zero training error, but gradient-based methods succeed in finding solutions that generalize well. It is conjectured that the optimization algorithm and the network architecture induce an implicit bias towards favorable solutions, so understanding such a bias has become a popular topic. We study the implicit bias of gradient flow (i.e., gradient descent with infinitesimal step size) applied on linear neural network training. We consider separable classification and underdetermined linear regression problems where there exist many solutions that achieve zero training error, and characterize how the network architecture and initialization affects the final solution found by gradient flow. Our results apply to a general tensor formulation of neural networks that includes linear fully-connected networks and linear convolutional networks as special cases, while removing convergence assumptions required by prior arts. We also provide experiments that corroborate our theoretical analysis.

%\title{Unifying perspectives on implicit bias of\\ gradient descent on linear neural networks}
\title{A unifying view on implicit bias in training linear neural networks}
% \title{How does initialization determine the implicit bias gradient descent on neural networks: a unifying view}
% \title{Unifying view on implicit bias of gradient descent: how initialization affects kernel vs active regimes}

% Authors must not appear in the submitted version. They should be hidden
% as long as the \iclrfinalcopy macro remains commented out below.
% Non-anonymous submissions will be rejected without review.

\author{Chulhee Yun\thanks{Based on work performed during internship at Google Research} \\
MIT\\
\texttt{chulheey@mit.edu} \\
\And
Shankar Krishnan\\
Google Research \\
\texttt{skrishnan@google.com} \\
\And
Hossein Mobahi \\
Google Research \\
\texttt{hmobahi@google.com}
}

% The \author macro works with any number of authors. There are two commands
% used to separate the names and addresses of multiple authors: \And and \AND.
%
% Using \And between authors leaves it to \LaTeX{} to determine where to break
% the lines. Using \AND forces a linebreak at that point. So, if \LaTeX{}
% puts 3 of 4 authors names on the first line, and the last on the second
% line, try using \AND instead of \And before the third author name.

\newcommand{\fix}{\marginpar{FIX}}
\newcommand{\new}{\marginpar{NEW}}

\iclrfinalcopy % Uncomment for camera-ready version, but NOT for submission.
\begin{document}


\maketitle

\begin{abstract}
% We study the implicit bias of gradient flow (i.e., gradient descent with infinitesimal step size) on linear neural network training. We consider separable classification and underdetermined linear regression problems where there exist many solutions that achieve zero training error, and characterize how the network architecture and initialization affects the final solution found by gradient flow. Our results apply to a general tensor formulation of neural networks that includes fully-connected networks, diagonal networks, and convolutional networks as special cases, while removing convergence assumptions required by prior research. We also provide experiments that corroborate our theoretical analysis.

%We study the implicit bias of gradient flow (i.e., gradient descent with infinitesimal step size) on linear neural network training. We propose a tensor formulation of neural networks that includes fully-connected, diagonal, and convolutional networks as special cases, and investigate the linear version of the formulation called linear tensor networks. With this formulation, we can characterize the convergence direction of the network parameters as singular vectors of a tensor defined by the network. For $L$-layer linear tensor networks that are orthogonally decomposable, we show that gradient flow on separable classification finds a stationary point of the $\ell_{2/L}$ max-margin problem in a ``transformed'' input space defined by the network. For underdetermined regression, we prove that gradient flow finds a global minimum which minimizes a norm-like function that interpolates between weighted $\ell_1$ and $\ell_2$ norms in the transformed input space. Our theorems subsume existing results in the literature while removing standard convergence assumptions. We also provide experiments that corroborate our analysis.

We study the implicit bias of gradient flow (i.e., gradient descent with infinitesimal step size) on linear neural network training. We propose a tensor formulation of neural networks that includes fully-connected, diagonal, and convolutional networks as special cases, and investigate the linear version of the formulation called linear tensor networks. With this formulation, we can characterize the convergence direction of the network parameters as singular vectors of a tensor defined by the network. For $L$-layer linear tensor networks that are orthogonally decomposable, we show that gradient flow on separable classification finds a stationary point of the $\ell_{2/L}$ max-margin problem in a ``transformed'' input space defined by the network. For underdetermined regression, we prove that gradient flow finds a global minimum which minimizes a norm-like function that interpolates between weighted $\ell_1$ and $\ell_2$ norms in the transformed input space. Our theorems subsume existing results in the literature while removing standard convergence assumptions. We also provide experiments that corroborate our analysis.

%With $L$-layer linear tensor networks, we consider solving classification and regression problems that have many global minima, and characterize how the network architecture and initialization affects the solution found by gradient flow. 

\end{abstract}

%\vspace*{-5pt}
\section{Introduction}
%\vspace*{-5pt}
Overparameterized neural networks have infinitely many solutions that achieve zero training error, and such global minima have different generalization performance. Moreover, training a neural network is a high-dimensional nonconvex problem, which is typically intractable to solve.
However, the success of deep learning indicates that first-order methods such as gradient descent or stochastic gradient descent (GD/SGD) not only (a) succeed in finding global minima, but also (b) are biased towards solutions that generalize well, which largely has remained a mystery in the literature.

To explain part (a) of the phenomenon, there is a growing literature studying the convergence of GD/SGD on overparameterized neural networks (e.g., \citet{du2018gradientB, du2018gradientA, allen2018convergence, zou2018stochastic,jacot2018neural,oymak2020towards}, and many more). 
There are also convergence results that focus on linear networks, without nonlinear activations \citep{bartlett2018gradient,arora2019convergence,wu2019global,du2019width,hu2020provable}.
These results typically focus on the convergence of loss, hence do not address \emph{which} of the many global minima is reached.

Another line of results tackles part (b), by studying the implicit bias or regularization of gradient-based methods on neural networks or related problems \citep{gunasekar2017implicit,gunasekar2018characterizing,gunasekar2018implicit,arora2018optimization,soudry2018implicit,ji2019gradient,arora2019implicit,woodworth2020kernel,chizat2020implicit,gissin2020implicit}. These results have shown interesting progress that even without explicit regularization terms in the training objective, algorithms such as GD applied on neural networks have an \emph{implicit bias} towards certain solutions among the many global minima. 
However, the results along this line are still in the preliminary steps, most of them pertaining only to multilinear models such as linear models, linear neural networks and matrix/tensor decomposition.

Our paper is motivated from two limitations that are common in the implicit bias literature. First, most analyses of implicit bias are done in a \emph{case-by-case} manner. A given theorem on a specific network does not provide useful insights on other architectures, which calls for a unifying framework that can incorporate different architectures in a single formulation.
Next, in proving implicit bias results, many existing theorems rely on \emph{convergence assumptions} such as global convergence of loss to zero and/or directional convergence of parameters and gradients. Ideally, such convergence assumptions should be removed because they cannot be tested \emph{a priori} and there are known examples where optimization algorithms do not converge to global minima under certain initializations \citep{bartlett2018gradient, arora2019convergence}.

%\vspace*{-4pt}
\subsection{Summary of our contributions}
%\vspace*{-4pt}
We study the implicit bias of gradient flow (GD with infinitesimal step size) on linear neural networks.
Following recent progress on this topic, we consider classification and regression problems that have multiple solutions attaining zero training error.
In light of the limitations discussed above, we provide theorems on a \textbf{general tensor framework} of networks that yield corollaries on specific architecture to recover known results.
We also make significant efforts to \textbf{remove convergence assumptions}; our theorems rely on less assumptions (if any) compared to the existing results in the literature.
Some key contributions are summarized below.
%Our analyses apply to a \textbf{general class of networks}, and prove \textbf{both convergence and implicit bias}, providing a more \emph{complete characterization} of the algorithm trajectory without relying on convergence assumptions.
%However, we emphasize that our implicit bias statements still holds true as long as the loss converges to zero.
% We formulate neural networks into a general tensor form, and identify how the architecture and initialization affects limit directions/points of the algorithm. 
% Our results extend existing results in the literature, while \emph{removing} the convergence assumptions required in prior research.
\begin{list}{\tiny{$\bullet$}}{\leftmargin=1.8em}
\setlength\itemsep{0em}
\item We propose a general \emph{tensor formulation} of nonlinear neural networks which includes many network architectures considered in the literature. For the purpose of implicit bias analysis, we focus on the linear version of this formulation (i.e., no nonlinear activations), called \emph{linear tensor networks}.
\item For linearly separable classification, we prove that linear tensor network parameters converge in direction to \emph{singular vectors} of a tensor defined by the network. 
As a \emph{corollary}, we show that linear fully-connected networks converge to the $\ell_2$ max-margin solution~\citep{ji2020directional}.
\item For separable classification, we further show that if the linear tensor network is orthogonally decomposable (Assumption~\ref{assm:complex-diag-decomp}), the gradient flow finds the $\ell_{2/{\rm depth}}$ max-margin solution in the singular value space, leading the parameters to converge to the \emph{top singular vectors} of the tensor when ${\rm depth} = 2$. This theorem subsumes known results on linear convolutional networks and diagonal networks proved in \citet{gunasekar2018implicit}, under fewer convergence assumptions.
%\item For underdetermined linear regression, we study the limit points of gradient flow on orthogonally decomposable networks (Assumption~\ref{assm:complex-diag-decomp}), and provide a full characterization of the limit points. This theorem covers results on deep matrix sensing \citep{arora2019implicit} as a special case, and extends a similar recent result \citep{woodworth2020kernel} to a broader class of networks.
\item For underdetermined linear regression, we characterize the limit points of gradient flow on orthogonally decomposable networks (Assumption~\ref{assm:complex-diag-decomp}). Proven without convergence assumptions, this theorem covers results on deep matrix factorization~\citep{arora2019implicit} as a special case, and extends a recent result \citep{woodworth2020kernel} to a broader class of networks.
\item For underdetermined linear regression with deep linear fully-connected networks, we prove that the network converges to the minimum $\ell_2$ norm solutions as we scale the initialization to zero.
\item Lastly, we present simple experiments that corroborate our theoretical analysis. Figure~\ref{fig:reg} shows that our predictions of limit points match with solutions found by GD. 
\end{list}

% \begin{figure}[h]
% \begin{center}
% %\framebox[4.0in]{$\;$}
% \fbox{\rule[-.5cm]{0cm}{4cm} \rule[-.5cm]{4cm}{0cm}}
% \end{center}
% \caption{Sample figure caption.}
% \end{figure}

% \begin{figure}[t]
%      \centering
%      \begin{subfigure}[t]{0.35\textwidth}
%          \centering
%          \includegraphics[width=\textwidth]{plot_reg_cropped.png}
%          \caption{Regression.}
%          \label{fig:reg}
%      \end{subfigure}
%      \quad\quad
%      \begin{subfigure}[t]{0.35\textwidth}
%          \centering
%          \includegraphics[width=\textwidth]{plot_clsf_cropped.png}
%          \caption{Classification.}
%          \label{fig:clsf}
%      \end{subfigure}
%     %  \hfill
%     %  \begin{subfigure}[t]{0.32\textwidth}
%     %      \centering
%     %      \includegraphics[width=\textwidth]{plot_1.png}
%     %      \caption{}
%     %      %\label{fig:three sin x}
%     %  \end{subfigure}
%     \caption{Gradient descent trajectories of linear fully-connected (purple), diagonal (green), and convolutional (blue) networks initialized at the same initial coefficients (circles on dotted lines). They highlight different implicit biases of networks induced by their architecture. Also,
%     (a) shows that our theoretical predictions (circles on red line, the set of global minima) agree with the solution found by GD. For details of the experimental setup, please refer to Section~\ref{sec:exp}. }
%     %     \label{fig:three graphs}
%     \vspace*{-12pt}
% \end{figure}

\begin{figure}[t]
     \centering
    %  \begin{subfigure}[t]{0.32\textwidth}
    %      \centering
    %      \includegraphics[width=\textwidth]{reg1_crop.png}
    %  \end{subfigure}
    %  \begin{subfigure}[t]{0.32\textwidth}
    %      \centering
    %      \includegraphics[width=\textwidth]{reg2_crop.png}
    %  \end{subfigure}
    %  \begin{subfigure}[t]{0.32\textwidth}
    %      \centering
    %      \includegraphics[width=\textwidth]{reg3_crop.png}
    %  \end{subfigure}
    
    \begin{subfigure}[t]{0.48\textwidth}
         \centering
         \includegraphics[width=\textwidth]{regression_bias_a_1en2.png}
     \end{subfigure}
     ~
    %  \begin{subfigure}[t]{0.32\textwidth}
    %      \centering
    %      \includegraphics[width=\textwidth]{regression_bias_a_5en1.png}
    %  \end{subfigure}
     \begin{subfigure}[t]{0.48\textwidth}
         \centering
         \includegraphics[width=\textwidth]{regression_bias_a_1e0.png}
     \end{subfigure}
    
    %  \begin{subfigure}[t]{0.32\textwidth}
    %      \centering
    %      \includegraphics[width=\textwidth]{clsf1_crop.png}
    %  \end{subfigure}
    %  \begin{subfigure}[t]{0.32\textwidth}
    %      \centering
    %      \includegraphics[width=\textwidth]{clsf2_crop.png}
    %  \end{subfigure}
    %  \begin{subfigure}[t]{0.32\textwidth}
    %      \centering
    %      \includegraphics[width=\textwidth]{clsf3_crop.png}
    %  \end{subfigure}
    
    \caption{Gradient descent trajectories of linear coefficients of linear fully-connected, diagonal, and convolutional networks on a regression task, initialized with different initial scales $\alpha = 0.01, 1$.
    Networks are initialized at the same coefficients (circles on purple lines), but follow different trajectories due to implicit biases of networks induced from their architecture. The figures show that our theoretical predictions on limit points (circles on yellow line, the set of global minima) agree with the solution found by GD. For details of the experimental setup, see Section~\ref{sec:exp}. }
    \label{fig:reg}
    %\vspace*{-14pt}
\end{figure}
%\vspace*{-6pt}
\section{Problem settings and related works}
%\vspace*{-5pt}
%\subsection{Notation}
We first define notation used in the paper.
Given a positive integer $a$, let $[a]\defeq \{1, \dots, a\}$.
We use $\mI_d$ to denote the $d \times d$ identity matrix.
Given a matrix $\mA$, we use $\vect(\mA)$ to denote its vectorization, i.e., the concatenation of all columns of $\mA$.
For two vectors $\va$ and $\vb$, let $\va \otimes \vb$ be their tensor product, $\va \odot \vb$ be their element-wise product, and $\va^{\odot k}$ be the element-wise $k$-th power of $\va$.
% We use $\va \cdot \vb$ to denote the dot product of $\va$ and $\vb$. 
% As we define later, the symbol $\cdot$ is also used to denote the multilinear product between a tensor and linear maps.
Given an order-$L$ tensor $\tA \in \reals^{k_1 \times \cdots \times k_L}$, we use $[\tA]_{j_1, \dots, j_L}$ to denote the $(j_1, j_2, \dots, j_L)$-th element of $\tA$, where $j_l \in [k_l]$ for all $l \in [L]$.
%We use the same indexing notation for elements of vectors $\va$ (order-1 tensors) and matrices $\mA$ (order-2 tensors).
In element indexing, we use $\cdot$ to denote all indices in the corresponding dimension, and $a:b$ to denote all indices from $a$ to $b$. For example, for a matrix $\mA$, $[\mA]_{\cdot,4:6}$ denotes a submatrix that consists of 4th--6th columns of $\mA$. 
The square bracket notation for indexing overloads with $[a]$ when $a \in \naturals$, but they will be distinguishable from the context.
Since element indices start from $1$, we re-define the modulo operation $a \bmod d \defeq a - \lfloor \frac{a-1}{d} \rfloor d \in [d]$ for $a > 0$.
We use $\ve^k_j$ to denote the $j$-th stardard basis vector of the vector space $\reals^{k}$.
Lastly, we define the multilinear multiplication between a tensor and linear maps, which can be viewed as a generalization of left- and right-multiplication on a matrix.
Given a tensor $\tA \in \reals^{k_1 \times \cdots \times k_L}$ and linear maps $\mB_l \in \reals^{p_l \times k_l}$ for $l \in [L]$, we define the multilinear multiplication $\circ$ between them as
\begin{align*}
    \tA \circ (\mB_1^T, \mB_2^T, \dots, \mB_L^T)
    &= \sum\nolimits_{j_1, \dots, j_L} [\tA]_{j_1, \dots, j_L} (\ve^{k_1}_{j_1} \otimes \cdots \otimes \ve^{k_L}_{j_L}) \circ (\mB_1^T, \dots, \mB_L^T)\\
    &\defeq \sum\nolimits_{j_1, \dots, j_L} [\tA]_{j_1, \dots, j_L} (\mB_1 \ve^{k_1}_{j_1} \otimes \cdots \otimes \mB_L \ve^{k_L}_{j_L}) 
    \in \reals^{p_1\times \cdots \times p_L}\,.
\end{align*}

%\vspace*{-18pt}
\subsection{Problem settings}
%\vspace*{-4pt}
\label{sec:prob-setting}
We are given a dataset $\{(\vx_i, y_i)\}_{i=1}^n$, where $\vx_i \in \reals^d$ and $y_i \in \reals$. We let $\mX \in \reals^{n \times d}$ and $\vy \in \reals^n$ be the data matrix and the label vector, respectively.
We study binary classification and linear regression in this paper, focusing on the settings where there exist \emph{many} global solutions.
For binary classification, we assume $y_i \in \{\pm 1\}$ and that the data is separable: there exists a unit vector $\vz$ and a constant $\gamma > 0$ such that $y_i \vx_i^T \vz \geq \gamma$ for all $i \in [n]$.
For regression, we consider the underdetermined case ($n \leq d$) where there are many parameters $\vz \in \reals^d$ such that $\mX \vz = \vy$. Throughout the paper, we assume that $\mX$ has full row rank.

We use $f(\cdot;\bTheta): \reals^d \to \reals$ to denote a neural network parametrized by $\bTheta$.
Given the network and the dataset, we consider minimizing the training loss $\mc L(\bTheta) \defeq \sum\nolimits_{i=1}^n \ell(f(\vx_i; \bTheta), y_i)$ over $\bTheta$.
% \begin{equation}
% \label{eq:erm}
%     \minimize_{\bTheta} \quad \mc L(\bTheta) \defeq \sum\nolimits_{i=1}^n
%     \ell(f(\vx_i; \bTheta), y_i),
% \end{equation}
Following previous results (e.g., \citet{lyu2020gradient, ji2020directional}), we use the exponential loss $\ell(\hat y, y) = \exp(-\hat y y)$ for classification problems. For regression, we use the squared error loss $\ell(\hat y, y) = \half (\hat y-y)^2$.
On the algorithm side, we minimize $\mc L$ using gradient flow, which can be viewed as GD with infinitesimal step size. The gradient flow dynamics is defined as $\frac{d}{dt} \bTheta = - \nabla_{\bTheta} \mc L(\bTheta)$.

%\vspace*{-4pt}
\subsection{Related works}
%\vspace*{-4pt}
\label{sec:related}
\textbf{Gradient flow/descent in separable classification.~~}
For linear models $f(\vx;\vz) = \vx^T \vz$ with separable data, \citet{soudry2018implicit} show that the GD run on $\mc L$ drives $\norm{\vz}$ to $\infty$, but $\vz$ converges in direction to the $\ell_2$ max-margin classifier. The limit direction of $\vz$ is aligned with the solution of
\begin{equation}
\label{eq:svm}
    \minimize\nolimits_{\vz \in \reals^d}\quad \norm{\vz} \quad \subjectto \quad y_i \vx_i^T\vz \geq 1 \text{ for } i \in [n],
\end{equation}
where the norm in the cost is the $\ell_2$ norm.
\citet{nacson2019convergence,nacson2019stochastic,gunasekar2018characterizing,ji2019implicit,ji2019refined} extend these results to other (stochastic) algorithms and non-separable settings.

\citet{gunasekar2018implicit} study the same problem on linear neural networks and show that GD exhibits different implicit biases depending on the architecture. 
The authors show that the linear coefficients of the network converges in direction to the solution of \eqref{eq:svm} with different norms: $\ell_2$ norm for linear fully-connected networks, $\ell_{2/L}$ (quasi-)norm for diagonal networks, and DFT-domain $\ell_{2/L}$ (quasi-)norm for convolutional networks with full-length filters. Here, $L$ denotes the depth. We note that \citet{gunasekar2018implicit} assume that GD globally minimizes the loss, and the network parameters and the gradient with respect to the linear coefficients converge in direction.
%; in contrast, our results in Section~\ref{sec:clsf} provide an extension to more general networks, while removing convergence assumptions. \todo{add more results by Ji, Nacson, Soudry, etc}
Subsequent results \citep{ji2019gradient,ji2020directional} remove such assumptions for linear fully-connected networks. 

A recent line of results \citep{nacson2019lexicographic,lyu2020gradient,ji2020directional} studies general homogeneous models and show divergence of parameters to infinity, monotone increase of smoothed margin, directional convergence and alignment of parameters (see Section~\ref{sec:clsf} for details). \citet{lyu2020gradient} also characterize the limit direction of parameters as the KKT point of a nonconvex max-margin problem similar to \eqref{eq:svm}, but this characterization does not provide useful insights for the functions $f(\cdot;\bTheta)$ represented by specific architectures, because the formulation is in the parameter space $\bTheta$. Also, these results require that gradient flow/descent has already reached 100\% training accuracy. Although we study a more restrictive set of networks (i.e., deep linear), we provide a more complete characterization of the implicit bias for the functions $f(\cdot;\bTheta)$, without assuming 100\% training accuracy.

\textbf{Gradient flow/descent in linear regression.~~}
It is known that for linear models $f(\vx;\vz) = \vx^T \vz$, GD converges to the global minimum that is closest in $\ell_2$ distance to the initialization (see e.g., \citet{gunasekar2018characterizing}). However, relatively less is known for deep networks, even for linear networks. This is partly because the parameters do not diverge to infinity, hence making limit points highly dependent on the initialization; this dependency renders analysis difficult.
A related problem of matrix sensing aims to minimize $\sum_{i=1}^n (y_i - \< \mA_i, \mW_1 \cdots \mW_L \> )^2$ over $\mW_1, \dots, \mW_L \in \reals^{d\times d}$. It is shown in \citet{gunasekar2017implicit, arora2019implicit} that if the sensor matrices $\mA_i$ commute and we initialize all $\mW_l$'s to $\alpha \mI$, GD finds the minimum nuclear norm solution as $\alpha \to 0$.

\citet{chizat2019lazy} show that if a network is zero at initialization, and we scale the network output by a factor of $\alpha \to \infty$, then the GD dynamics enters a ``lazy regime'' where the network behaves like a first-order approximation at its initialization, as also seen in results studying kernel approximations of neural networks and convergence of GD in the corresponding RKHS (e.g., \citet{jacot2018neural}). 

%To bridge such a gap between the two extremes, i.e., ``active/rich'' regime in $\alpha \to 0$ vs.\ ''lazy/kernel'' regime in $\alpha \to \infty$,
\citet{woodworth2020kernel} study linear regression with a diagonal network of the form $f(\vx;\vw_+,\vw_-) = \vx^T(\vw_+^{\odot L}-\vw_-^{\odot L})$, where $\vw_+$ and $\vw_-$ are identically initialized $\vw_+(0) = \vw_-(0) = \alpha \bar \vw$. The authors show that the global minimum reached by GD minimizes a norm-like function which interpolates between (weighted) $\ell_1$ norm ($\alpha \to 0$) and $\ell_2$ norm ($\alpha \to \infty$).
In our paper, we consider a more general class of orthogonally decomposable networks, and obtain similar results interpolating between weighted $\ell_1$ and $\ell_2$ norms. We also remark that our results include the results in \citet{arora2019implicit} as a special case, and we do not assume convergence to global minima, as done in \citet{gunasekar2017implicit,arora2019implicit,woodworth2020kernel}.

%\citet{gissin2020implicit} study a similar model and characterized the phenomenon of the ``incremental learning'' but their analysis is limited to the case where the global minimum is unique.

% In light of \citep{woodworth2020kernel}, these binary classification results may sound confusing because it seems like they do not depend on the scale of initialization.
% In fact, we empirically observe that a similar kernel vs active phenomenon occurs in the binary classification setup too.

% \subsection{Convergence of gradient descent/flow on linear neural networks}
% There are...
% \citet{bartlett2018gradient,arora2018optimization,arora2019convergence,wu2019global,hu2020provable,zou2020global,lyu2020gradient,ji2020directional}
% and many more... (to be added)

% In the exponential/logistic loss case, \citet{ji2020directional} show that if a model $f(\vx;\bTheta)$ is locally Lipschitz, $L$-homogeneous as a function of $\bTheta$, and definable in some o-minimal structure (Assumption~2.1 of \citep{ji2020directional}) and the gradient flow starts at a point with perfect traning accuracy, then the gradient flow $\dot \bTheta = - \nabla_{\bTheta} \mc L(\bTheta)$ dynamics of the parameters $\bTheta$ converge in direction. Under additional condition of locally Lipschitz gradients $\nabla_{\bTheta} f(\vx;\bTheta)$, the parameters $\bTheta$ align with their gradient direction $-\nabla_{\bTheta} \mc L(\bTheta)$.
%\vspace*{-5pt}
\section{Tensor formulation of neural networks}
%\vspace*{-5pt}
\label{sec:tensorform}
In this section, we present a general tensor formulation of neural networks.
Given an input $\vx \in \reals^d$, the network uses a linear map $\tM$ that maps $\vx$ to an order-$L$ tensor $\tM(\vx) \in \reals^{k_1 \times \cdots \times k_L}$,
where $L \geq 2$. Using parameters $\vv_l \in \reals^{k_l}$ and activation $\phi$, the network computes its layers as the following:
\begin{equation}
\label{eq:tensorrep}
\begin{aligned}
    \tH_1(\vx) &= \phi\left (\tM(\vx) \circ (\vv_1, \mI_{k_2}, \dots, \mI_{k_L}) \right ) \in \reals^{k_2 \times \cdots \times k_L},\\
    \tH_l(\vx) &= \phi\left (\tH_{l-1}(\vx) \circ (\vv_l, \mI_{k_{l+1}}, \dots, \mI_{k_L}) \right ) \in \reals^{k_{l+1} \times \dots, k_L},~~\text{ for } l = 2, \dots, L-1,\\
    f(\vx;\bTheta) &= \tH_{L-1}(\vx) \circ \vv_L \in \reals.
\end{aligned}
\end{equation}
\begin{figure}[t]
	\centering
	\includegraphics[width=0.8\textwidth]{tensor_illust.png}
	\caption{Illustration of tensor formulation, for $L = 3, k_1 = 5, k_2 = 4, k_3 = 3$.}
	\label{fig:tensor}
\end{figure}
% we focus on neural networks where each layer can be represented as a multilinear multiplication followed by activation. More concretely, we consider a $L$-layer neural network with parameters $\vv_1 \in \reals^{k_1}, \vv_2 \in \reals^{k_2}, \dots, \vv_L \in \reals^{k_L}$ and activation $\phi$ applied element-wise. Given an input vector $\vx \in \reals^d$, the network uses a linear map $\tM$ that maps $\vx$ to a order-$L$ tensor $\tM(\vx) \in \reals^{k_1 \times \cdots \times k_L}$, and then computes each layer in the following way:
% % \begin{equation}
% % \label{eq:tensorrep}
% % \begin{aligned}
% %     \tH_1(\vx) &= \phi\left ((\vv_1^T, \mI_{k_2}, \dots, \mI_{k_L}) \cdot \tM(\vx) \right ) \in \reals^{k_2 \times \cdots \times k_L},\\
% %     \tH_2(\vx) &= \phi\left ((\vv_2^T, \mI_{k_3}, \dots, \mI_{k_L}) \cdot \tH_1(\vx) \right ) \in \reals^{k_3 \times \cdots \times k_L},\\
% %     &\dots,\\
% %     \tH_{L-1}(\vx) &= \phi\left ((\vv_{L-1}^T, \mI_{k_L}) \cdot \tH_{L-2}(\vx) \right ) \in \reals^{k_L},\\
% %     f(\vx;\bTheta) &= \vv_L^T \tH_{L-1}(\vx) \in \reals.
% % \end{aligned}
% % \end{equation}
% \begin{equation}
% \label{eq:tensorrep}
% \begin{aligned}
%     \tH_1(\vx) &= \phi\left (\tM(\vx) \cdot (\vv_1, \mI_{k_2}, \dots, \mI_{k_L}) \right ) \in \reals^{k_2 \times \cdots \times k_L},\\
%     \tH_2(\vx) &= \phi\left (\tH_1(\vx) \cdot (\vv_2, \mI_{k_3}, \dots, \mI_{k_L}) \right ) \in \reals^{k_3 \times \cdots \times k_L},\\
%     &\dots,\\
%     \tH_{L-1}(\vx) &= \phi\left (\tH_{L-2}(\vx) \cdot (\vv_{L-1}, \mI_{k_L}) \right ) \in \reals^{k_L},\\
%     f(\vx;\bTheta) &= \tH_{L-1}(\vx) \cdot \vv_L \in \reals.
% \end{aligned}
% \end{equation}

We use $\bTheta$ to denote the collection of all parameters $(\vv_1, \dots, \vv_L)$. We call $\tM(\vx)$ the \emph{data tensor}.
Figure~\ref{fig:tensor} illustrates how our new tensor formulation calculates its scalar output in a feedforward manner. Each row in the figure represents a layer. At the $l$-th hidden layer, the parameter vector $\vv_l$ takes inner products with fibers in the order-$(L+1-l)$ tensor $\tH_{l-1}(\vx)$ along the corresponding dimension. The result is an order-$(L-l)$ tensor, which goes through the entry-wise activation function $\phi$ and becomes the output $\tH_l(\vx)$ of the hidden layer.

Although this new formulation may look a bit odd in the first glance, it is general enough to capture many network architectures considered in the literature, including fully-connected networks, diagonal networks, and circular convolutional networks. 
We formally define these architectures below.
%We define these architectures below; due to space limits, we defer the description of how they can be written in the tensor form~\eqref{eq:tensorrep} to Appendix~\ref{sec:tensorform-ex}.

\textbf{Diagonal networks.}~~
An $L$-layer diagonal network is written as
\begin{align}
\label{eq:diagnet}
    \fdiag (\vx; \bThetadiag) = \phi(\cdots\phi (\phi(\vx \odot \vw_1) \odot \vw_2)\cdots \odot \vw_{L-1})^T \vw_L,
\end{align}
where $\vw_l \in \reals^d$ for $l \in [L]$. The representation of $\fdiag$ as the tensor form~\eqref{eq:tensorrep} is straightforward. Let $\tMdiag(\vx) \in \reals^{d \times \cdots \times d}$ have 
%\begin{align*}
   $[\tMdiag(\vx)]_{j,j,\dots,j} = [\vx]_j$,
%\end{align*}
while all the remaining entries of $\tMdiag(\vx)$ are set to zero.
We can set $\vv_l = \vw_l$ for all $l$, and $\tM = \tMdiag$ to verify that \eqref{eq:tensorrep} and \eqref{eq:diagnet} are equivalent.

% This linear diagonal network is similar to the simple homogeneous models considered in \citep{gissin2020implicit, woodworth2020kernel, moroshko2020implicit}.


\textbf{Circular convolutional networks.}~~
The tensor formulation~\eqref{eq:tensorrep} includes convolutional networks
\begin{align}
\label{eq:convnet}
    \fconv(\vx; \bThetaconv) = \phi(\cdots\phi(\phi(\vx \star \vw_1) \star \vw_2 ) \cdots  \star \vw_{L-1})^T \vw_L,
\end{align}
where $\vw_l \in \reals^{k_l}$ with $k_l \leq d$ and $k_L = d$, and $\star$ defines the circular convolution: for any $\va \in \reals^d$ and $\vb \in \reals^k$ ($k \leq d$), we have $\va \star \vb \in \reals^d$ defined as
$[\va \star \vb]_i = \sum_{j=1}^k [\va]_{(i+j-1) \bmod d} [\vb]_j$, for $i \in [d]$.
%Note that in \eqref{eq:convnet}, we have $L-1$ circular convolutional layers of filter size $k_1, \dots, k_{L-1}$, followed by a final fully-connected layer. 
Define $\tMconv(\vx) \in \reals^{k_1 \times \cdots \times k_L}$ as
%\begin{equation*}
$[\tMconv(\vx)]_{j_1,j_2,\dots,j_L} = 
[\vx]_{(\sum_{l=1}^L j_l-L+1) \bmod d}$
%\end{equation*}
for $j_l \in [k_l]$, $l \in [L]$. 
% For example, if $L = 2$, $\tMconv(\vx)$ is a $k_1 \times d$ ``Hankel-like'' matrix:
% \begin{equation*}
%     \tMconv(\vx)
%     =
%     \begin{bmatrix}
%     [\vx]_1 & [\vx]_2 & \hdots & [\vx]_{d}\\
%     [\vx]_2 & [\vx]_3 & \hdots & [\vx]_{1}\\
%     \vdots & \vdots & \ddots & \vdots\\
%     [\vx]_{k_1} & [\vx]_{k_1+1} & \hdots & [\vx]_{k_1-1}
%     \end{bmatrix}.
% \end{equation*}
Setting $\vv_l = \vw_l$ and $\tM = \tMconv$, one can verify that \eqref{eq:tensorrep} and \eqref{eq:convnet} are identical.

\textbf{Fully-connected networks.}~~
An $L$-layer fully-connected network is defined as
\begin{equation}
\label{eq:fcnet}
    \ffc (\vx; \bThetafc) = \phi(\cdots\phi (\phi(\vx^T \mW_1) \mW_2)\cdots \mW_{L-1}) \vw_L,
\end{equation}
where $\mW_l \in \reals^{d_{l} \times d_{l+1}}$ for $l \in [L-1]$ (we use $d_1 = d$) and $\vw_L \in \reals^{d_{L}}$.
One can represent $\ffc$ as the tensor form~\eqref{eq:tensorrep} by defining parameters $\vv_l = \vect(\mW_l)$ for $l \in [L-1]$ and $\vv_L = \vw_L$, and constructing the tensor $\tMfc(\vx)$ by a recursive ``block diagonal'' manner.
For example, if $L = 2$, we can define $\tMfc(\vx) \in \reals^{d_1 d_2 \times d_2}$ to be the Kronecker product of $\mI_{d_2}$ and $\vx$. For deeper networks, we defer the full description of $\tMfc(\vx)$ to Appendix~\ref{sec:tensorform-ex}.

\textbf{Our focus: linear tensor networks.}~~
Throughout this section, we have used the activation $\phi$ to motivate our tensor formulation~\eqref{eq:tensorrep} for neural networks with nonlinear activations. For the remaining of the paper, we study the case whose activation is \emph{linear}, i.e., $\phi(t) = t$. In this case, 
\begin{equation}
\label{eq:tensorrep-linear}
    f(\vx;\bTheta) = \tM(\vx) \circ (\vv_1, \vv_2, \dots, \vv_L).
\end{equation}
We will refer to~\eqref{eq:tensorrep-linear} as \emph{linear tensor networks}, where ``linear'' is to indicate that the activation is linear.
Note that as a function of parameters $\vv_1, \dots, \vv_L$, $f(\vx;\bTheta)$ is in fact multilinear.
We also remark that when depth $L=2$, the data tensor $\tM(\vx)$ is a $k_1 \times  k_2$ matrix and the network formulation boils down to $f(\vx;\bTheta) = \vv_1 ^T \tM(\vx) \vv_2$.

Since the data tensor $\tM(\vx)$ is a linear function of $\vx$, the linear tensor network is also a linear function of $\vx$. Thus, the output of the network can also be written as $f(\vx;\bTheta) = \vx^T \vbeta(\bTheta)$, where $\vbeta(\bTheta) \in \reals^d$ denotes the \emph{linear coefficients} computed as a function of the network parameters $\bTheta$. 
%In the following sections, we will present results for general linear tensor networks~\eqref{eq:tensorrep-linear} which also apply to special cases such as linear fully-connected, linear diagonal, and linear convolutional networks.
% Whenever we refer to a specific architecture, we use corresponding subscripts for linear coefficients e.g., $\vbetafc(\bThetafc) = \mW_1 \cdots \mW_{L-1} \vw_L$.
% The vectors $\vbetadiag(\bThetadiag)$ and $\vbetaconv(\bThetaconv)$ can also be computed similarly, as functions of the corresponding network parameters.
Since the linear tensor network $f(\vx;\bTheta)$ is linear in $\vx$, the expressive power of $f$ is at best a linear model $\vx \mapsto \vx^T \vz$. 
%For the specific architectures (fully-connected, diagonal, and convolutional) considered above, one can easily check that regardless of the depth and the architecture, we can choose suitable parameters $\bTheta$ such that $\vbeta(\bTheta) = \vz$ for any $\vz \in \reals^d$. Hence, the expressive power of these architectures are all identical to simple linear models.
However, even though the models have the same expressive power, their architectural differences lead to different implicit biases in training, which is the focus of our investigation in this paper.
Studying separable classification and underdetermined regression is useful for highlighting such biases because there are \emph{infinitely many} coefficients that perfectly classify or fit the dataset.

For our linear tensor network, the evolution of parameters $\vv_l$ under gradient flow dynamics reads
\begin{align*}
    \dot \vv_l 
    &= -\nabla_{\vv_l} \mc L(\bTheta) 
    = -\sum\nolimits_{i=1}^n \ell'(f(\vx_i;\bTheta),y_i) \tM(\vx_i) \circ (\vv_1, \dots, \vv_{l-1}, \mI_{k_l}, \vv_{l+1}, \dots, \vv_L)\\
    &= \tM(-\mX^T \vr) \circ (\vv_1, \dots, \vv_{l-1}, \mI_{k_l}, \vv_{l+1}, \dots, \vv_L),~~\forall l \in [L],
%    \vv_l(0) &= \alpha \bar \vv_l,,
\end{align*}
where we initialize $\vv_l(0) = \alpha \bar \vv_l$, for $l \in [L]$. We refer to $\alpha$ and $\bar \vv_l$ as the \emph{initial scale} and \emph{initial direction}, respectively. We note that we do not restrict $\bar \vv_l$'s to be unit vectors, in order to allow different scaling (at initialization) over different layers. The vector $\vr \in \reals^{n}$ is the \emph{residual vector}, and each component of $\vr$ is defined as
\begin{equation}
\label{eq:residual-vec}
    [\vr]_i = \ell'(f(\vx_i;\bTheta),y_i) =
    \begin{cases}
    -y_i \exp(-y_i f(\vx_i;\bTheta)) & \text{ for classification, }\\
    f(\vx_i;\bTheta) - y_i & \text{ for regression. }
    \end{cases}
\end{equation}
%\vspace*{-10pt}
\section{Implicit bias of gradient flow in separable classification}
%\vspace*{-5pt}
\label{sec:clsf}
In this section, we present our results on the implicit bias of gradient flow in binary classification with linearly separable data. Recent papers~\citep{lyu2020gradient, ji2020directional} on this separable classification setup prove that after 100\% training accuracy has been achieved by gradient flow (along with other technical conditions), the parameters of $L$-homogeneous models diverge to infinity, while converging in direction that aligns with the direction of the negative gradient.
Mathematically,
% \begin{align*}
%     \lim_{t\to\infty} \norm{\vv_l(t)} = \infty,
%     ~~
%     \lim_{t\to\infty}\frac{\vv_l(t)}{\norm{\vv_l(t)}} = \vv_l^\infty,
%     ~~
%     \lim_{t\to\infty} \frac{\vv_l(t)^T \nabla_{\vv_l} \mc L(\bTheta(t))}{\norm{\vv_l(t)} \norm{\nabla_{\vv_l} \mc L(\bTheta(t))}}
%     = -1.
% \end{align*}
\begin{align*}
    \lim_{t\to\infty} \norm{\bTheta(t)} = \infty,
    ~~
    \lim_{t\to\infty}\tfrac{\bTheta(t)}{\norm{\bTheta(t)}} = \bTheta^\infty,
    ~~
    \lim_{t\to\infty} \tfrac{\bTheta(t)^T \nabla_{\bTheta} \mc L(\bTheta(t))}{\norm{\bTheta(t)} \norm{\nabla_{\bTheta} \mc L(\bTheta(t))}}
    = -1.
\end{align*}
Since the linear tensor network satisfies the technical assumptions in the prior works, we apply these results to our setting and develop a new characterization of the limit directions of the parameters.
Here, we present theorems on separable classification with general linear tensor networks. Corollaries for specific networks are deferred to Appendix~\ref{sec:deferred-corollaries}. 

%We note that our theorems subsume the existing results on linear fully-connected, diagonal, and convolutional networks, while making less convergence assumptions.

%\vspace*{-4pt}
\subsection{Limit directions of parameters are singular vectors}
%\vspace*{-4pt}
\label{sec:meta-thms-clsf-1}
Consider the singular value decomposition (SVD) of a matrix $\mA = \sum_{j=1}^m s_j (\vu_j \otimes \vv_j)$, where $m$ is the rank of $\mA$. Note that the tuples $(\vu_j, \vv_j, s_j)$ are solutions to the system of equations
$s \vu = \mA \vv$ and $s \vv = \mA^T \vu$.
\citet{lim2005singular} generalizes this definition of singular vectors and singular values to higher-order tensors: given an order-$L$ tensor $\tA \in \reals^{k_1 \times \cdots \times k_L}$, we define the singular vectors $\vu_1, \vu_2, \dots, \vu_L$ and singular value $s$ to be the solution of the following system of equations:
\begin{align}
\label{eq:tensorsingvec}
    s \vu_l = \tA \circ (\vu_1, \dots, \vu_{l-1}, \mI_{k_l}, \vu_{l+1}, \dots, \vu_L), \text{ for } l \in [L].
\end{align}
% For order-$L$ tensors ($L\geq 3$), computing such a tensor rank decomposition is NP-hard in general. Similarly, computing the spectral norm (or equivalently, the largest singular value) of a given tensor is also NP-hard.
Using the definition of singular vectors of tensors, we can characterize the limit direction of parameters after reaching 100\% training accuracy. In Appendix~\ref{sec:proof-meta-thm-clsf-1}, we prove the following:
%(the only result where we assume convergence):
\begin{theorem}
\label{thm:meta-thm-clsf-1}
Consider an $L$-layer linear tensor network~\eqref{eq:tensorrep-linear}. 
Assume that the gradient flow satisfies $\mc L(\bTheta(t_0)) < 1$ for some $t_0 \geq 0$ and $\mX^T \vr(t)$ converges in direction, say $\vu^\infty \defeq \lim_{t\to\infty} \frac{\mX^T \vr(t)}{\ltwos{\mX^T\vr(t)}}$.
Then, $\vv_1, \dots, \vv_L$ converge in direction to the singular vectors of $\tM(-\vu^\infty)$.
\end{theorem}
%\vspace*{-5pt}
%For this theorem, we make some convergence assumptions, because the network is fully general; this is the \emph{only} result where we assume convergence.
Theorem~\ref{thm:meta-thm-clsf-1} shows that the limit directions of parameter vectors $\vv_1, \dots, \vv_L$ under gradient flow dynamics must be singular vectors of the data tensor $\tM(-\vu^\infty)$.
For this theorem, we do make some convergence assumptions. First, we assume that the gradient flow finds a parameter $\bTheta(t_0)$ with 100\% training accuracy (i.e., $\mc L(\bTheta(t_0)) < 1$); however, this is because the network is fully general, without any structure to exploit. In the remaining theorems, convergence of loss $\mc L(\bTheta(t)) \to 0$ will be explicitly proven under initial conditions.
The next assumption is that $\mX^T \vr(t)$ converges in direction, which is equivalent to directional convergence of the gradient of $\mc L$ with respect to linear coefficients (also assumed in \citet{gunasekar2018implicit}).
It fact, for the special case of linear fully-connected networks, the directional convergence assumption is \emph{not} required, and the linear coefficients $\vbetafc(\bThetafc)$ converge in direction to the $\ell_2$ max-margin classifier. We state this corollary in Appendix~\ref{sec:cor-meta-thm-clsf-1};
this result appears in~\citet{ji2020directional}, but we provide an alternative proof.

%\vspace*{-4pt}
\subsection{Limit directions in orthogonally decomposable networks}
%\vspace*{-4pt}
\label{sec:meta-thms-clsf-2}
Admittedly, Theorem~\ref{thm:meta-thm-clsf-1} is not a \emph{full} characterization of the limit directions, because there are usually multiple solutions that satisfy~\eqref{eq:tensorsingvec}.
For example, in case of $L=2$, the data tensor $\tM(-\vu^\infty)$ is a matrix and the number of possible limit directions (up to scaling) of $(\vv_1, \vv_2)$ is at least the rank of $\tM(-\vu^\infty)$.
Singular vectors of high order tensors are much less understood than the matrix counterparts, and are much harder to deal with.
Although their existence is implied from the variational formulation \citep{lim2005singular}, they are intractable to compute. Testing if a given number is a singular value, approximating the corresponding singular vectors, and computing the best rank-1 approximation are all NP-hard \citep{hillar2013most}; let alone orthogonal decompositions.

% are both NP-hard \citep{hillar2013most}.
% computing the best rank-1 approximation is also NP-hard \citep{hillar2013most}, 

Given this intractability, it might be reasonable to make some assumptions on the ``structure'' of the data tensor $\tM(\vx)$, so that they are easier to handle. The following assumption defines a class of \emph{orthogonally decomposable} data tensors, which includes \textbf{linear diagonal networks} and \textbf{linear full-length convolutional networks} as special cases (for the proof, see Appendix~\ref{sec:proof-cor-clsf-diag} and \ref{sec:proof-cor-clsf-conv-full}).
\begin{assumption}
\label{assm:complex-diag-decomp}
For the data tensor $\tM(\vx) \in \R^{k_1 \times \cdots \times k_L}$ of a linear tensor network~\eqref{eq:tensorrep-linear}, there exist  a full column rank matrix $\mS \in \C^{m \times d}$ \textup{($d \leq m \leq \min_l k_l$)} and matrices $\mU_1 \in \C^{k_1 \times m}, \dots, \mU_L \in \C^{k_L \times m}$ such that $\mU_l^H \mU_l = \mI_{m}$ for all $l \in [L]$, and the data tensor $\tM(\vx)$ can be written as
%\vspace*{-2pt}
\begin{align}
\label{eq:tensor-orth-decomp}
\tM(\vx) = \sum\nolimits_{j=1}^m [\mS \vx]_j ([\mU_1]_{\cdot,j} \otimes [\mU_2]_{\cdot,j} \otimes \cdots \otimes [\mU_L]_{\cdot,j}
).
\end{align}
% An order-$L$ tensor $\tA \in \R^{k_1 \times \cdots \times k_L}$ is \emph{orthogonally decomposable} with unitary matrices $\mU_1 \in \C^{k_1 \times k_1}, \dots, \mU_L \in \C^{k_L \times k_L}$ if there exists an integer $0 \leq m \leq \min\{k_1, \dots, k_L\}$ and a vector $\vs \in \C^m$ with nonzero entries such that
\end{assumption}
%\vspace*{-5pt}
In this assumption, we allow $\mU_1, \dots, \mU_L$ and $\mS$ to be complex matrices, although $\tM(\vx)$ and parameters $\vv_l$ stay real, as defined earlier.
For a complex matrix $\mA$, we use 
$\mA^*$ to denote its entry-wise complex conjugate,
$\mA^T$ to denote its transpose (without conjugating), and $\mA^H$ to denote its conjugate transpose.
In case of $L = 2$, Assumption~\ref{assm:complex-diag-decomp} requires that the data tensor $\tM(\vx)$ (now a matrix) has singular value decomposition $\tM(\vx) = \mU_1 \diag(\mS \vx) \mU_2^T$; i.e., the left and right singular vectors are independent of $\vx$, and the singular values are linear in $\vx$.
Using Assumption~\ref{assm:complex-diag-decomp}, the following theorem characterizes the limit directions.

% We note that this type of orthogonal decomposition~\eqref{eq:tensor-orth-decomp} does not always exist for general tensors as long as $L > 2$. It is also known that, even without the unitary matrix requirements, computing the tensor rank decomposition for generic tensors is NP-hard (citations needed).

% For example, any matrix ($L=2$) is orthogonally decomposable because they have singular value decomposition.
% Also, $\tMdiag(\vx)$ is orthogonally decomposable.(details req)
% (mention that general order-3 tensors are not orthogonally decomposable)

% We next define an assumption on linear tensor networks whose data tensor $\tT$ is orthogonally decomposable with unitary matrices irrespective of the residual vector. Then, we present a unified theorem applies to networks that satisfy such an assumption.
% \begin{assumption}
% \label{assm:diag-decomp}
% For any residual vector $\vr \in \reals^n$, the data tensor $\tM(\mX^T\vr)$ of the linear tensor network~\eqref{eq:tensorrep-linear} is orthogonally decomposable with fixed unitary matrices $\mU_1, \dots, \mU_L$ independent of $\vr$. Furthermore, the vector $\vs$ of singular values satisfies $\vs = \mS \mX^T \vr$ for some matrix $\mS \in \C^{m \times d}$ such that the transformed dataset $\{(\mS \vx_i,y_i)\}_{i=1}^n$ is linearly separable.
% \end{assumption}
% One sufficient condition under which the dataset $\{(\mS\vx_i,y_i)\}_{i=1}^n$ remains linearly separable after transformation by $\mS$ is that $\mS$ has full column rank. Another sufficient condition is an extremely underdetermined setup, where $n=1$ and $L=2$. In this case, there is only one point $(\vx,y)$ in the dataset and the data tensor $\tM(r \vx) = r\tM(\vx)$ has a fixed direction, only scaling as $r$ changes.
% As we will see later, the diagonal networks and convolutional networks correspond to the cases where $\tM$ satisfies Assumption~\ref{assm:diag-decomp} and $\mS$ is full rank. We will also investigate the cases of $n=1$ and $L=2$ because it allows us to characterize a wider variety of networks, including the convolutional networks with filter size smaller than $d$.

\begin{theorem}
\label{thm:meta-thm-clsf-2}
Suppose a linear tensor network~\eqref{eq:tensorrep-linear} satisfies Assumption~\ref{assm:complex-diag-decomp}. 
If there exists $\lambda > 0$ such that the initial directions $\bar \vv_1, \dots, \bar \vv_L$ of the network parameters satisfy $|[\mU_l^T \bar \vv_l]_j|^2 - |[\mU_L^T \bar \vv_L]_j|^2 \geq \lambda$ for all $l \in [L-1]$ and $j \in [m]$, then the training loss $\mc L(\bTheta(t)) \to 0$. If we additionally assume that $\mX^T \vr(t)$ converges in direction, then $\vbeta(\bTheta(t))$ converges in a direction that aligns with $\mS^T \vrho^\infty$, where $\vrho^\infty \in \C^m$ denotes a stationary point of
\begin{equation*}
    \minimize\nolimits_{\vrho \in \C^m}\quad \norms{\vrho}_{2/L}
    \quad
    \subjectto \quad y_i \vx_i^T\mS^T \vrho \geq 1,~~\forall i\in[n].
\end{equation*}
In case of invertible $\mS$, $\vbeta(\bTheta(t))$ converges in a direction that aligns with a stationary point $\vz^\infty$ of 
\begin{equation*}
    \minimize\nolimits_{\vz \in \R^d}\quad \norms{\mS^{-T} \vz}_{2/L}
    \quad
    \subjectto \quad y_i \vx_i^T \vz \geq 1,~~\forall i\in[n].
\end{equation*}
\end{theorem}
%\vspace*{-4pt}
Theorem~\ref{thm:meta-thm-clsf-2} shows that the gradient flow finds sparse $\vrho^\infty$ that minimizes the $\ell_{2/L}$ norm in the ``singular value space,'' where the data points $\vx_i$ are transformed into vectors $\mS \vx_i$ consisting of singular values of $\tM(\vx_i)$. Also, the proof of Theorem~\ref{thm:meta-thm-clsf-2} reveals that in case of $L=2$, the parameters $\vv_l(t)$ in fact converge in direction to the \emph{top} singular vectors of the data tensor; thus, compared to Theorem~\ref{thm:meta-thm-clsf-1}, we have a more complete characterization of ``which'' singular vectors to converge to. 

The proof of Theorem~\ref{thm:meta-thm-clsf-2} is in Appendix~\ref{sec:proof-meta-thm-clsf-2}. Since the orthogonal decomposition (Assumption~\ref{assm:complex-diag-decomp}) of $\tM(\vx)$ tells us that the singular vectors $\tM(\vx)$ in $\mU_1, \dots, \mU_L$ are independent of $\vx$, we can transform the network parameters $\vv_l$ to $\mU_l^T\vv_l$ and show that the network behaves like a linear diagonal network. This observation comes in handy in the characterization of limit directions.

\begin{remark}[Removal of some convergence assumptions]
%\label{rmk:1}
Recall that one of our aims was to present implicit bias results without relying on various convergence assumptions. Theorem~\ref{thm:meta-thm-clsf-2} extends \citet{gunasekar2018implicit} to a general framework (see Appendix~\ref{sec:cor-meta-thm-clsf-2} for corollaries), while removing the directional convergence assumption on parameters and also removing the assumption that the loss converges to zero by directly proving it under certain initial conditions.
Please note that the theorem does not remove the directional convergence assumption on $\mX^T\vr$. We noticed recently that the ICLR~2021 version of this paper had an erroneous claim that this directional convergence assumption was \emph{not} needed in Theorem~\ref{thm:meta-thm-clsf-2}. In our previous proof, we stated that directional convergence of parameters implies directional convergence of $\mX^T \vr$, which was not fully correct. We leave the removal of this convergence assumption for future work.
\end{remark}

\begin{remark}[Necessity of initialization assumptions]
In removing the assumption on loss, we emphasize that at least some conditions on initialization are \emph{necessary}, because there are examples showing non-convergence of gradient flow for certain initializations \citep{bartlett2018gradient,arora2019convergence}.
The assumptions on $\bar \vv_l$ we pose in Theorem~\ref{thm:meta-thm-clsf-2} are sufficient conditions for the loss $\mc L(\bTheta(t))$ to converge to zero. Due to its sufficiency, the conditions are ``stronger'' than assuming $\mc L(\bTheta(t)) \to 0$; however, they are useful because they can be easily checked \emph{a priori}, i.e., before running gradient flow.
%We also note an important fact that in the implicit bias part (Theorems~\ref{thm:meta-thm-clsf-2}.\ref{item:thm-meta-thm-clsf-2-2}) holds whenever the loss converges to zero, even for initializations that do not satisfy our conditions in Theorems~\ref{thm:meta-thm-clsf-2}.\ref{item:thm-meta-thm-clsf-2-1}.
%This is also the case for other theorems and corollaries that consist of separate convergence and bias parts.
In addition, we argue that our initialization assumptions are not too restrictive; $\lambda$ can be arbitrarily small, so the conditions are satisfied \emph{with probability~1} if we set $\bar \vv_L = \zeros$ and randomly sample other $\bar \vv_l$'s. Setting one layer to zero to prove convergence is also studied in \citet{wu2019global}. In fact, the condition that $\bar \vv_L$ is ``small'' can be replaced with any layer; e.g., convergence still holds if $|[\mU_l^T \bar \vv_l]_j|^2 - |[\mU_1^T \bar \vv_1]_j|^2 \geq \lambda$ for all $l = 2, \dots, L$ and $j \in [m]$.
\end{remark}

%\begin{remark}[Comparison to existing results]
%Theorem~\ref{thm:meta-thm-clsf-2} leads to corollaries (stated in Appendix~\ref{sec:cor-meta-thm-clsf-2}) on linear diagonal and full-length convolutional networks, showing that diagonal (or convolutional) networks converge to the stationary point of the max-margin problem with respect to the $\ell_{2/L}$ norm (or DFT-domain $\ell_{2/L}$ norm).
%Theorem~\ref{thm:meta-thm-clsf-2} recovers the results in \citet{gunasekar2018implicit} without relying on assumptions such as directional convergence of parameters and gradients.
%%Although this was already shown in \citet{gunasekar2018implicit}, note that Theorem~\ref{thm:meta-thm-clsf-2} does not rely on assumptions such as directional convergence of parameters and gradients, which is different from the prior results. 
%\end{remark}

\begin{remark}[Implications to architecture design]
Theorem~\ref{thm:meta-thm-clsf-2} shows that the gradient flow finds a solution that is sparse in a ``transformed'' input space where all data points are transformed with $\mS$. This implies something interesting about architecture design: if the sparsity of the solution under a certain linear transformation $\mT$ is needed, one can design a network using Assumption~\ref{assm:complex-diag-decomp} by setting $\mS = \mT$. Training such a network will give us a solution that has the desired sparsity property.
\end{remark}


% (For a complex matrix $\mA$, we use $\mA^*$ to denote its element-wise conjugate, $\mA^T$ to denote its transpose (without conjugation, and $\mA^H$ to denote its conjugate transpose.)

%The key motivation for Assumption~\ref{assm:complex-diag-decomp} is that under this assumption, the singular vectors of $\tM(-\mX^T \vr)$ are independent of the residual vector $\vr$, which allows us to derive conditions on initialization that ensures convergence to top singular vectors.
%Other than Assumption~\ref{assm:complex-diag-decomp}, there is one more case $\tM(-\mX^T \vr)$ is orthogonally decomposable irrespective of $\vr$: when there is only one data point ($n=1$) and the network is 2-layered ($L=2$). This ``extremely overparametrized'' case is also motivated by the empirical studies in \citep{zhang2019identity}.

\subsection{Limit directions in extremely overparameterized settings}
%\vspace*{-3pt}
Other than Assumption~\ref{assm:complex-diag-decomp}, there is another setting where we can prove a full characterization of limit directions: when there is one data point ($n=1$) and the network is 2-layer ($L=2$). This ``extremely overparameterized'' case is motivated by an experimental paper \citep{zhang2019identity} which studies generalization performance of different architectures when there is only one training data point. Please note that from this theorem onward, we do not require any convergence assumptions.
\begin{theorem}
\label{thm:meta-thm-clsf-3}
Suppose we have a 2-layer linear tensor network~\eqref{eq:tensorrep-linear} and a single data point $(\vx, y)$. Consider the singular value decomposition $\tM(\vx) = \mU_1 \diag(\vs) \mU_2^T$, where $\mU_1 \in \reals^{k_1 \times m}$, $\mU_2 \in \reals^{k_2 \times m}$, and $\vs \in \reals^{m}$ for $m \leq \min\{k_1, k_2\}$.
Let $\vrho^\infty \in \R^m$ be a solution of the following optimization problem
\begin{equation*}
	\minimize\nolimits_{\vrho \in \R^m}\quad \norm{\vrho}_{1}
	\quad
	\subjectto \quad y \vs^T \vrho \geq 1.
\end{equation*}
If there exists $\lambda > 0$ such that the initial directions $\bar \vv_1, \bar \vv_2$ of the network parameters satisfy $[\mU_1^T \bar \vv_1]_j^2 - [\mU_2^T \bar \vv_2]_j^2 \geq \lambda$ for all $j \in [m]$, then the training loss $\mc L(\bTheta(t)) \to 0$. Also, $\vv_1$ and $\vv_2$ converge in direction to $\mU_1 \veta_1^\infty$ and $\mU_2 \veta_2^\infty$, where $\veta_1^\infty, \veta_2^\infty \in \reals^m$ are vectors satisfying $|\veta_1^\infty| = |\veta_2^\infty| = |\vrho^\infty|^{\odot 1/2}$, and $\sign(\veta_1^\infty) = \sign(y)\odot\sign(\veta_2^\infty)$.
\end{theorem}
%\vspace*{-5pt}
The proof of Theorem~\ref{thm:meta-thm-clsf-3} can be found in Appendix~\ref{sec:proof-meta-thm-clsf-3}. Let us parse Theorem~\ref{thm:meta-thm-clsf-3} a bit. Since $\vrho^\infty$ is the minimum $\ell_1$ norm solution in the singular value space, the parameters $\vv_1$ and $\vv_2$ converge in direction to the top singular vectors. We would like to emphasize that this theorem can be applied to \emph{any} network architecture that can be represented as a linear tensor network. For example, recall that the known results on convolutional networks only consider full-length filters ($k_1=d$), hence providing limited insights on networks with small filters, e.g., $k_1=2$. In light of this, we present a corollary in Appendix~\ref{sec:cor-meta-thm-clsf-3} characterizing the convergence directions of linear coefficients for convolutional networks with filter size $k_1 = 1$ and $k_2 = 2$.
% For the special cases of $k_1 = 1$ and $2$, we provide...
% \begin{corollary}
% \label{cor:clsf-conv}
% Suppose we have a 2-layer convolutional network~\eqref{eq:convnet} with filter size $k_1$ and a single data point $(\vx, y)$.
% If if gradient flow satisfies $\mc L(\bThetafc(t_0)) < 1$ for some $t_0 \geq 0$, then $\vbetafc(\bThetafc(t))$ converges in a direction that aligns with the solution of the following optimization problem
% \begin{list}{{\tiny$\bullet$}}{\leftmargin=1.8em}
% \item If $k_1=1$, then $\vbetaconv^\infty = \vx$, which is the minimum $\ell_2$-norm solution. This is \textbf{independent} of the initialization scale $\alpha$ or direction $\bar \vw_1$.
% \item If $k_1=2$,
% \begin{list}{{\tiny$\bullet$}}{\leftmargin=1.8em}
% \item If $\alpha \to \infty$ or $\sum_{i=1}^d [\vx]_i [\vx]_{i+1} = 0$, then
% %\begin{equation*}
%     $[\vbetaconv^\infty]_i \propto \bar \vw_1^T \mc D_i^2(\vx) \bar \vw_1$.
% %\end{equation*}
% \item If $\alpha \to 0$ and \textup{((}$\begin{bmatrix} 1 & 1 \end{bmatrix} \bar \vw_1 \neq 0$ and $\sum_{i=1}^d [\vx]_i [\vx]_{i+1} > 0$\textup{)} or \textup{(}$\begin{bmatrix} 1 & -1 \end{bmatrix} \bar \vw_1 = 0$ and $\sum_{i=1}^d [\vx]_i [\vx]_{i+1} < 0$\textup{))}, then 
% %\begin{equation*}
%     $[\vbetaconv^\infty]_i \propto [\vx]_i + \half [\vx]_{i-1} + \half [\vx]_{i+1}$.
% %\end{equation*}
% \item If $\alpha \to 0$ and \textup{((}$\begin{bmatrix} 1 & -1 \end{bmatrix} \bar \vw_1 \neq 0$ and $\sum_{i=1}^d [\vx]_i [\vx]_{i+1} < 0$\textup{)} or \textup{(}$\begin{bmatrix} 1 & 1 \end{bmatrix} \bar \vw_1 = 0$ and $\sum_{i=1}^d [\vx]_i [\vx]_{i+1} > 0$\textup{))}, then 
% %\begin{equation*}
%     $[\vbetaconv^\infty]_i \propto [\vx]_i - \half [\vx]_{i-1} - \half [\vx]_{i+1}$.
% %\end{equation*} 
% \end{list}
% \end{list}
% \end{corollary}


% We next move on to regression.
%In the next section, we switch gears and study underdetermined regression.
%\vspace*{-5pt}
\section{Implicit bias of gradient flow in underdetermined regression}
%\vspace*{-5pt}
In Section~\ref{sec:clsf}, the limit directions of parameters we characterized do not depend on initialization. This is due to the fact that the parameters diverge to infinity in separable classification problems, so that the initialization becomes unimportant in the limit.
This is not the case in regression setting, because parameters do not diverge to infinity. As we show in this section, the limit points are closely tied to initialization, and our theorems characterize the dependency between them.

%\vspace*{-4pt}
\subsection{Limit point characterization in orthogonally decomposable networks}
%\vspace*{-4pt}
\label{sec:meta-thms-reg}
For the orthogonally decomposable networks satisfying Assumption~\ref{assm:complex-diag-decomp} with real $\mS$ and $\mU_l$'s, we consider how limit points of gradient flow change according to initialization. We consider a specific initialization scheme that, in the special case of diagonal networks, corresponds to setting $\vw_l(0) = \alpha \bar \vw$ for $l \in [L-1]$ and $\vw_L(0) = \zeros$. We use the following lemma on a relevant system of ODEs:
\begin{restatable}{lemma}{lemodereal}
\label{lem:ode-real}
Consider the system of ODEs, where $p, q: \reals \to \reals$:
\begin{align*}
    \dot p = p^{L-2} q, \quad 
    \dot q = p^{L-1}, \quad
    p(0) = 1, \quad
    q(0) = 0.
\end{align*}
Then, the solutions $p_L(t)$ and $q_L(t)$ are continuous on their maximal interval of existence of the form $(-c,c) \subset \reals$ for some $c \in (0, \infty]$. %Moreover, $p_L(t)$ is even and $q_L(t)$ is odd. 
Define $h_L(t) = p_L(t)^{L-1}q_L(t)$; then, $h_L(t)$ is odd and strictly increasing, satisfying $\lim_{t \uparrow c} h_L(t) = \infty$ and $\lim_{t \downarrow -c} h_L(t)= -\infty$.
\end{restatable}

Using the function $h_L(t)$ from Lemma~\ref{lem:ode-real}, we can obtain the following theorem that characterizes the limit points as the minimizer of a norm-like function $Q_{L,\alpha,\bar\veta}$ among the global minima.

\begin{theorem}
\label{thm:meta-thm-reg-1}
Suppose a linear tensor network~\eqref{eq:tensorrep-linear} satisfies Assumption~\ref{assm:complex-diag-decomp}. Assume also that the matrices $\mU_1, \dots, \mU_L$ and $\mS$ from Assumption~\ref{assm:complex-diag-decomp} are all real matrices. 
For some $\lambda > 0$, choose any vector $\bar \veta \in \reals^m$ satisfying $[\bar \veta]_j^2 \geq \lambda$ for all $j \in [m]$, and choose initial directions $\bar \vv_l = \mU_l \bar \veta$ for $l \in [L-1]$ and $\bar \vv_L = \zeros$.
Then, the linear coefficients $\vbeta(\bTheta(t))$ converge to a global minimum $\mS^T \vrho^\infty$, where $\vrho^\infty$ is the solution to
\begin{align*}
	\minimize\nolimits_{\vrho \in \R^m}&\quad Q_{L, \alpha, \bar \veta}(\vrho) 
	\defeq
	\alpha^2  \sum\nolimits_{j=1}^m [\bar \veta]_j^2 H_L \left ( 
	\tfrac{[\vrho]_j}{\alpha^L |[\bar \veta]_j|^L}
	\right )\quad
	\subjectto \quad \mX \mS^T \vrho= \vy,
\end{align*}
where $Q_{L, \alpha, \bar \veta} : \R^m \to \R$ is a norm-like function defined using $H_L (t) \defeq \int_0^t h_L^{-1} (\tau) d\tau$.
In case of invertible $\mS$, $\vbeta(\bTheta(t))$ converges to a global minimum $\vz^\infty$, the solution of
\begin{equation*}
	\minimize\nolimits_{\vz \in \R^d}\quad Q_{L, \alpha, \bar \veta}(\mS^{-T} \vz) 
	\quad
	\subjectto \quad \mX \vz = \vy.
\end{equation*}
\end{theorem}
%\vspace*{-5pt}
The proofs of Lemma~\ref{lem:ode-real} and Theorem~\ref{thm:meta-thm-reg-1} are deferred to Appendix~\ref{sec:proof-meta-thm-reg-1}. 

\begin{remark}[Interpolation between $\ell_1$ and $\ell_2$]
It can be checked that $H_L(t)$ grows like the absolute value function if $t$ is large, and grows like a quadratic function if $t$ is close to zero. This means that 
\begin{align*}
    \lim_{\alpha \to 0} Q_{L, \alpha, \bar \veta}(\vrho) \propto \sum\nolimits_{j=1}^m \tfrac{|[\vrho]_j|}{|[\bar \veta]_j|^{L-2}},\quad
    \lim_{\alpha \to \infty} Q_{L, \alpha, \bar \veta}(\vrho) \propto \sum\nolimits_{j=1}^m \tfrac{[\vrho]_j^2}{[\bar \veta]_j^{2L-2}},
\end{align*}
so $Q_{L, \alpha, \bar \veta}$ interpolates between the weighted $\ell_1$ and weighted $\ell_2$ norms of $\vrho$. Also, the weights in the norm are \emph{dependent} on the initialization direction $\bar \veta$ unless $L=2$ and $\alpha \to 0$. In general, $Q_{L, \alpha, \bar \veta}$ interpolates the standard $\ell_1$ and $\ell_2$ norms only if $|[\bar \veta]_j|$ is the same for all $j \in [m]$. 
This result is similar to the observations made in \citet{woodworth2020kernel} which considers a diagonal network with a ``differential'' structure $f(\vx;\vw_+,\vw_-) = \vx^T(\vw_+^{\odot L}-\vw_-^{\odot L})$. In contrast, our results apply to a more general class of networks, without the need to have the differential structure. In Appendix~\ref{sec:cor-meta-thm-reg-1}, we state corollaries of Theorem~\ref{thm:meta-thm-reg-1} for linear diagonal networks and linear full-length convolutional networks with even data points. There, we also show that deep matrix sensing with commutative sensor matrices~\citep{arora2019implicit} is a special case of our setting. We note that we explicitly prove convergence of loss to zero, instead of assuming it (as done in existing results).
\end{remark}

\subsection{Limit point characterization in extremely overparameterized settings}
Next, we present the regression counterpart of Theorem~\ref{thm:meta-thm-clsf-3}, for 2-layer linear tensor networks~\eqref{eq:tensorrep-linear} with a single data point. For this extremely overparameterized setup, we can fully characterize the limit points as functions of initialization $\vv_1(0) = \alpha \bar \vv_1$ and $\vv_2(0) = \alpha \bar \vv_2$, for \emph{any} linear tensor networks including linear convolutional networks with filter size smaller than input dimension.
\begin{theorem}
\label{thm:meta-thm-reg-2}
Suppose we have a 2-layer linear tensor network~\eqref{eq:tensorrep-linear} and a single data point $(\vx, y)$. Consider the compact SVD $\tM(\vx) = \mU_1 \diag(\vs) \mU_2^T$, where $\mU_1 \in \reals^{k_1 \times m}$, $\mU_2 \in \reals^{k_2 \times m}$, and $\vs \in \reals^{m}$ for $m \leq \min\{k_1, k_2\}$. Assume that there exists $\lambda > 0$ such that the initial directions $\bar \vv_1, \bar \vv_2$ of the network parameters satisfy $[\mU_1^T \bar \vv_1]_j^2 - [\mU_2^T \bar \vv_2]_j^2 \geq \lambda$ for all $j \in [m]$. Then, gradient flow converges to a global minimizer of the loss $\mc L$, and $\vv_1(t)$ and $\vv_2(t)$ converge to the limit points:
\begin{align*}
    \vv_1^\infty \!&=\! \alpha \mU_1\!
    \left ( \mU_1^T\bar \vv_1 \odot \cosh \left (g^{-1} \left (\frac{y}{\alpha^2} \right)\vs\right)
    + \mU_2^T\bar \vv_2 \odot \sinh \left (g^{-1} \left (\frac{y}{\alpha^2} \right) \vs \right)
    \right )
    \!+\!\alpha (\mI_{k_1} - \mU_1 \mU_1^T) \bar \vv_1,\\
    \vv_2^\infty \!&=\! \alpha \mU_2\!
    \left ( \mU_1^T\bar \vv_1 \odot \sinh \left (g^{-1} \left (\frac{y}{\alpha^2} \right) \vs \right)
    +
    \mU_2^T\bar \vv_2 \odot \cosh \left (g^{-1} \left (\frac{y}{\alpha^2} \right)\vs\right)
    \right )
    \!+\!\alpha (\mI_{k_2} - \mU_2 \mU_2^T) \bar \vv_2,
\end{align*}
where $g^{-1}$ is the inverse of the following strictly increasing function
\begin{equation*}
    g(\nu) = \sum\nolimits_{j=1}^m [\vs]_j \Big (\tfrac{[\mU_1^T\bar \vv_1]_j^2+[\mU_2^T\bar \vv_2]_j^2}{2} \sinh(2[\vs]_j \nu) + [\mU_1^T\bar \vv_1]_j [\mU_2^T\bar \vv_2]_j \cosh(2[\vs]_j \nu) \Big).
\end{equation*}
\end{theorem}
The proof can be found in Appendix~\ref{sec:proof-meta-thm-reg-2}. 
We can observe that as $\alpha \to 0$, we have $g^{-1} \left (\frac{y}{\alpha^2} \right) \to \infty$, which results in exponentially faster growth of the $\sinh(\cdot)$ and $\cosh(\cdot)$ for the top singular values. As a result, the top singular vectors dominate the limit points $\vv_1^\infty$ and $\vv_2^\infty$ as $\alpha \to 0$, and the limit points become independent of the initial directions $\bar \vv_1, \bar \vv_2$.
Experiment results in Section~\ref{sec:exp} support this observation.
%The leftmost plot in Figure~\ref{fig:reg} confirms this.
%In contrast, for all $\alpha > 0$, the limit points do depend on initial directions.


%\vspace*{-4pt}
\subsection{Implicit bias in fully-connected networks: the $\alpha \to 0$ limit}
%\vspace*{-4pt}
We state our last theoretical element of this paper, which proves that the linear coefficients $\vbetafc(\bThetafc)$ of deep linear fully-connected networks converge to the minimum $\ell_2$ norm solution as $\alpha \to 0$.
We assume for simplicity that $d_1 = d_2 = \dots = d_L = d$ in this section, but we can extend it for $d_l \geq d$ without too much difficulty, in a similar way as described in \citet{wu2019global}.
Recall $\ffc(\vx; \bThetafc) = \vx^T \mW_1 \cdots \mW_{L-1} \vw_L$. We consider minimizing the training loss $\mc L$ with initialization $\mW_l(0) = \alpha \bar \mW_l$ for $l \in [L-1]$ and $\vw_L(0) = \alpha \bar \vw_L$.

\begin{theorem}
\label{thm:fc-reg}
Consider an $L$-layer linear fully-connected network.
\begin{enumerate}[leftmargin=0.5cm]
	\item \label{item:thm-fc-reg-1} (convergence) If (1)~$\bar \mW_l^T \bar \mW_l \succeq \bar \mW_{l+1} \bar \mW_{l+1}^T$ for $l \in [L-2]$, and (2)~there exists $\lambda > 0$ such that $\bar \mW_{L-1}^T \bar \mW_{L-1} - \bar \vw_L \bar \vw_L^T \succeq \lambda \mI_d$, then the training loss $\mc L(\bThetafc(t)) \to 0$.
	\item \label{item:thm-fc-reg-2} (bias) If $\mc L(\bThetafc(t)) \to 0$ for some fixed initial directions $\bar \mW_1, \dots, \bar \mW_{L-1}, \bar \vw_L$, then
	\begin{equation*}
		\lim_{\alpha \to 0} \lim_{t \to \infty} \vbetafc(\bThetafc(t)) = \mX^T (\mX \mX^T)^{-1} \vy.
	\end{equation*}
\end{enumerate}
\end{theorem}
%\vspace*{-5pt}
The proof is presented in Appendix~\ref{sec:proof-thm-fc-reg}.
Theorem~\ref{thm:fc-reg} shows that in the limit $\alpha \to 0$, linear fully-connected networks have bias towards the minimum $\ell_2$ norm solution, regardless of the depth. This is consistent with the results shown for classification. We note that the convergence part (Theorem~\ref{thm:fc-reg}.\ref{item:thm-fc-reg-1}) holds for any $\alpha > 0$, not necessarily in the limit $\alpha \to 0$. Our sufficient conditions for global convergence stated in Theorem~\ref{thm:fc-reg}.\ref{item:thm-fc-reg-1} is a generalization of the zero-asymmetric initialization scheme ($\bar \mW_1 = \dots = \bar \mW_{L-1} = \mI_d$ and $\bar \vw_L = \zeros$) proposed in \citet{wu2019global}.
We also emphasize that the bias part (Theorem~\ref{thm:fc-reg}.\ref{item:thm-fc-reg-2}) holds for any initial directions that lead to convergence of loss to zero, not just the ones satisfying conditions in Theorem~\ref{thm:fc-reg}.\ref{item:thm-fc-reg-1}.

%\vspace*{-8pt}
\section{Experiments}
%\vspace*{-5pt}
\label{sec:exp}
%We describe our experiments that corroborate our analysis and suggest future research directions. 

\textbf{Regression.~~}
To fully visualize the trajectory of linear coefficients, we run simple experiments with 2-layer linear fully-connected/diagonal/convolutional networks with a single 2-dimensional data point $(\vx,y) = ([1~~2],1)$. For this dataset, the minimum $\ell_2$ norm solution (corresponding to fully-connected networks) of the regression problem is $[0.2~~0.4]$, whereas the minimum $\ell_1$ norm solution (corresponding to diagonal) is $[0~~0.5]$ and the minimum DFT-domain $\ell_1$ norm solution (corresponding to convolutional) is $[0.33~~0.33]$.
We randomly pick four directions $\bar \vz_1, \dots \bar \vz_4 \in \R^2$, and choose initial directions of the network parameters in a way that their linear coefficients at initialization are exactly $\vbeta(\bTheta(0)) = \alpha^2 \bar \vz_j$. With varying initial scales $\alpha \in \{0.01, 0.5, 1\}$, we run GD with small step size $\eta = 10^{-3}$ for large enough number of iterations $T = 5 \times 10^3$. Figures~\ref{fig:reg} and \ref{fig:clsf} plot the trajectories of $\vbeta(\bTheta)$ (appropriately clipped for visual clarity) as well as the predicted limit points (Theorem~\ref{thm:meta-thm-reg-2}). We observe that even though the networks start at the same linear coefficients $\alpha^2 \bar \vz_j$, they evolve differently due to different architectures. Note that the prediction of limit points is accurate, and the solution found by GD is less dependent on initial directions when $\alpha$ is small.

\textbf{Classification.~~}
It is shown in the existing works as well as in Section~\ref{sec:clsf} that the limit directions of linear coefficients are independent of the initialization. Is this also true in practice? To see this, we run a set of toy experiments on classification with two data points $(\vx_1,y_1) = ([1~~2],+1)$ and $(\vx_2,y_2) = ([0~-\!3],-1)$. One can check that the max-margin classifiers for this problem are in the same directions to the corresponding min-norm solutions in the regression problem above. We use the same networks as in regression, and the same set of initial directions satisfying $\vbeta(\bTheta(0)) = \alpha^2 \bar \vz_j$. With initial scales $\alpha \in \{0.01, 0.5, 1\}$, we run GD with step size $\eta = 5 \times 10^{-4}$ for $T = 2 \times 10^6$ iterations. All experiments reached $\mc L(\bTheta) \lesssim 10^{-5}$ at the end. The trajectories are plotted in Figure~\ref{fig:clsf} in the Appendix. We find that, in contrast to our theoretical characterization, the actual coefficients are quite dependent on initialization, because we do not train the network all the way to zero loss. This observation is also consistent with a recent analysis \citep{moroshko2020implicit} for diagonal networks, and suggests that understanding the behavior of iterates after a finite number of steps is an important future work.
%\vspace*{-5pt}
\section{Conclusion}
%\vspace*{-5pt}
This paper studies the implicit bias of gradient flow on training linear tensor networks. Under a general tensor formulation of linear networks, we provide theorems characterizing how the network architectures and initializations affect the limit directions/points of gradient flow.
Our work provides a unified framework that connects multiple existing results on implicit bias of gradient flow as special cases.
%We hope our unifying framework will be a stepping stone for future research.

\bibliographystyle{iclr2021_conference}
\bibliography{cite}

\newpage
\appendix

\begin{figure}[t]
     \centering
    %  \begin{subfigure}[t]{0.32\textwidth}
    %      \centering
    %      \includegraphics[width=\textwidth]{reg1_crop.png}
    %  \end{subfigure}
    %  \begin{subfigure}[t]{0.32\textwidth}
    %      \centering
    %      \includegraphics[width=\textwidth]{reg2_crop.png}
    %  \end{subfigure}
    %  \begin{subfigure}[t]{0.32\textwidth}
    %      \centering
    %      \includegraphics[width=\textwidth]{reg3_crop.png}
    %  \end{subfigure}
     \begin{subfigure}[t]{0.48\textwidth}
         \centering
         \includegraphics[width=\textwidth]{regression_bias_a_5en1.png}
     \end{subfigure}
     ~
     \begin{subfigure}[t]{0.48\textwidth}
         \centering
         \includegraphics[width=\textwidth]{classification_bias_a_1en2.png}
     \end{subfigure}
     
     \vspace{5pt}
     
     \begin{subfigure}[t]{0.48\textwidth}
         \centering
         \includegraphics[width=\textwidth]{classification_bias_a_5en1.png}
     \end{subfigure}
     ~
     \begin{subfigure}[t]{0.48\textwidth}
         \centering
         \includegraphics[width=\textwidth]{classification_bias_a_1e0.png}
     \end{subfigure}
    
    \caption{Gradient descent trajectories of linear coefficients of linear fully-connected, diagonal, and convolutional networks on a regression task with initial scale $\alpha = 0.5$ (top left), and networks on a classification task with initial scales $\alpha = 0.01, 0.5, 1$ (rest).
    Networks are initialized at the same coefficients (circles on purple lines), but follow different trajectories due to different implicit biases of networks induced from their architecture. 
    The top left figure shows that our theoretical predictions on limit points (circles on yellow line, the set of global minima) agree with the solution found by GD.
    For details of the experimental setup, please refer to Section~\ref{sec:exp}. }
    \label{fig:clsf}
\end{figure}

\section{Corollaries on specific network architectures}
\label{sec:deferred-corollaries}
We present corollaries obtained by specializing the theorems in the main text to specific network architectures. We briefly review the linear neural network architectures studied in this section.

\textbf{Linear fully-connected networks.}~~
An $L$-layer linear fully-connected network is defined as
\begin{equation}
\label{eq:fcnet-linear}
    \ffc (\vx; \bThetafc) = \vx^T \mW_1 \cdots \mW_{L-1} \vw_L,
\end{equation}
where $\mW_l \in \reals^{d_{l} \times d_{l+1}}$ for $l \in [L-1]$ (we use $d_1 = d$) and $\vw_L \in \reals^{d_{L}}$.

\textbf{Linear diagonal networks.}~~
An $L$-layer linear diagonal network is written as
\begin{align}
\label{eq:diagnet-linear}
    \fdiag (\vx; \bThetadiag) = (\vx \odot \vw_1 \odot \cdots \odot \vw_{L-1})^T \vw_L,
\end{align}
where $\vw_l \in \reals^d$ for $l \in [L]$. 

\textbf{Linear (circular) convolutional networks.}~~
An $L$-layer linear convolutional network is written as 
\begin{align}
\label{eq:convnet-linear}
    \fconv(\vx; \bThetaconv) = (\cdots((\vx \star \vw_1) \star \vw_2 ) \cdots  \star \vw_{L-1})^T \vw_L,
\end{align}
where $\vw_l \in \reals^{k_l}$ with $k_l \leq d$ and $k_L = d$, and $\star$ defines the circular convolution: for any $\va \in \reals^d$ and $\vb \in \reals^k$ ($k \leq d$), we have $\va \star \vb \in \reals^d$ defined as
$[\va \star \vb]_i = \sum_{j=1}^k [\va]_{(i+j-1) \bmod d} [\vb]_j$, for $i \in [d]$.
In case of $k_l = d$ for all $l \in [L]$, we refer to this network as \emph{full-length} convolutional networks.

\textbf{Deep matrix sensing.}~~
The deep matrix sensing problem considered in \citet{gunasekar2017implicit,arora2019implicit} aims to minimize the following problem
\begin{equation}
\label{eq:matrix-sensing}
    \minimize_{\mW_1, \dots, \mW_L \in \reals^{d \times d}}
    \quad
    \mc L_{\rm ms}(\mW_1 \cdots \mW_L) \defeq
    \sum\nolimits_{i=1}^n (y_i - \< \mA_i, \mW_1\cdots \mW_L \> )^2,
\end{equation}
where the sensor matrices $\mA_1, \dots, \mA_n \in \reals^{d \times d}$ are symmetric.
Following \citet{gunasekar2017implicit,arora2019implicit}, we consider sensor matrices $\mA_1, \dots, \mA_n \in \reals^{d \times d}$ that commute. To make the problem underdetermined, we assume that $n \leq d$, and $\mA_i$'s are linearly independent.

\subsection{Corollary of Theorem~\ref{thm:meta-thm-clsf-1}}
\label{sec:cor-meta-thm-clsf-1}
\begin{corollary}
\label{cor:clsf-fc}
Consider an $L$-layer linear fully-connected network~\eqref{eq:fcnet-linear}. If the training loss satisfies $\mc L(\bThetafc(t_0)) < 1$ for some $t_0 \geq 0$, then $\vbetafc(\bThetafc(t))$ converges in a direction that aligns with the solution of the following optimization problem
\begin{equation*}
    \minimize\nolimits_{\vz \in \reals^d} \quad \ltwo{\vz}^2 
    \quad
    \subjectto \quad y_i \vx_i^T \vz \geq 1,~~\forall i\in[n].
\end{equation*}
\end{corollary}
Corollary~\ref{cor:clsf-fc} shows that whenever the network separates the data correctly, the direction of linear coefficients $\vbetafc(\bThetafc)$ of linear fully-connected networks converges to the $\ell_2$ max-margin classifier.
Note that this corollary does not require the directional convergence of $\mX^T\vr$, which is different from Theorem~\ref{thm:meta-thm-clsf-1}. In fact, this corollary also appears in~\citet{ji2020directional}, but we provide an alternative proof based on our tensor formulation.
The proof of Corollary~\ref{cor:clsf-fc} can be found in Appendix~\ref{sec:proof-meta-thm-clsf-1}.
%There are many papers studying convergence of gradient flow/descent on fully connected networks, so this corollary can be combined with any of these results to prove that linear fully-connected networks converge in direction to the $\ell_2$ max-margin classifier.

\subsection{Corollaries of Theorem~\ref{thm:meta-thm-clsf-2}}
\label{sec:cor-meta-thm-clsf-2}
Theorem~\ref{thm:meta-thm-clsf-2} leads to corollaries on linear diagonal and full-length convolutional networks, showing that diagonal (or convolutional) networks converge to the stationary point of the max-margin problem with respect to the $\ell_{2/L}$ norm (or DFT-domain $\ell_{2/L}$ norm).
We state the corollary on linear diagonal networks below:
\begin{corollary}
\label{cor:clsf-diag}
Consider an $L$-layer linear diagonal network~\eqref{eq:diagnet-linear}.
If there exists $\lambda > 0$ such that the initial directions $\bar \vw_1, \dots, \bar \vw_L$ of the network parameters satisfy $[\bar \vw_l]_j^2 - [\bar \vw_L]_j^2 \geq \lambda$ for all $l \in [L-1]$ and $j \in [d]$, then the training loss $\mc L(\bThetadiag(t)) \to 0$. If we additionally assume that $\mX^T \vr(t)$ converges in direction, then $\vbetadiag(\bThetadiag(t))$ converges in a direction that aligns with a stationary point $\vz^\infty$ of 
\begin{equation*}
    \minimize\nolimits_{\vz \in \R^d}\quad \norms{\vz}_{2/L}
    \quad
    \subjectto \quad y_i \vx_i^T \vz \geq 1,~~\forall i\in[n].
\end{equation*}
\end{corollary}
For the corollary on full-length convolutional networks, we define $\mF \in \C^{d \times d}$ to be the matrix of discrete Fourier transform basis $[\mF]_{j,k} = \frac{1}{\sqrt{d}} \exp(-\frac{\sqrt{-1} \cdot 2 \pi (j-1)(k-1)}{d})$. Note that $\mF^* = \mF^{-1}$, and both $\mF$ and $\mF^*$ are symmetric, but not Hermitian.
\begin{corollary}
\label{cor:clsf-conv-full}
Consider an $L$-layer linear full-length convolutional network~\eqref{eq:convnet-linear}. 
If there exists $\lambda > 0$ such that the initial directions $\bar \vw_1, \dots, \bar \vw_L$ of the network parameters satisfy $|[\mF \bar \vw_l]_j|^2 - |[\mF \bar \vw_L]_j|^2 \geq \lambda$ for all $l \in [L-1]$ and $j \in [d]$, then the training loss $\mc L(\bThetaconv(t)) \to 0$. If we additionally assume that $\mX^T \vr(t)$ converges in direction, then $\vbetaconv(\bThetaconv(t))$ converges in a direction that aligns with a stationary point $\vz^\infty$ of 
\begin{equation*}
    \minimize\nolimits_{\vz \in \R^d}\quad \norms{\mF \vz}_{2/L}
    \quad
    \subjectto \quad y_i \vx_i^T \vz \geq 1,~~\forall i\in[n].
\end{equation*}
\end{corollary}
Corollary~\ref{cor:clsf-diag} shows that in the limit, linear diagonal network finds a sparse solution $\vz$ that is a stationary point of the $\ell_{2/L}$ max-margin classification problem. Corollary~\ref{cor:clsf-conv-full} has a similar conclusion except that the standard $\ell_{2/L}$ norm is replaced with DFT-domain $\ell_{2/L}$ norm. These corollaries remove two of the convergence assumptions required in \citet{gunasekar2018implicit}.
The proofs of Corollaries~\ref{cor:clsf-diag} and \ref{cor:clsf-conv-full} are in Appendix~\ref{sec:proof-meta-thm-clsf-2}.

\subsection{Corollary of Theorem~\ref{thm:meta-thm-clsf-3}}
\label{sec:cor-meta-thm-clsf-3}

Recall that Theorem~\ref{thm:meta-thm-clsf-3} can be applied to any 2-layer networks that can be represented as linear tensor networks. Examples include the convolutional networks that are not full-length (i.e., filter size $k_1 < d$), which are not covered by the previous result \citep{gunasekar2018implicit}.
Here, we present the characterization of convergence directions of $\vbetaconv(\bThetaconv(t))$ for 2-layer linear convolutional networks, with filter size $k_1 = 1$ and $k_1 = 2$.
\begin{corollary}
\label{cor:clsf-conv-small}
Consider a 2-layer linear convolutional network~\eqref{eq:convnet-linear} with $k_1 = 1$ and a single data point $(\vx, y)$. If there exists $\lambda > 0$ such that the initial directions $\bar \vw_1$ and $\bar \vw_2$ of the network parameters satisfy $\norm{\vx}^2\bar \vv_1^2 - (\vx^T \bar \vv_2)^2 \geq \norm{\vx}^2 \lambda$, then the training loss $\mc L(\bThetaconv(t)) \to 0$. Also, $\vbetaconv(\bThetaconv(t))$ converges in direction that aligns with $y \vx$.

Consider a 2-layer linear convolutional network~\eqref{eq:convnet-linear} with $k_1 = 2$ and a single data point $(\vx, y)$. Let $\overleftarrow \vx \defeq \begin{bmatrix} [\vx]_2 & \cdots & [\vx]_d & [\vx]_1 \end{bmatrix}$, and $\overrightarrow \vx \defeq \begin{bmatrix} [\vx]_d & [\vx]_1 & \cdots & [\vx]_{d-1} \end{bmatrix}$. If there exists $\lambda > 0$ such that the initial directions $\bar \vw_1$ and $\bar \vw_2$ of the network parameters satisfy 
\begin{align*}
	([\bar \vv_1]_1 + [\bar \vv_1]_2 )^2 - \frac{((\vx+\overleftarrow \vx)^T \bar \vv_2)^2}{\ltwos{\vx}^2 + \vx^T \overleftarrow \vx} \geq \lambda,~~\text{ and }~~
	([\bar \vv_1]_1 - [\bar \vv_1]_2 )^2 - \frac{((\vx-\overleftarrow \vx)^T \bar \vv_2)^2}{\ltwos{\vx}^2 - \vx^T \overleftarrow \vx} \geq \lambda,
\end{align*}
then the training loss $\mc L(\bThetaconv(t)) \to 0$. Also, $\vbetaconv(\bThetaconv(t))$ converges in a direction that aligns with a ``filtered'' version of $\vx$:
\begin{equation*}
	\lim_{t \to \infty} \frac{\vbetaconv(\bThetaconv(t))}{\ltwos{\vbetaconv(\bThetaconv(t))}} \propto
	\begin{cases}
		2y\vx + y\overleftarrow \vx + y\overrightarrow \vx & \text{ if } \vx^T \overleftarrow \vx > 0,\\
		2y\vx - y\overleftarrow \vx - y\overrightarrow \vx & \text{ if } \vx^T \overleftarrow \vx < 0.
	\end{cases}
\end{equation*}
\end{corollary}
Corollary~\ref{cor:clsf-conv-small} shows that if the filter size is $k_1 = 1$, then the limit direction of $\vbetaconv(\bThetaconv)$ is the $\ell_2$ max-margin classifier. Note that this is quite different from the case $k_1 = d$ which converges to the DFT-domain $\ell_1$ max-margin classifier.
However, for $1 < k_1 < d$, it is difficult to characterize the limit direction as the max-margin classifier of some commonly-used norms. Rather, the limit directions of $\vbetaconv(\bThetaconv)$ correspond to a ``filtered'' version of the data point, and the weights of the filter depend on the data point $\vx$. For $k_1 = 2$, the filter is a low-pass filter if the autocorrelation $\vx^T \overleftarrow \vx$ of $\vx$ is positive, and high-pass if the autocorrelation is negative. For $k_1 > 2$, the filter weights are more complicated to characterize in terms of $\vx$, and the ``filter length'' increases as $k_1$ increases. We prove Corollary~\ref{cor:clsf-conv-small} in Appendix~\ref{sec:proof-meta-thm-clsf-3}. For a more detailed investigation on the 2-layer convolutional network settings, see \citet{jagadeesan2021inductive}.

\subsection{Corollaries of Theorem~\ref{thm:meta-thm-reg-1}}
\label{sec:cor-meta-thm-reg-1}

To illustrate the versatility of Theorem~\ref{thm:meta-thm-reg-1}, we state its corollaries for three different settings: linear diagonal networks, linear full-length convolutional networks with even data, and deep matrix sensing with commutative sensor matrices. The proofs of the corollaries can be found in Appendix~\ref{sec:proof-meta-thm-reg-1}.
\begin{corollary}
\label{cor:reg-diag}
Consider an $L$-layer linear diagonal network~\eqref{eq:diagnet-linear}.
For some $\lambda > 0$, choose any vector $\bar \vw \in \reals^d$ satisfying $[\bar \vw]_j^2 \geq \lambda$ for all $j \in [d]$, and choose initial directions $\bar \vw_l = \bar \vw$ for $l \in [L-1]$ and $\bar \vw_L = \zeros$.
Then, the linear coefficients $\vbetadiag(\bThetadiag(t))$ converge to a global minimum $\vz^\infty$, which is the solution of
\begin{equation*}
    \minimize\nolimits_{\vz \in \R^d}\quad Q_{L, \alpha, \bar \vw}(\vz)
    \defeq
    \alpha^2  \sum\nolimits_{j=1}^d [\bar \vw]_j^2 H_L \left ( 
    \tfrac{[\vz]_j}{\alpha^L |[\bar \vw]_j|^L}
    \right )
    \quad
    \subjectto \quad \mX \vz = \vy.
\end{equation*}
\end{corollary}

Recall that the statement of Assumption~\ref{assm:complex-diag-decomp} allows the matrices $\mS, \mU_1, \dots, \mU_L$ to be complex, but Theorem~\ref{thm:meta-thm-reg-1} poses another assumption that these matrices are real.
In applying Theorem~\ref{thm:meta-thm-clsf-2} to convolutional networks to get Corollary~\ref{cor:clsf-conv-full}, we used the fact that the data tensor $\tMconv(\vx)$ of a linear full-length convolutional network satisfies Assumption~\ref{assm:complex-diag-decomp} with $\mS = d^{\frac{L-1}{2}} \mF$ and $\mU_1 = \dots = \mU_L = \mF^*$, where $\mF \in \C^{d \times d}$ is the matrix of discrete Fourier transform basis $[\mF]_{j,k} = \frac{1}{\sqrt{d}} \exp(-\frac{\sqrt{-1} \cdot 2 \pi (j-1)(k-1)}{d})$ and $\mF^*$ is the complex conjugate of $\mF$.
Note that these are complex matrices, so one cannot directly apply Theorem~\ref{thm:meta-thm-reg-1} to convolutional networks. However, it turns out that if the data and initialization are even, we can derive a corollary for convolutional networks.


We say that a vector is \emph{even} when it satisfies the even symmetry, as in even functions. More concretely,
a vector $\vx \in \reals^d$ is even if $[\vx]_{j+2} = [\vx]_{d-j}$ for $j = 0, \dots, \lfloor \frac{d-3}{2} \rfloor$; i.e., the vector has the even symmetry around its ``origin'' $[\vx]_1$. From the definition of the matrix $\mF \in \C^{d \times d}$, it is straightforward to check that if $\vx$ is real and even, then its DFT $\mF \vx$ is also real and even (see Appendix~\ref{sec:proof-reg-conv-full} for details).

\begin{corollary}
\label{cor:reg-conv-full}
Consider an $L$-layer linear full-length convolutional network~\eqref{eq:convnet-linear}. Assume that the data points $\{\vx_i\}_{i=1}^n$ are all even. For some $\lambda > 0$, choose any even vector $\bar \vw$ satisfying $[\mF \bar \vw]_j^2 \geq \lambda$ for all $j \in [d]$, and choose initial directions $\bar \vw_l = \bar \vw$ for $l \in [L-1]$ and $\bar \vw_L = \zeros$. Then, the linear coefficients $\vbetaconv(\bThetaconv(t))$ converge to a global minimum $\vz^\infty$, which is the solution of
\begin{equation*}
    \minimize_{\vz \in \R^d,\textup{ even}}\quad Q_{L, \alpha, \mF \bar \vw}(\mF \vz)
    \defeq
    \alpha^2 \sum\nolimits_{j=1}^d [\mF \bar \vw]_j^2 H_L \left ( 
    \tfrac{[\mF \vz]_j}{\alpha^L |[\mF \bar \vw]_j|^L}
    \right )
    \quad
    \subjectto \quad \mX \vz = \vy.
\end{equation*}
\end{corollary}
Corollaries~\ref{cor:reg-diag} and \ref{cor:reg-conv-full} show that the interpolation between minimum weighted $\ell_1$ and weighted $\ell_2$ solutions occurs for diagonal networks, and also for convolutional networks (in DFT domain, with the restriction of even symmetry).
The conclusion of Corollary~\ref{cor:reg-diag} is similar to the results in \citet{woodworth2020kernel}, but the network architecture~\eqref{eq:diagnet-linear} considered in our corollary is different from the ``differential'' network $f(\vx;\vw_+,\vw_-) = \vx^T(\vw_+^{\odot L}-\vw_-^{\odot L})$ in \citet{woodworth2020kernel}.

As mentioned in the main text, we can actually show that the matrix sensing result in \citet{arora2019implicit} is a special case of our Theorem~\ref{thm:meta-thm-reg-1}.
Given any symmetric matrix $\mM \in \reals^{d \times d}$, let $\eig(\mM) \in \reals^d$ be the $d$-dimensional vector containing the eigenvalues of $\mM$.
\begin{corollary}
\label{cor:reg-matrix-fac}
Consider the depth-$L$ deep matrix sensing problem~\eqref{eq:matrix-sensing}. Let $\mA_i$'s be symmetric, and assume that $\mA_1, \dots, \mA_n$ commute. For $\alpha > 0$, choose initialization $\mW_l(0) = \alpha \mI_d$ for $l \in [L-1]$ and $\mW_L(0) = \zeros$. Then, the product $\mW_1(t)\cdots\mW_L(t)$ converges to the solution $\mM^\infty$ of
\begin{equation*}
    \minimize_{\mM \in \R^{d \times d}, \textup{ symmetric}}\quad Q_{L, \alpha}(\eig(\mM))
    \defeq
    \alpha^2  \sum\nolimits_{j=1}^d H_L \left ( 
    \tfrac{[\eig(\mM)]_j}{\alpha^L}
    \right )
    \quad
    \subjectto
    \quad
    \mc L_{\rm ms}(\mM) = 0.
\end{equation*}
\end{corollary}
Under an additional assumption that $\mA_i$'s are positive semidefinite, Theorem~2 in \citet{arora2019implicit} studies the initialization $\mW_l(0) = \alpha \mI_d$ for all $l \in [L]$, and shows that the limit point of $\mW_1 \cdots \mW_L$ converges to the minimum nuclear norm solution as $\alpha \to 0$. We remove the assumption of positive definiteness of $\mA_i$'s and let $\mW_L(0) = \zeros$, to show a complete characterization of the solution found by gradient flow, which interpolates between the minimum nuclear norm (i.e., Schatten 1-norm) solution (when $\alpha \to 0$) and the minimum Frobenius norm (i.e., Schatten 2-norm) solution (when $\alpha \to \infty$).

\section{Tensor representation of fully-connected networks}
\label{sec:tensorform-ex}
In Section~\ref{sec:tensorform}, we only defined the data tensor $\tMfc(\vx)$ of fully-connected networks for $L=2$. Here, we describe an iterative procedure constructing the data tensor for deep fully-connected networks.

We start with $\tT_1(\vx) \defeq \vx \in \reals^{d_1}$. Next, define a block diagonal matrix $\tT_2(\vx) \in \reals^{d_1 d_2 \times d_2}$ where the ``diagonals'' $[\tT_2(\vx)]_{d_1(j-1)+1:d_1 j,j} \defeq \tT_1(\vx)$ for $j\in[d_2]$, while all the other entries are filled with 0. We continue this ``block diagonal'' procedure, as the following. Having defined $\tT_{l-1}(\vx) \in \reals^{d_1d_2 \times \cdots \times d_{l-2}d_{l-1} \times d_{l-1}}$,
\begin{enumerate}
    \item Define $\tT_{l}(\vx) \in \reals^{d_1d_2 \times \cdots \times d_{l-1}d_{l} \times d_{l}}$.
    \item Set $[\tT_l(\vx)]_{\cdot,\dots,\cdot,d_{l-1}(j-1)+1:d_{l-1}j,j} \defeq \tT_{l-1}(\vx), \forall j \in [d_l]$.
    \item Set all the remaining entries of $\tT_l(\vx)$ to zero.
\end{enumerate}
We iterate this process for $l = 2, \dots, L$, and set $\tMfc(\vx) \defeq \tT_{L}(\vx)$.
By defining the parameters of the tensor formulation $\vv_l = \vect(\mW_l)$ for $l \in [L-1]$ and $\vv_L = \vw_L$, and using the tensor $\tM(\vx) = \tMfc(\vx)$, we can show the equivalence of \eqref{eq:tensorrep} and \eqref{eq:fcnet}.


% \begin{equation}
% \label{eq:fcnet-tensor}
% \begin{aligned}
% &\tT_1(\vx) \defeq \vx \in \reals^{d_1}\\
% &\tT_l(\vx) \in \reals^{d_1 d_2 \times \cdots \times d_{l-1} d_l \times d_l},
% [\tT_l(\vx)]_{\cdot,\dots,\cdot,d_{l-1}(j-1)+1:d_{l-1} j,j} = \tT_{l-1}(\vx), \forall j \in [d_l],
% \text{ all the remaining entries of $\tT_l(\vx)$ filled with zero.} 
% \tMfc(\vx) = \tT_L(\vx).
% \end{aligned}
% \end{equation}


\section{Proofs of Theorem~\ref{thm:meta-thm-clsf-1} and Corollary~\ref{cor:clsf-fc}}
\label{sec:proof-meta-thm-clsf-1}
\subsection{Proof of Theorem~\ref{thm:meta-thm-clsf-1}}
%If the gradient flow satisfies $\mc L(\bTheta(t_0)) < 1$ for some $t_0 \geq 0$ and $\mX^T \vr(t)$ converges in direction, then $\vv_1, \dots, \vv_L$ converge to the singular vectors of $\tM(-\lim_{t\to\infty}\nicefrac{\mX^T \vr(t)}{\norms{\mX^T \vr(t)}})$.

The proof of Theorem~\ref{thm:meta-thm-clsf-1} is outlined as follows. First, using the directional convergence and alignment results in \citet{ji2020directional}, we prove that each of our network parameters $\vv_l$ converges in direction, and it aligns with its corresponding negative gradient $-\nabla_{\vv_l} \mc L$. Then, we prove that the directions of $\vv_l$'s are actually singular vectors of $\tM(-\vu^\infty)$, where $\vu^\infty \defeq \lim_{t\to\infty} \frac{\mX^T \vr(t)}{\ltwos{\mX^T\vr(t)}}$.

Since a linear tensor network is an $L$-homogeneous polynomial of $\vv_1, \dots, \vv_L$, it satisfies the assumptions required for Theorems~3.1 and 4.1 in \citet{ji2020directional}. 
These theorems imply that if the gradient flow satisfies $\mc L(\bTheta(t_0)) < 1$ for some $t_0 \geq 0$, then $\bTheta(t)$ converges in direction, and the direction aligns with $-\nabla_{\bTheta} \mc L(\bTheta(t))$; that is,
\begin{equation}
\label{eq:dirconv}
    \lim_{t\to\infty} \ltwos{\bTheta(t)} = \infty,
    ~~
    \lim_{t\to\infty}\frac{\bTheta(t)}{\ltwos{\bTheta(t)}} = \bTheta^\infty,
    ~~
    \lim_{t\to\infty} \frac{\bTheta(t)^T \nabla_{\bTheta} \mc L(\bTheta(t))}{\ltwos{\bTheta(t)} \ltwos{\nabla_{\bTheta} \mc L(\bTheta(t))}}
    = -1.
\end{equation}
For linear tensor networks~\eqref{eq:tensorrep-linear}, the parameter $\bTheta$ is the concatenation of all parameter vectors $\vv_1, \dots, \vv_L$, so \eqref{eq:dirconv} holds for $\bTheta = \begin{bmatrix} \vv_1^T & \dots & \vv_L^T \end{bmatrix}^T$.

Now, recall that by the definition of the linear tensor network, we have the following gradient flow dynamics
\begin{align*}
    \dot \vv_l
    = \tM(-\mX^T\vr) \circ (\vv_1, \dots, \vv_{l-1}, \mI_{k_l}, \vv_{l+1}, \dots, \vv_L).
\end{align*}
Note that we can apply this to calculate the rate of growth of $\ltwos{\vv_l}^2$:
\begin{align*}
    \frac{d}{dt} \ltwos{\vv_l}^2
    &= 2 \vv_l^T \dot \vv_l
    = 2\vv_l^T \tM(-\mX^T\vr) \circ (\vv_1, \dots, \vv_{l-1}, \mI_{k_l}, \vv_{l+1}, \dots, \vv_L)\\
    &= 2 \tM(-\mX^T\vr) \circ (\vv_1, \dots, \vv_{l-1}, \vv_{l}, \vv_{l+1}, \dots, \vv_L)\\
    &= \frac{d}{dt} \ltwos{\vv_{l'}}^2 \quad\text{ for any $l' \in [L]$, }
\end{align*}
so the rate at which $\ltwos{\vv_l}^2$ grows over time is the same for all layers $l \in [L]$. By the definition of $\bTheta$ and \eqref{eq:dirconv}, we have 
\begin{equation*}
    \ltwos{\bTheta}^2 = \sum_{l=1}^L \ltwos{\vv_l}^2 \to \infty,
\end{equation*}
which then implies
\begin{equation*}
    \lim_{t\to\infty} \ltwos{\vv_l(t)} \to \infty,~~
    \lim_{t\to\infty} \frac{\ltwos{\bTheta(t)}}{\ltwos{\vv_l(t)}}
    = \sqrt{\frac{\ltwos{\bTheta(t)}^2}{\ltwos{\vv_l(t)}^2}}
    = \sqrt{L},
\end{equation*}
for all $l \in [L]$.
Now, let $\mc I_l$ be the set of indices that correspond to the components of $\vv_l$ in $\bTheta$. 
It follows from~\eqref{eq:dirconv} that
\begin{align*}
    \lim_{t \to \infty}
    \frac{\vv_l(t)}{\ltwos{\vv_l(t)}}
    =
    \lim_{t \to \infty}
    \frac{\vv_l(t)}{\ltwos{\bTheta(t)}}
    \frac{\ltwos{\bTheta(t)}}{\ltwos{\vv_l(t)}}
    =
    \lim_{t \to \infty}
    \frac{[\bTheta(t)]_{\mc I_l}}{\ltwos{\bTheta(t)}}
    \frac{\ltwos{\bTheta(t)}}{\ltwos{\vv_l(t)}}
    =
    \sqrt{L} [\bTheta^\infty]_{\mc I_l},
\end{align*}
thus showing the directional convergence of $\vv_l$'s.

Next, it follows from directional convergence of $\bTheta$ and the fact that it aligns with $-\nabla_{\bTheta} \mc L(\bTheta)$~\eqref{eq:dirconv} that $\nabla_{\bTheta} \mc L(\bTheta)$ also converges in direction, in the opposite direction of $\bTheta$. By comparing the components in $\mc I_l$'s, we get that $\nabla_{\vv_l} \mc L(\bTheta)$ converges in the opposite direction of $\vv_l$.

For any $l\in[L]$, now let $\vv_l^\infty\defeq \lim_{t \to \infty} \frac{\vv_l(t)}{\ltwos{\vv_l(t)}}$. Also recall the assumption that $\mX^T\vr(t)$ converges in direction, to a unit vector $\vu^\infty \defeq \lim_{t\to\infty} \frac{\mX^T \vr(t)}{\ltwos{\mX^T\vr(t)}}$. By the gradient flow dynamics of $\vv_l$, we have
\begin{align*}
    \vv_l^\infty
    &\propto -\nabla_{\vv_l} \mc L(\bTheta^\infty)
    = \tM(-\vu^\infty) \circ (\vv_1^\infty, \dots, \vv_{l-1}^\infty, \mI_{k_l}, \vv_{l+1}^\infty, \dots, \vv_L^\infty),
\end{align*}
for all $l \in [L]$. Note that this equation has the same form as~\eqref{eq:tensorsingvec}, the definition of singular vectors in tensors. So this proves that $(\vv_1^\infty, \dots, \vv_L^\infty)$ are singular vectors of $\tM(-\vu^\infty)$.

\subsection{Proof of Corollary~\ref{cor:clsf-fc}}
\label{sec:proof-cor-clsf-fc}
The proof proceeds as follows. First, we will show using the structure of the data tensor $\tMfc$ that the limit direction of linear coefficients $\vbetafc(\bThetafc^\infty)$ is proportional to $c \vu^\infty$, where $c$ is a nonzero scalar and $\vu^\infty$ is the limit direction of $\mX^T\vr$. Then, through a closer look at $\vu^\infty$ and $c$, we will prove that $\vbetafc(\bThetafc^\infty)$ is in fact a conic combination of the support vectors (i.e., the data points with the minimum margins). Finally, we will compare $\vbetafc(\bThetafc^\infty)$ with the KKT conditions of the $\ell_2$ max-margin classification problem and conclude that $\vbetafc(\bThetafc^\infty)$ must be in the same direction as the $\ell_2$ max-margin classifier.

Due to the way how the data tensor $\tMfc$ is constructed for fully-connected networks (Appendix~\ref{sec:tensorform-ex}), we always have
\begin{equation*}
    -\nabla_{\vv_1}\mc L(\bThetafc) = \tMfc(-\mX^T\vr) \circ (\mI_{k_1}, \vv_2, \dots, \vv_L) \in
    \spann \left \{ 
    \begin{bmatrix}
    \mX^T\vr \\ \zeros \\ \vdots \\ \zeros
    \end{bmatrix},
    \begin{bmatrix}
    \zeros \\ \mX^T\vr \\ \vdots \\ \zeros
    \end{bmatrix},
    \dots,
    \begin{bmatrix}
    \zeros \\ \zeros \\ \vdots \\ \mX^T\vr
    \end{bmatrix}
    \right \}.
\end{equation*}
From Theorem~\ref{thm:meta-thm-clsf-1}, we established directional convergence of $\vv_1$ and its alignment with $-\nabla_{\vv_1}\mc L(\bThetafc)$. This means that the limit direction $\vv_1^\infty$, which is a fixed vector, must be also in the span of vectors written above. This implies that $\mX^T \vr$ must also converge to some direction, say $\vu^\infty \defeq \lim_{t \to \infty} \frac{\mX^T \vr(t)}{\ltwo{\mX^T \vr(t)}}$.

Now recall the definition of $\vv_1$ in case of the fully-connected network: $\vv_1 = \vect(\mW_1)$. So, by reshaping $\vv_1^\infty$ into its original $d_1 \times d_2$ matrix form $\mW_1^\infty$, we have
\begin{equation*}
    \mW_1^\infty \propto \vu^\infty \vq^T,
\end{equation*}
for some $\vq \in \reals^{d_2}$.
This implies that the linear coefficients $\vbetafc(\bThetafc)$ of the network converge in direction to
\begin{equation}
\label{eq:proof-clsf-fc-3}
    \vbetafc(\bThetafc^{\infty}) = \mW_1^\infty \mW_2^\infty \dots \mW_{L-1}^\infty \vw_L^\infty
    \propto \vu^\infty \vq^T \mW_2^\infty \dots \mW_{L-1}^\infty \vw_L^\infty
    = c \vu^\infty,
\end{equation}
where $c$ is some nonzero real number.

Let us now take a closer look at the vector $\vu^\infty$, the limit direction of $\mX^T \vr$.
Recall from Section~\ref{sec:tensorform} that for any $i \in [n]$,
\begin{equation*}
[\vr]_i = -y_i \exp(-y_i \ffc(\vx_i;\bThetafc)) = -y_i \exp(-y_i \vx_i^T \vbetafc(\bThetafc)),
\end{equation*} 
in case of classification. Recall that $\ltwos{\vbetafc(\bThetafc(t))} \to \infty$ while converging to a certain direction $\vbetafc(\bThetafc^\infty)$. This means that if 
\begin{equation*}
    y_j \vx_j^T \vbetafc(\bThetafc^\infty) > y_i\vx_i^T \vbetafc(\bThetafc^\infty)
\end{equation*}
for any $i,j \in [n]$, then
\begin{equation}
\label{eq:proof-clsf-fc-2}
    \lim_{t \to \infty} \frac{\exp(-y_j \vx_j^T \vbetafc(\bThetafc(t)))}{\exp(-y_i \vx_i^T \vbetafc(\bThetafc(t)))} = 0.
\end{equation}
Take $i$ to be the index of any support vector, i.e., any $i$ that attains the minimum $y_i x_i^T \vbetafc(\bThetafc^\infty)$ among all data points. Using such an $i$, the observation~\eqref{eq:proof-clsf-fc-2} implies that $\lim_{t \to\infty}[\vr(t)]_j = 0$ for any $\vx_j$ that is not a support vector.
Thus, by the argument above, $\vu^\infty$ can in fact be written as
\begin{equation}
\label{eq:proof-clsf-fc-4}
    \vu^\infty = \lim_{t \to \infty} \frac{\sum_{i=1}^n \vx_i [\vr(t)]_i}
    {\ltwo{\sum_{i=1}^n \vx_i [\vr(t)]_i}}
    = - \sum_{i=1}^n \nu_i y_i \vx_i,
\end{equation}
where $\nu_i \geq 0$ for all $i \in [n]$, and $\nu_j = 0$ for $\vx_j$'s that are \emph{not} support vectors.
Combining \eqref{eq:proof-clsf-fc-4} and \eqref{eq:proof-clsf-fc-3},
\begin{equation}
\label{eq:proof-clsf-fc-1}
    \vbetafc(\bThetafc^{\infty}) \propto -c \sum_{i = 1}^n \nu_i y_i \vx_i.
\end{equation}

Recall that we do not yet know whether $c$, introduced in~\eqref{eq:proof-clsf-fc-3}, is positive or negative; we will now show that $c$ has to be negative.
From \citet{lyu2020gradient}, we know that $\mc L(\bThetafc(t)) \to 0$, which implies that $y_i \vx_i^T \vbetafc(\bThetafc^\infty) > 0$ for all $i \in [n]$. However, if $c > 0$, then \eqref{eq:proof-clsf-fc-1} implies that $\vbetafc(\bThetafc^{\infty})$ is inside a cone $\mc K$ defined as
\begin{equation*}
    \mc K \defeq \left \{\sum_{i = 1}^n \gamma_i y_i \vx_i \mid \gamma_i \leq 0, \forall i \in [n] \right \}.
\end{equation*}
Note that the polar cone of $\mc K$, denoted as $\mc K^\circ$, is
\begin{equation*}
    \mc K^\circ \defeq \left \{ \vz \mid \vbeta^T \vz \leq 0, \forall \vbeta \in \mc K \right \}
    = \{\vz \mid y_i \vx_i^T \vz \geq 0, \forall i \in [n]\}.
\end{equation*}
It is known that $\mc K \cap \mc K^\circ = \{\zeros\}$ for any convex cone $\mc K$ and its polar cone $\mc K^\circ$. Therefore, having $c>0$ implies that $\vbetafc(\bThetafc^{\infty}) \in \mc K \setminus\mc K^\circ$, which means that there exists some $i \in [n]$ such that $y_i \vx_i^T \vbetafc(\bThetafc^\infty) < 0$; this contradicts the fact that the loss goes to zero as $t \to \infty$. Therefore, $c$ in \eqref{eq:proof-clsf-fc-3} and \eqref{eq:proof-clsf-fc-1} must be negative:
\begin{equation}
\label{eq:proof-clsf-fc-5}
    \vbetafc(\bThetafc^{\infty}) \propto \sum_{i = 1}^n \nu_i y_i \vx_i,
\end{equation}
for $\nu_i \geq 0$ for all $i \in [n]$ and $\nu_j = 0$ for all $\vx_j$'s that are not suport vectors.

Finally, compare \eqref{eq:proof-clsf-fc-5} with the KKT conditions of the following optimization problem:
\begin{align*}
    &\minimize_{\vz}\quad \ltwo{\vz}^2 
    \quad
    \subjectto \quad y_i \vx_i^T \vz \geq 1,~~\forall i\in[n].
    % &\minimize_{\vz}\quad \ltwo{\vz}^2 
    % \quad
    % \subjectto \quad y_i \vx_i^T \vz \leq -1,~~\forall i\in[n].
\end{align*}
The KKT conditions of this problem are
\begin{equation*}
    \vz = \sum_{i=1}^n \mu_i y_i \vx_i,~\text{ and }~\mu_i \geq 0,~ \mu_i(1-y_i\vx_i^T \vz) = 0 \text{ for all } i \in [n],
\end{equation*}
where $\mu_1, \dots, \mu_n$ are the dual variables.
Note that this is (up to scaling) satisfied by $\vbetafc(\bThetafc^\infty)$~\eqref{eq:proof-clsf-fc-5}, if we replace $\mu_i$'s with $\nu_i$'s.
This finishes the proof that $\vbetafc(\bThetafc^{\infty})$ is aligned with the $\ell_2$ max-margin classifier.


\section{Proofs of Theorem~\ref{thm:meta-thm-clsf-2} and Corollaries~\ref{cor:clsf-diag} \& \ref{cor:clsf-conv-full}}
\label{sec:proof-meta-thm-clsf-2}
\subsection{Proof of Theorem~\ref{thm:meta-thm-clsf-2}}
\subsubsection{Convergence of loss to zero}
\label{sec:proof-meta-thm-clsf-2-conv}
We first show that given the conditions on initialization, the training loss $\mc L(\bTheta(t))$ converges to zero.
Recall from Section~\ref{sec:tensorform} that
\begin{align*}
    \dot \vv_l 
    = -\nabla_{\vv_l} \mc L(\bTheta)
    = \tM(-\mX^T \vr) \circ (\vv_1, \dots, \vv_{l-1}, \mI_{k_l}, \vv_{l+1}, \dots, \vv_L).
\end{align*}
Applying the structure~\eqref{eq:tensor-orth-decomp} in Assumption~\ref{assm:complex-diag-decomp}, we get
\begin{align*}
    \dot \vv_l 
    &= \tM(-\mX^T \vr) \circ (\vv_1, \dots, \vv_{l-1}, \mI_{k_l}, \vv_{l+1}, \dots, \vv_L)\\
    &= -\!\sum_{j=1}^m [\mS\mX^T\vr]_j (\vv_1^T[\mU_1]_{\cdot,j} \!\otimes \cdots \otimes \vv_{l-1}^T[\mU_{l-1}]_{\cdot,j} \!\otimes [\mU_l]_{\cdot, j} \!\otimes \vv_{l+1}^T[\mU_{l+1}]_{\cdot, j} \!\otimes \cdots \otimes \vv_{L}^T[\mU_{L}]_{\cdot, j})\\
    &= -\sum_{j=1}^m [\mS\mX^T\vr]_j \bigg(\prod_{k\neq l} [\mU_k^T \vv_k]_{j} \bigg) [\mU_l]_{\cdot, j}.
\end{align*}
Left-multiplying $\mU_l^H$ (the conjugate transpose of $\mU_l$) to both sides, we get
\begin{equation}
\label{eq:proof-clsf-meta2-1}
    \mU_l^H \dot \vv_l = - \mS \mX^T \vr \odot \prod\nolimits_{k \neq l}^{\odot} \mU_k^T \vv_k,
\end{equation}
where $\prod^\odot$ denotes the product using entry-wise multiplication $\odot$.

Now consider the rate of growth for the absolute value squared of the $j$-th component of $\mU_l^T \vv_l$:
\begin{align}
    \frac{d}{dt} |[\mU_l^T \vv_l]_j|^2
    &= \frac{d}{dt} [\mU_l^T \vv_l]_j [\mU_l^T \vv_l]_j^*
    = \frac{d}{dt} [\mU_l^T \vv_l]_j [\mU_l^H \vv_l]_j\notag\\
    &= [\mU_l^T \dot \vv_l]_j [\mU_l^H \vv_l]_j
    + [\mU_l^H \dot \vv_l]_j [\mU_l^T \vv_l]_j\notag\\
    &= 2\Re\left ([\mU_l^H \dot \vv_l]_j [\mU_l^T \vv_l]_j \right )\notag\\
    &= 2\Re\left (-[\mS\mX^T \vr]_j \prod\nolimits_{k=1}^L [\mU_k^T \vv_k]_j \right )\notag\\
    &= \frac{d}{dt} |[\mU_{l'}^T \vv_{l'}]_j|^2\quad\text{ for any $l' \in [L]$, }\notag
\end{align}
so for any $j\in [m]$, the squared absolute value of the $j$-th components in $\mU_l^T \vv_l$ grow at the same rate for each layer $l \in [L]$. This means that the gap between any two different layers stays constant for all $t \geq 0$. Combining this with our conditions on initial directions, we have
\begin{equation}
\label{eq:proof-clsf-meta2-11}
\begin{aligned}
    |[\mU_l^T \vv_l(t)]_j|^2 - |[\mU_L^T \vv_L(t)]_j|^2
    &= |[\mU_l^T \vv_l(0)]_j|^2 - |[\mU_L^T \vv_L(0)]_j|^2\\
    &= \alpha^2 |[\mU_l^T \bar \vv_l]_j|^2 - \alpha^2 |[\mU_L^T \bar \vv_L]_j|^2 \geq \alpha^2 \lambda,
\end{aligned}
\end{equation}
for any $j \in [m]$, $l \in [L-1]$, and $t \geq 0$. This inequality also implies
\begin{equation}
\label{eq:proof-clsf-meta2-2}
    |[\mU_l^T \vv_l(t)]_j|^2 
    \geq 
    |[\mU_L^T \vv_L(t)]_j|^2 + \alpha^2 \lambda
    \geq
    \alpha^2 \lambda.
\end{equation}

Let us now consider the time derivative of $\mc L(\bTheta(t))$.
We have the following chain of upper bounds on the time derivative:
\begin{align}
    \frac{d}{dt} \mc L(\bTheta(t))
    &= \nabla_{\bTheta} \mc L(\bTheta(t))^T \dot \bTheta(t)
    = -\ltwos{\nabla_{\bTheta} \mc L(\bTheta(t))}^2\notag\\
    &\leq -\ltwos{\nabla_{\vv_L} \mc L(\bTheta(t))}^2
    = -\ltwos{\dot \vv_L(t)}^2\notag\\
    &\leqwtxt{(a)} -\ltwos{\mU_L^H \dot \vv_L(t)}^2
    \eqwtxt{(b)} -\ltwo{\mS \mX^T \vr(t) \odot \prod\nolimits_{k \neq L}^{\odot} \mU_k^T \vv_k(t)}^2\notag\\
    &=-\sum\nolimits_{j=1}^m |[\mS \mX^T \vr(t)]_j|^2 \prod\nolimits_{k \neq L} |[\mU_k^T \vv_k(t)]_j|^2\notag\\
    &\leqwtxt{(c)} -\alpha^{2L-2} \lambda^{L-1} \sum\nolimits_{j=1}^m |[\mS \mX^T \vr(t)]_j|^2 \notag\\
    &= -\alpha^{2L-2} \lambda^{L-1} \ltwos{\mS \mX^T \vr(t)}^2\notag\\
    &\leqwtxt{(d)} -\alpha^{2L-2} \lambda^{L-1} s_{\min}(\mS)^2 \ltwos{\mX^T \vr(t)}^2,\label{eq:proof-clsf-meta2-3}
\end{align}
where (a) used the fact that $\ltwos{\dot \vv_L(t)}^2 \geq \ltwos{\mU_L \mU_L^H \dot \vv_L(t)}^2$ because it is a projection onto a subspace, and $\ltwos{\mU_L \mU_L^H \dot \vv_L(t)}^2 = \ltwos{\mU_L^H \dot \vv_L(t)}^2$ because $\mU_L^H \mU_L = \mI_{k_L}$; (b) is due to~\eqref{eq:proof-clsf-meta2-1}; (c) is due to~\eqref{eq:proof-clsf-meta2-2}; and (d) used the fact that $\mS \in \C^{m \times d}$ is a matrix that has full column rank, so for any $\vz \in \C^d$, we can use $\ltwos{\mS \vz} \geq s_{\min}(\mS) \ltwos{\vz}$ where $s_{\min}(\mS)$ is the minimum singular value of $\mS$.

We now prove a lower bound on the quantity $\ltwos{\mX^T \vr(t)}^2$. Recall from Section~\ref{sec:tensorform} the definition of $[\vr(t)]_i = -y_i \exp(-y_i f(\vx_i;\bTheta(t)))$ for classification problems. Also, recall the assumption that the dataset is linearly separable, which means that there exists a unit vector $\vz \in \R^d$ such that
\begin{equation*}
    y_i \vx_i^T \vz \geq \gamma > 0
\end{equation*}
holds for all $i \in [n]$, for some $\gamma > 0$.
Using these,
\begin{align*}
    \ltwos{\mX^T \vr(t)}^2 
    &= \ltwos{\sum\nolimits_{i=1}^n y_i \vx_i \exp(-y_i f(\vx_i;\bTheta(t)))}^2 \\
    &\geq [\vz^T \sum\nolimits_{i=1}^n y_i \vx_i \exp(-y_i f(\vx_i;\bTheta(t)))]^2\\
    &\geq \gamma^2 [\sum\nolimits_{i=1}^n \exp(-y_i f(\vx_i;\bTheta(t)))]^2
    = \gamma^2 \mc L(\bTheta(t))^2.
\end{align*}
Combining this with~\eqref{eq:proof-clsf-meta2-3}, we get
\begin{equation*}
    \frac{d}{dt} \mc L(\bTheta(t)) \leq - \alpha^{2L-2} \lambda^{L-1} s_{\min}(\mS)^2 \gamma^2 \mc L(\bTheta(t))^2,
\end{equation*}
which implies
\begin{equation*}
    \mc L(\bTheta(t)) \leq 
    \frac{\mc L(\bTheta(0))}{1+\alpha^{2L-2} \lambda^{L-1} s_{\min}(\mS)^2 \gamma^2 t}.
\end{equation*}
Therefore, $\mc L(\bTheta(t)) \to 0$ as $t \to \infty$.

\subsubsection{Characterizing the limit direction}
Since we have $\mc L(\bTheta(t)) \to 0$, the argument in the proof of Theorem~\ref{thm:meta-thm-clsf-1} applies to this case, and it shows that the parameters $\vv_l$ converge in direction and align with $\dot \vv_l = -\nabla_{\vv_l} \mc L(\bTheta)$.
Let $\vv_l^\infty \defeq \lim_{t \to \infty} \frac{\vv_l(t)}{\ltwos{\vv_l(t)}}$ be the limit direction of $\vv_l$.
Recall also that we additionally assumed that $\mX^T \vr(t)$ converges in direction. Let $\vu^\infty \defeq \lim_{t \to \infty} \frac{\mS \mX^T \vr(t)}{\ltwos{\mS \mX^T \vr(t)}}$, which exists due to the directional convergence of $\mX^T \vr(t)$.

For the remaining steps of the proof, we derive a number of conditions that has to be satisfied by the limit directions of the parameters. Next, we compare these conditions with the KKT conditions of the minimization problem, and finish the proof.

By Assumption~\ref{assm:complex-diag-decomp}, we have
\begin{align*}
    f(\vx; \bTheta) 
    &= \tM(\vx) \circ (\vv_1, \dots, \vv_L)
    = \sum_{j=1}^m [\mS \vx]_j \prod_{l=1}^L [\mU_l^T \vv_l]_j\\
    &= \bigg[ \sum_{j=1}^m \bigg(\prod_{l=1}^L [\mU_l^T \vv_l]_j \bigg) [\mS]_{j,\cdot} \bigg ]\vx
    = \vx^T \mS^T \bigg(\prod\nolimits_{l \in [L]}^{\odot} \mU_l^T \vv_l \bigg)
    = \vx^T \mS^T \vrho.
\end{align*}
Here, we defined $\vrho \defeq \prod\nolimits_{l \in [L]}^{\odot} \mU_l^T \vv_l \in \C^m$.
Since the linear coefficients must be real, we have $\mS^T \vrho \in \R^d$ for any real $\vv_l$'s.
Since $\vv_l$'s converge in direction, $\vrho$ also converges in direction, to $\vrho^\infty \defeq \prod\nolimits_{l \in [L]}^{\odot} \mU_l^T \vv_l^\infty$.
So we can express the limit direction of $\vbeta(\bTheta)$ as
\begin{equation}
\label{eq:proof-clsf-meta2-19}
    \vbeta(\bTheta^\infty) \propto \mS^T \bigg(\prod\nolimits_{l \in [L]}^{\odot} \mU_l^T \vv_l^\infty \bigg) = \mS^T \vrho^\infty.
\end{equation}

Below, we would like to show the following three conditions hold for $\vu^\infty$ and $\vrho^\infty$.
\begin{enumerate}[label=\textup{(\alph*)}]
	\item \label{cond:IBproof1} $|[\vrho^\infty]_j| \neq 0 \implies \Arg(-[\vu^\infty]_j) + \Arg([\vrho^\infty]_j) = 0$,
	\item \label{cond:IBproof2} $|[\vrho^\infty]_j| \neq 0 \implies |[\vu^\infty]_j| \propto |[\vrho^\infty]_j|^{\frac{2}{L}-1}$,
	\item \label{cond:IBproof3} If $L = 2$, then $|[\vrho^\infty]_j| = 0, |[\vrho^\infty]_{j'}| \neq 0 \implies
	|[\vu^\infty]_j| \leq |[\vu^\infty]_{j'}|$,
\end{enumerate}
for any $j, j' \in [m]$.

To prove the first two Conditions~\ref{cond:IBproof1} and \ref{cond:IBproof2}, assume that all the components in $\vrho^\infty$ are nonzero, i.e., $|[\vrho^\infty]_j| \neq 0$. This is without loss of generality because the conditions do not require anything about the case $|[\vrho^\infty]_j| = 0$.
Having $|[\vrho^\infty]_j| \neq 0$ for all $j \in [m]$ implies that $\mU_{l}^T \vv_l^\infty$'s also do not have any zero components. 

From~\eqref{eq:proof-clsf-meta2-1} and alignment of $\vv_l$ and $\dot \vv_l$, we have
\begin{equation}
\label{eq:proof-clsf-meta2-17}
    \lim_{t \to \infty} \mU_{l}^H \vv_l(t)
    = 
    \lim_{t \to \infty} (\mU_{l}^T \vv_l(t))^*
    %= \mU_{l}^H \lim_{t \to \infty} \frac{\vv_l(t)}{\ltwos{\vv_l(t)}}
    \propto -\lim_{t \to \infty}
    \mS\mX^T \vr(t) \odot \prod\nolimits_{k \neq l}^{\odot} \mU_k^T \vv_k(t).
\end{equation}

Using $\vu^\infty$, we can rewrite \eqref{eq:proof-clsf-meta2-17} as 
\begin{equation*}
    \mU_{l}^H \vv_l^\infty \propto - \vu^\infty \odot \prod\nolimits_{k \neq l}^{\odot} \mU_k^T \vv_k^\infty,
\end{equation*}
for all $l \in [L]$. Element-wise multiplying $\mU_{l}^T \vv_l^\infty$ to both sides gives
\begin{equation}
\label{eq:proof-clsf-meta2-4}
    \mU_{l}^T \vv_l^\infty \odot \mU_{l}^H \vv_l^\infty = |\mU_{l}^T \vv_l^\infty|^{\odot 2} \propto - \vu^\infty \odot \prod\nolimits_{k \in [L]}^{\odot} \mU_k^T \vv_k^\infty 
    = -\vu^\infty \odot \vrho^\infty,
\end{equation}
where $\va^{\odot b}$ denotes element-wise $b$-th power of the vector $\va$.
% Define $\vrho^\infty \in \C^m$ as the element-wise product of $\mU_k^T \vv_k^\infty$'s
% \begin{equation*}
%     \vrho^\infty \defeq \prod\nolimits_{k \in [L]}^{\odot} \mU_k^T \vv_k^\infty.
% \end{equation*}
Since the LHS of~\eqref{eq:proof-clsf-meta2-4} is a positive real number, we have
\begin{equation*}
	\Arg(|[\mU_l^T \vv_l^\infty]_j|^2) = 0 = \Arg([-\vu^\infty]_j) + \Arg([\vrho^\infty]_j),
\end{equation*}
hence proving Condition~\ref{cond:IBproof1}.
Using this, \eqref{eq:proof-clsf-meta2-4} becomes
\begin{equation}
\label{eq:proof-clsf-meta2-5}
    |\mU_{l}^T \vv_l^\infty|^{\odot 2} \propto  |\vu^\infty| \odot |\vrho^\infty|.
\end{equation}
Now element-wise multiply \eqref{eq:proof-clsf-meta2-5} for all $l \in [L]$, then we get
\begin{equation}
\label{eq:proof-clsf-meta2-6}
    |\vrho^\infty|^{\odot 2} \propto |\vu^\infty|^{\odot L} \odot |\vrho^\infty|^{\odot L},
\end{equation}
from which we conclude
\begin{align*}
%\label{eq:proof-clsf-meta2-7}
|[\vrho^\infty]_j| \neq 0 \implies |[\vu^\infty]_j| \propto |[\vrho^\infty]_j|^{\frac{2}{L}-1},
\end{align*}
for all $j \in [m]$. This proves Condition~\ref{cond:IBproof2}.

Now, we are left with Condition~\ref{cond:IBproof3}, which only concerns the case $L = 2$.
First, consider the time derivative of $[\vrho]_j = [\mU_1^T \vv_1]_j [\mU_2^T \vv_2]_j$.
\begin{align}
    \frac{d}{dt} [\vrho(t)]_j
    &=[\mU_1^T \vv_1(t)]_j \frac{d}{dt} [\mU_2^T \vv_2(t)]_j
    + [\mU_2^T \vv_2(t)]_j \frac{d}{dt} [\mU_1^T \vv_1(t)]_j\notag\\
    &\eqwtxt{(a)}
    -[\mS\mX^T \vr(t)]_j^* (|[\mU_1^T \vv_1(t)]_j|^2 + |[\mU_2^T \vv_2(t)]_j|^2),\label{eq:proof-clsf-meta2-9}
\end{align}
where (a) used \eqref{eq:proof-clsf-meta2-1}.
Since $|[\mU_1^T \vv_1(t)]_j|^2 \geq \alpha^2 \lambda$ \eqref{eq:proof-clsf-meta2-2} by our assumption on initialization, \eqref{eq:proof-clsf-meta2-9} implies that whenever $[\vu^\infty]_j \neq 0$, the derivative $\frac{d}{dt} [\vrho(t)]_j$ is nonzero and has phase equal to $\Arg(-[\vu^\infty]_j^*)$ in the limit $t \to \infty$. This also implies that $[\vrho(t)]_j$ does not stay stuck at zero forever, provided that $[\vu^\infty]_j \neq 0$. 

Now consider
\begin{align}
\label{eq:proof-clsf-meta2-10}
    \frac{\left | \frac{d}{dt} [\vrho(t)]_j \right |}{\ltwos{\mS \mX^T \vr(t)}|[\vrho(t)]_j|} 
    = \frac{|[\mS\mX^T \vr(t)]_j|}{\ltwos{\mS \mX^T \vr(t)}}\frac{|[\mU_1^T \vv_1(t)]_j|^2 + |[\mU_2^T \vv_2(t)]_j|^2}{|[\vrho(t)]_j|}.
\end{align}
We want to compare this quantity for different $j,j' \in [m]$ satisfying $[\vu^\infty]_j \neq 0$ and $[\vu^\infty]_{j'} \neq 0$. Before we do that, we take a look at the last term in the RHS of~\eqref{eq:proof-clsf-meta2-10}. Recall from \eqref{eq:proof-clsf-meta2-11} that 
\begin{equation}
\label{eq:proof-clsf-meta2-12}
|[\mU_1^T \vv_1(t)]_j|^2 = |[\mU_2^T \vv_2(t)]_j|^2+|[\mU_1^T \vv_1(0)]_j|^2- |[\mU_2^T \vv_2(0)]_j|^2.
\end{equation}
For simplicity, let $\delta_j \defeq |[\mU_1^T \vv_1(0)]_j|^2- |[\mU_2^T \vv_2(0)]_j|^2$, which is a fixed positive number due to our assumption on initialization. Then, we can use \eqref{eq:proof-clsf-meta2-12} and $|[\vrho(t)]_j|=|[\mU_1^T \vv_1(t)]_j||[\mU_2^T \vv_2(t)]_j|$ to show that
\begin{align}
	&\frac{|[\mU_1^T \vv_1(t)]_j|^2 + |[\mU_2^T \vv_2(t)]_j|^2}{|[\vrho(t)]_j|}
	=
	\frac{2 |[\mU_2^T \vv_2(t)]_j|^2 + \delta_j}
	{|[\mU_2^T \vv_2(t)]_j| \sqrt{|[\mU_2^T \vv_2(t)]_j|^2+\delta_j}}
	\geq 2,\label{eq:fix1}\\
	&\lim_{t \to \infty}
	\frac{|[\mU_1^T \vv_1(t)]_j|^2 + |[\mU_2^T \vv_2(t)]_j|^2}{|[\vrho(t)]_j|}
	= 2 \text{~~~if } \lim_{t \to \infty} |[\mU_2^T \vv_2(t)]_j| = \infty.\label{eq:fix2}
\end{align}
Recall that we want to prove Condition~\ref{cond:IBproof3}, namely
\begin{align*}
	|[\vrho^\infty]_j| = 0, |[\vrho^\infty]_{j'}| \neq 0 \implies
	|[\vu^\infty]_j| \leq |[\vu^\infty]_{j'}|.
\end{align*}
For the sake of contradiction, suppose that there exists $j \in [m]$ that satisfies $|[\vrho^\infty]_j| = 0$ but $|[\vu^\infty]_j| > |[\vu^\infty]_{j'}|$, for some $j' \in [m]$ satisfying $|[\vrho^\infty]_{j'}| \neq 0$.
Note from Condition~\ref{cond:IBproof2} that $|[\vu^\infty]_{j'}| > 0$.
Having $|[\vrho^\infty]_j| = 0$ and $|[\vrho^\infty]_{j'}| \neq 0$ implies that $|[\vrho(t)]_{j'}| \to \infty$ and $\frac{|[\vrho(t)]_j|}{|[\vrho(t)]_{j'}|} \to 0$.
We now want to compute the ratio of \eqref{eq:proof-clsf-meta2-10} for $j$ and $j'$. First, note that
\begin{align}
	\label{eq:proof-clsf-meta2-13}
	&\lim_{t \to \infty} 
	\frac{|[\mS\mX^T \vr(t)]_j|/\ltwos{\mS \mX^T \vr(t)}}
	{|[\mS\mX^T \vr(t)]_{j'}|/\ltwos{\mS \mX^T \vr(t)}}
	= \frac{|[\vu^\infty]_j|}{|[\vu^\infty]_{j'}|} > 1.
\end{align}
Next, using $|[\vrho(t)]_{j'}| \to \infty$, \eqref{eq:fix1}, and \eqref{eq:fix2}, we have
\begin{align}
	\label{eq:proof-clsf-meta2-14}
	\lim_{t\to\infty} \frac{|[\mU_1^T \vv_1(t)]_{j}|^2 + |[\mU_2^T \vv_2(t)]_{j}|^2}{|[\vrho(t)]_{j}|}
	\geq
	\lim_{t\to\infty} \frac{|[\mU_1^T \vv_1(t)]_{j'}|^2 + |[\mU_2^T \vv_2(t)]_{j'}|^2}{|[\vrho(t)]_{j'}|} = 2.
\end{align}
Combining \eqref{eq:proof-clsf-meta2-13} and \eqref{eq:proof-clsf-meta2-14} to compute the ratio of \eqref{eq:proof-clsf-meta2-10} for $j$ and $j'$, we get that there exists some $t_0' \geq 0$ such that for any $t \geq t_0'$, we have
\begin{align}
	\label{eq:proof-clsf-meta2-15}
	\frac{\left | \frac{d}{dt} [\vrho(t)]_j \right |/|[\vrho(t)]_j|}
	{\left | \frac{d}{dt} [\vrho(t)]_{j'} \right |/|[\vrho(t)]_{j'}|}
	> 1.
\end{align}
This implies that the ratio of the absolute value of time derivative of $[\vrho(t)]_j$ to the absolute value of current value of $[\vrho(t)]_j$ is strictly bigger than that of $[\vrho(t)]_{j'}$. 
Moreover, we saw in \eqref{eq:proof-clsf-meta2-9} that the phase of $\frac{d}{dt} [\vrho(t)]_j$ converges to that of $-[\vu^\infty]_j^*$. Since this holds for all $t \geq t_0'$, \eqref{eq:proof-clsf-meta2-15} results in a growth of $|[\vrho(t)]_j|$ that is exponentially faster than that of $|[\vrho(t)]_{j'}|$, so $[\vrho(t)]_j$ becomes a dominant component in $\vrho(t)$ as $t \to \infty$. This contradicts that $[\vrho^\infty]_j = 0$, hence Condition~\ref{cond:IBproof3} has to be satisfied.

In addition to the three conditioned proven above, one can use the same argument as in Appendix~\ref{sec:proof-cor-clsf-fc}, more specifically \eqref{eq:proof-clsf-fc-2} and \eqref{eq:proof-clsf-fc-4}, to show that $\vu^\infty$ can be written as
\begin{equation}
	\label{eq:proof-clsf-meta2-16}
	\vu^\infty = \lim_{t \to \infty} 
	\frac{\mS \sum_{i=1}^n \vx_i [\vr(t)]_i}
	{\ltwo{\mS \sum_{i=1}^n \vx_i [\vr(t)]_i}}
	= - \mS \sum_{i=1}^n \nu_i y_i \vx_i,
\end{equation}
where $\nu_i \geq 0$ for all $i \in [n]$, and $\nu_j = 0$ for $\vx_j$'s that are \emph{not} support vectors, i.e., those satisfying $y_j \vx_j^T \mS^T \vrho^\infty > \min_{i \in [n]} y_i \vx_i^T \mS^T \vrho^\infty$.

So far, we have shown that Conditions~\ref{cond:IBproof1}, \ref{cond:IBproof2}, and \ref{cond:IBproof3} as well as \eqref{eq:proof-clsf-meta2-16} are satisfied by the limit directions $\vu^\infty$ and $\vrho^\infty$ of $\mS \mX^T \vr$ and $\vrho$. We now consider the following optimization problem and prove that these conditions are in fact the KKT conditions of the optimization problem.
Consider 
\begin{equation}
	\label{eq:proof-clsf-meta2-20}
	\minimize_{\vrho \in \C^m}\quad \norm{\vrho}_{2/L}
	\quad
	\subjectto \quad y_i \vx_i^T\mS^T \vrho \geq 1,~~\forall i\in[n].
\end{equation}
The KKT conditions of this problem are
\begin{equation*}
	\partial \norm{\vrho}_{2/L} \ni \mS^* \sum_{i=1}^n \mu_i y_i \vx_i,~\text{ and }~\mu_i \geq 0,~ \mu_i(1-y_i\vx_i^T \mS^T \vrho) = 0 \text{ for all } i \in [n],
\end{equation*}
where $\mu_1, \dots, \mu_n$ are the dual variables.
The symbol $\partial \norm{\cdot}_{2/L}$ denotes the (local) subdifferential of the $\ell_{2/L}$ norm\footnote{We use the definition of subdifferentials from \citet{gunasekar2018implicit}.}, which can be written as
\begin{align*}
	\partial \norm{\vrho}_{1}
	=
	\{ \vu \in \C^m \mid &|[\vu]_j| \leq 1 \text{ for all } j \in [m], \text{ and }\\ &[\vrho]_j \neq 0 \implies [\vu]_j = \exp(\sqrt{-1} \Arg([\vrho]_j))\},
\end{align*}
if $L = 2$ (in this case $\partial \norm{\vrho}_{1}$ is the global subdifferential), and
\begin{equation*}
	\partial \norm{\vrho}_{2/L}
	=
	\left \{ \vu \in \C^m \mid [\vrho]_j \neq 0 \implies [\vu]_j = \frac{2}{L} |[\vrho]_j|^{\frac{2}{L}-1} \exp(\sqrt{-1} \Arg([\vrho]_j)) \right \}, 
\end{equation*}
if $L > 2$.
By replacing $\mu_i$'s with $\nu_i$'s defined in \eqref{eq:proof-clsf-meta2-16}, we can check from \eqref{eq:proof-clsf-meta2-16}, Conditions~\ref{cond:IBproof1}, \ref{cond:IBproof2}, and \ref{cond:IBproof3} that the that $\vrho^\infty$ and $\vu^\infty$ satisfy the KKT conditions up to scaling. 
Therefore, by \eqref{eq:proof-clsf-meta2-19}, $\vbeta(\bTheta(t))$ converges in direction aligned with $\mS^T \vrho^{\infty}$, where $\vrho^\infty$ is aligned with a stationary point (global minimum in case of $L=2$) of the optimization problem~\eqref{eq:proof-clsf-meta2-20}.

If $\mS$ is invertible, we can get $\mS^{-T} \vbeta(\bTheta^\infty) \propto \vrho^\infty$. Plugging this into the optimization problem~\eqref{eq:proof-clsf-meta2-20} gives the last statement of the theorem.

\subsection{Proof of Corollary~\ref{cor:clsf-diag}}
\label{sec:proof-cor-clsf-diag}
It suffices to prove that linear diagonal networks satisfy Assumption~\ref{assm:complex-diag-decomp}, with $\mS = \mI_d$.
The proof is very straightforward, since $\tMdiag(\vx) \in \R^{d \times \cdots \times d}$ has $[\tMdiag(\vx)]_{j,j,\dots,j} = [\vx]_j$ while all the remaining entries are zero. It is straightforward to verify that $\tMdiag(\vx)$ satisfies Assumption~\ref{assm:complex-diag-decomp} with $\mS = \mU_1 = \dots = \mU_L = \mI_d$.
A direct substitution into Theorem~\ref{thm:meta-thm-clsf-2} gives the corollary. 


\subsection{Proof of Corollary~\ref{cor:clsf-conv-full}}
\label{sec:proof-cor-clsf-conv-full}
For full-length convolutional networks ($k_1 = \dots = k_L = d$), we will prove that they satisfy Assumption~\ref{assm:complex-diag-decomp} with
$\mS = d^{\frac{L-1}{2}} \mF$ and $\mU_1 = \dots = \mU_L = \mF^*$, where $\mF \in \C^{d \times d}$ is the matrix of discrete Fourier transform basis $[\mF]_{j,k} = \frac{1}{\sqrt{d}} \exp(-\frac{\sqrt{-1} \cdot 2 \pi (j-1)(k-1)}{d})$ and $\mF^*$ is the complex conjugate of $\mF$.

For simplicity of notation, define $\psi=\exp(-\frac{\sqrt{-1}\cdot2\pi}{d})$.
With matrices $\mS$ and $\mU_1, \dots, \mU_L$ chosen as above, we can write $\tM(\vx)$ as
\begin{align*}
    \tM(\vx) &= \sum_{j=1}^d [\mS \vx]_j ([\mU_1]_{\cdot,j} \otimes [\mU_2]_{\cdot,j} \otimes \cdots \otimes [\mU_L]_{\cdot,j})\\
    &= \sum_{j=1}^d \left [d^{\frac{L-2}{2}} \sum_{k=1}^d [\vx]_k \psi^{(j-1)(k-1)} \right] 
    \begin{bmatrix}
    \psi^0 / \sqrt{d} \\ \psi^{-(j-1)} / \sqrt{d} \\ \psi^{-2(j-1)} / \sqrt{d} \\ \vdots \\ \psi^{-(d-1)(j-1)} / \sqrt{d}
    \end{bmatrix}^{\otimes L},
\end{align*}
where $\va^{\otimes L}$ denotes the $L$-times tensor product of $\va$.
We will show that $\tM(\vx) = \tMconv(\vx)$.

For any $j_1, \dots, j_L \in [d]$,
\begin{align*}
    [\tM(\vx)]_{j_1, \dots, j_L} &= 
    \frac{1}{d} \sum_{l=1}^d \left [\sum_{k=1}^d [\vx]_k \psi^{(l-1)(k-1)} \right] \psi^{-(l-1)(\sum_{q=1}^L j_q - L)}\\
    &= \frac{1}{d} \sum_{k=1}^d [\vx]_k 
    \sum_{l=1}^d \psi^{(l-1)(k-1-\sum_{q=1}^L j_q + L)}.
\end{align*}
Recall that 
\begin{align*}
    \sum_{l=1}^d \psi^{(l-1)(k-1-\sum_{q=1}^L j_q + L)}
    = \begin{cases}
    d & \text{ if } k-1-\sum_{q=1}^L j_q + L \text{ is a multiple of $d$},\\
    0 & \text{ otherwise}.
    \end{cases}
\end{align*}
Using this, we have
\begin{align*}
    [\tM(\vx)]_{j_1, \dots, j_L} 
    &= \frac{1}{d} \sum_{k=1}^d [\vx]_k 
    \sum_{l=1}^d \psi^{(l-1)(k-1-\sum_{q=1}^L j_q + L)}\\
    &= [\vx]_{\sum_{q=1}^L j_q -L+1 \bmod d}
    = [\tMconv(\vx)]_{j_1, \dots, j_L}.
\end{align*}
Hence, linear full-length convolutional networks satisfy Assumption~\ref{assm:complex-diag-decomp} with $\mS = d^{\frac{L-1}{2}} \mF$. A direct substitution into Theorem~\ref{thm:meta-thm-clsf-2} and then using the fact that $|[\mF \vz]_j| = |[\mF^* \vz]_j|$ for any real vector $\vz \in \R^d$ gives the corollary.



\section{Proofs of Theorem~\ref{thm:meta-thm-clsf-3} and Corollary~\ref{cor:clsf-conv-small}}
\label{sec:proof-meta-thm-clsf-3}
\subsection{Proof of Theorem~\ref{thm:meta-thm-clsf-3}}
\subsubsection{Convergence of loss to zero}
We first show that given the conditions on initialization, the training loss $\mc L(\bTheta(t))$ converges to zero. 
Since $L=2$ and $\tM(\vx) = \mU_1\diag(\vs)\mU_2^T$, we can write the gradient flow dynamics from Section~\ref{sec:tensorform} as
\begin{equation}
\label{eq:proof-clsf-meta3-1}
\begin{aligned}
    \dot \vv_1
    &= -\tM(\mX^T \vr) \circ (\mI_{k_1}, \vv_2)
    = -r \mU_1 \diag(\vs) \mU_2^T \vv_2,\\
    \dot \vv_2
    &= -\tM(\mX^T \vr) \circ (\vv_1, \mI_{k_2})
    = -r \mU_2 \diag(\vs) \mU_1^T \vv_1,
\end{aligned}
\end{equation}
where $r(t) = -y \exp(-y f(\vx;\bTheta(t)))$ is the residual of the data point $(\vx,y)$.
From \eqref{eq:proof-clsf-meta3-1} we get
\begin{equation}
\label{eq:proof-clsf-meta3-2}
    \mU_1^T \dot \vv_1 = -r \vs \odot \mU_2^T \vv_2,~~
    \mU_2^T \dot \vv_2 = -r \vs \odot \mU_1^T \vv_1.
\end{equation}

Now consider the rate of growth for the $j$-th component of $\mU_1^T \vv_1$ squared:
\begin{equation}
\label{eq:proof-clsf-meta3-15}
    \frac{d}{dt} [\mU_1^T \vv_1]_j^2
    = 2 [\mU_1^T \vv_1]_j [\mU_1^T \dot \vv_1]_j
    = -2 r [\vs]_j [\mU_1^T \vv_1]_j [\mU_2^T \vv_2]_j
    = \frac{d}{dt} [\mU_{2}^T \vv_{2}]_j^2.
\end{equation}
So for any $j\in [m]$, $[\mU_1^T \vv_1]_j^2$ and $[\mU_2^T \vv_2]_j^2$ grow at the same rate. This means that the gap between the two layers stays constant for all $t \geq 0$. Combining this with our conditions on initial directions,
\begin{equation}
\label{eq:proof-clsf-meta3-9}
\begin{aligned}
    [\mU_1^T \vv_1(t)]_j^2 - [\mU_2^T \vv_2(t)]_j^2
    &= [\mU_1^T \vv_1(0)]_j^2 - [\mU_2^T \vv_2(0)]_j^2\\
    &= \alpha^2 [\mU_1^T \bar \vv_1]_j^2 - \alpha^2 [\mU_2^T \bar \vv_2]_j^2 \geq \alpha^2 \lambda,
\end{aligned}
\end{equation}
for any $j \in [m]$ and $t \geq 0$. This inequality implies
\begin{equation}
\label{eq:proof-clsf-meta3-3}
    [\mU_1^T \vv_1(t)]_j^2 
    \geq 
    [\mU_2^T \vv_2(t)]_j^2 + \alpha^2 \lambda
    \geq
    \alpha^2 \lambda.
\end{equation}

Let us now consider the time derivative of $\mc L(\bTheta(t))$.
We have the following chain of upper bounds on the time derivative:
\begin{align}
    \frac{d}{dt} \mc L(\bTheta(t))
    &= \nabla_{\bTheta} \mc L(\bTheta(t))^T \dot \bTheta(t)
    = -\ltwos{\nabla_{\bTheta} \mc L(\bTheta(t))}^2\notag\\
    &\leq -\ltwos{\nabla_{\vv_2} \mc L(\bTheta(t))}^2
    = -\ltwos{\dot \vv_2(t)}^2\notag\\
    &\leqwtxt{(a)} -\ltwos{\mU_2^T \dot \vv_2(t)}^2
    \eqwtxt{(b)} -r(t)^2 \ltwo{\vs \odot \mU_1^T \vv_1(t)}^2\notag\\
    &=-r(t)^2 \sum\nolimits_{j=1}^m [\vs]_j^2 [\mU_1^T \vv_1(t)]_j^2\notag\\
    &\leqwtxt{(c)} -\alpha^{2} \lambda r(t)^2 \sum\nolimits_{j=1}^m [\vs]_j^2 \notag\\
    &= -\alpha^{2} \lambda \ltwos{\vs}^2 \mc L(\bTheta(t))^2\notag,
\end{align}
where (a) used the fact that $\ltwos{\dot \vv_2(t)}^2 \geq \ltwos{\mU_2 \mU_2^T \dot \vv_2(t)}^2$ because it is a projection onto a subspace, and $\ltwos{\mU_2 \mU_2^T \dot \vv_L(t)}^2 = \ltwos{\mU_2^T \dot \vv_2(t)}^2$ because $\mU_2^T \mU_2 = \mI_{k_2}$; (b) is due to~\eqref{eq:proof-clsf-meta3-2}; (c) is due to~\eqref{eq:proof-clsf-meta3-3}.
From this, we get
\begin{equation*}
    \mc L(\bTheta(t)) \leq 
    \frac{\mc L(\bTheta(0))}{1+\alpha^{2} \lambda \ltwos{\vs}^2 t}.
\end{equation*}
Therefore, $\mc L(\bTheta(t)) \to 0$ as $t \to \infty$.

\subsubsection{Characterizing the limit direction}
Since we proved that $\mc L(\bTheta(t)) \to 0$, the argument in the proof of Theorem~\ref{thm:meta-thm-clsf-1} applies to this case, and shows that the parameters $\vv_l$ converge in direction and align with $\dot \vv_l = -\nabla_{\vv_l} \mc L(\bTheta)$.
Let $\vv_l^\infty \defeq \lim_{t \to \infty} \frac{\vv_l(t)}{\ltwos{\vv_l(t)}}$ be the limit direction of $\vv_l$.
As done in the proof of Theorem~\ref{thm:meta-thm-clsf-2}, define $\vrho(t) = \mU_1^T \vv_1(t) \odot \mU_2^T \vv_2(t)$ and $\vrho^\infty = \mU_1^T \vv_1^\infty \odot \mU_2^T \vv_2^\infty$.
% The remaining steps of the proof are as follows. We first prove that $\mS \mX^T \vr(t)$ converges in direction $\vu^\infty$. Using this $\vu^\infty$, we derive a number of conditions that has to be satisfied by the limit directions of the parameters. Finally, we compare these conditions with the KKT conditions of the minimization problem, and finish the proof.

It follows from $r(t) = -y \exp(-y f(\vx;\bTheta(t)))$ that we have $\sign(r(t)) = -\sign(y)$. Using this, \eqref{eq:proof-clsf-meta3-2}, and alignment of $\vv_l$ and $\dot \vv_l$, we have
\begin{equation}
\label{eq:proof-clsf-meta3-14}
    \mU_{1}^T \vv_1^\infty \propto y \vs \odot \mU_2^T \vv_2^\infty,~~
    \mU_{2}^T \vv_2^\infty \propto y \vs \odot \mU_1^T \vv_1^\infty.
\end{equation}

Element-wise multiplying LHSs to both sides gives
\begin{equation}
\label{eq:proof-clsf-meta3-4}
    (\mU_{1}^T \vv_1^\infty)^{\odot 2} \propto y\vs \odot \vrho^\infty,~~
    (\mU_{2}^T \vv_2^\infty)^{\odot 2} \propto y\vs \odot \vrho^\infty.
\end{equation}
Since the LHSs are positive and $\vs$ is positive, the following equations have to be satisfied for all $j \in [m]$:
\begin{equation}
\label{eq:proof-clsf-meta3-16}
    \sign(y) = \sign([\vrho^\infty]_j).
\end{equation}
Now, multiplying both sides of the two equations~\eqref{eq:proof-clsf-meta3-4}, we get
\begin{equation}
\label{eq:proof-clsf-meta3-5}
    (\vrho^{\infty})^{\odot 2} \propto \vs^{\odot 2} \odot (\vrho^\infty)^{\odot 2}.
\end{equation}
From \eqref{eq:proof-clsf-meta3-5}, $\vrho^\infty$ must satisfy that
\begin{align}
\label{eq:proof-clsf-meta3-6}
[\vrho^\infty]_j \neq 0, [\vrho^\infty]_{j'} \neq 0 \implies [\vs]_j = [\vs]_{j'} > 0,
\end{align}
for all $j, j' \in [m]$.
As in the proof of Theorem~\ref{thm:meta-thm-clsf-2}, there is another condition that has to be satisfied:
\begin{align}
\label{eq:proof-clsf-meta3-7}
[\vrho^\infty]_j = 0, [\vrho^\infty]_{j'} \neq 0 &\implies
[\vs]_j \leq [\vs]_{j'},
\end{align}
for any $j, j'\in [m]$; let us prove why. 
First, consider the time derivative of $[\vrho]_j = [\mU_1^T \vv_1]_j [\mU_2^T \vv_2]_j$.
\begin{align}
	\frac{d}{dt} [\vrho(t)]_j
	&=[\mU_1^T \vv_1(t)]_j \frac{d}{dt} [\mU_2^T \vv_2(t)]_j
	+ [\mU_2^T \vv_2(t)]_j \frac{d}{dt} [\mU_1^T \vv_1(t)]_j\notag\\
	&\eqwtxt{(a)}
	-r(t) [\vs]_j ([\mU_1^T \vv_1(t)]_j^2 + [\mU_2^T \vv_2(t)]_j^2),\label{eq:proof-clsf-meta3-17}
\end{align}
where (a) used \eqref{eq:proof-clsf-meta3-2}.
Since $|[\mU_1^T \vv_1(t)]_j|^2 \geq \alpha^2 \lambda$ \eqref{eq:proof-clsf-meta3-3} by our assumption on initialization, \eqref{eq:proof-clsf-meta3-17} implies that whenever $[\vs]_j \neq 0$, the derivative $\frac{d}{dt} [\vrho(t)]_j$ is nonzero and has sign equal to $y$. This also implies that $[\vrho(t)]_j$ does not stay stuck at zero forever, provided that $[\vs]_j \neq 0$. 

Now consider
\begin{align}
\label{eq:proof-clsf-meta3-8}
    \frac{\left | \frac{d}{dt} [\vrho(t)]_j \right |}{|r(t)||[\vrho(t)]_j|} 
    = [\vs]_j \frac{[\mU_1^T \vv_1(t)]_j^2 + [\mU_2^T \vv_2(t)]_j^2}{|[\vrho(t)]_j|}.
\end{align}
We want to compare this quantity for different $j,j' \in [m]$ satisfying $[\vs]_j \neq 0$ and $[\vs]_{j'} \neq 0$. Before we do that, we take a look at the last term in the RHS of~\eqref{eq:proof-clsf-meta3-8}. Recall from \eqref{eq:proof-clsf-meta3-9} that 
\begin{equation}
\label{eq:proof-clsf-meta3-10}
[\mU_1^T \vv_1(t)]_j^2 = [\mU_2^T \vv_2(t)]_j^2+[\mU_1^T \vv_1(0)]_j^2- [\mU_2^T \vv_2(0)]_j^2.
\end{equation}
For simplicity, let $\delta_j \defeq [\mU_1^T \vv_1(0)]_j^2- [\mU_2^T \vv_2(0)]_j^2$, which is a fixed positive number due to our assumption on initialization. Then, we use $|[\vrho(t)]_j|=|[\mU_1^T \vv_1(t)]_j||[\mU_2^T \vv_2(t)]_j|$ and \eqref{eq:proof-clsf-meta3-10} to show that
\begin{align}
	&\frac{[\mU_1^T \vv_1(t)]_j^2 + [\mU_2^T \vv_2(t)]_j^2}{|[\vrho(t)]_j|}
	=
	\frac{2 [\mU_2^T \vv_2(t)]_j^2 + \delta_j}
	{|[\mU_2^T \vv_2(t)]_j| \sqrt{[\mU_2^T \vv_2(t)]_j^2+\delta_j}}
	\geq 2,\label{eq:fix3}\\
	&\lim_{t \to \infty}
	\frac{[\mU_1^T \vv_1(t)]_j^2 + [\mU_2^T \vv_2(t)]_j^2}{|[\vrho(t)]_j|}
	= 2 \text{~~~if } \lim_{t \to \infty} |[\mU_2^T \vv_2(t)]_j| = \infty.\label{eq:fix4}
\end{align}
Recall that we want to prove that \eqref{eq:proof-clsf-meta3-7} should necessarily hold. For the sake of contradiction, suppose that there exists $j \in [m]$ that satisfies $[\vrho^\infty]_j = 0$ but $[\vs]_j > [\vs]_{j'}$, for some $j' \in [m]$ satisfying $[\vrho^\infty]_{j'} \neq 0$. 
Note from \eqref{eq:proof-clsf-meta3-6} that $[\vs]_{j'} > 0$.
Having $[\vrho^\infty]_j = 0$ and $[\vrho^\infty]_{j'} \neq 0$ implies that $|[\vrho(t)]_{j'}| \to \infty$ and $\frac{|[\vrho(t)]_j|}{|[\vrho(t)]_{j'}|} \to 0$.
We now want to compute the ratio of \eqref{eq:proof-clsf-meta3-8} for $j$ and $j'$.
Using $|[\vrho(t)]_{j'}| \to \infty$, \eqref{eq:fix3}, and \eqref{eq:fix4}, we have 
\begin{align}
	\label{eq:proof-clsf-meta3-11}
	\lim_{t\to\infty} \frac{[\mU_1^T \vv_1(t)]_{j}^2 + [\mU_2^T \vv_2(t)]_{j}^2}{|[\vrho(t)]_{j}|}
	\geq
	\lim_{t\to\infty} \frac{[\mU_1^T \vv_1(t)]_{j'}^2 + [\mU_2^T \vv_2(t)]_{j'}^2}{|[\vrho(t)]_{j'}|} = 2.
\end{align}
%Using $\frac{|[\vrho(t)]_j|}{|[\vrho(t)]_{j'}|} \to 0$ and the fact that $x \mapsto \frac{2x^2 + \delta}{x \sqrt{x^2+\delta}}$ is a decreasing function of $x\geq 0$ for any $\delta > 0$, we have
%\begin{align}
%\label{eq:proof-clsf-meta3-11}
%    \frac{([\mU_1^T \vv_1(t)]_{j}^2 + [\mU_2^T \vv_2(t)]_{j}^2)/|[\vrho(t)]_{j}|}
%    {([\mU_1^T \vv_1(t)]_{j'}^2 + [\mU_2^T \vv_2(t)]_{j'}^2))/|[\vrho(t)]_{j'}|} \geq 1,
%\end{align}
%for any $t \geq t_0$, when $t_0$ is large enough. 
Combining $\frac{[\vs]_j}{[\vs]_{j'}} > 1$ and \eqref{eq:proof-clsf-meta3-11} to compute the ratio of \eqref{eq:proof-clsf-meta3-8} for $j$ and $j'$, there exists some $t_0 \geq 0$ such that for any $t \geq t_0$, we have
\begin{align}
\label{eq:proof-clsf-meta3-12}
    \frac{\left | \frac{d}{dt} [\vrho(t)]_j \right |/|[\vrho(t)]_j|}
    {\left | \frac{d}{dt} [\vrho(t)]_{j'} \right |/|[\vrho(t)]_{j'}|}
    > 1.
\end{align}
This implies that the ratio of the absolute value of time derivative of $[\vrho(t)]_j$ to the absolute value of current value of $[\vrho(t)]_j$ is strictly bigger than that of $[\vrho(t)]_{j'}$. 
Moreover, by the definition of $r(t)$, $\frac{d}{dt} [\vrho(t)]_j$ always has sign equal to $y$~\eqref{eq:proof-clsf-meta3-17}. Since this holds for all $t \geq t_0$, \eqref{eq:proof-clsf-meta3-12} results in a growth of $|[\vrho(t)]_j|$ that is exponentially faster than that of $|[\vrho(t)]_{j'}|$, so $[\vrho(t)]_j$ becomes a dominant component in $\vrho(t)$ as $t \to \infty$. This contradicts that $[\vrho^\infty]_j = 0$, hence the condition~\eqref{eq:proof-clsf-meta3-7} has to be satisfied.

So far, we have characterized some conditions~\eqref{eq:proof-clsf-meta3-16}, \eqref{eq:proof-clsf-meta3-6}, \eqref{eq:proof-clsf-meta3-7} that have to be satisfied by the limit direction $\vrho^\infty$ of $\vrho$. We now consider the following optimization problem and prove that these conditions are in fact the KKT conditions of the optimization problem.
Consider 
\begin{equation}
\label{eq:proof-clsf-meta3-13}
    \minimize_{\vrho \in \R^m}\quad \norm{\vrho}_{1}
    \quad
    \subjectto \quad y \vs^T \vrho \geq 1.
\end{equation}
The KKT condition of this problem is
\begin{equation*}
    \partial \norm{\vrho}_{1} \ni y \vs,
\end{equation*}
where the global subdifferential $\partial \norm{\cdot}_{1}$ is defined as
\begin{equation*}
    \partial \norm{\vrho}_{1}
    =
    \{ \vu \in \R^m \mid |[\vu]_j| \leq 1 \text{ for all } j \in [m], \text{ and } [\vrho]_j \neq 0 \implies [\vu]_j = \sign([\vrho]_j)\}.
\end{equation*}
We can check from~\eqref{eq:proof-clsf-meta3-16}, \eqref{eq:proof-clsf-meta3-6}, \eqref{eq:proof-clsf-meta3-7} that the that $\vrho^\infty$ satisfies the KKT condition up to scaling.

Now, how do we characterize $\vv_1^\infty$ and $\vv_2^\infty$ in terms of $\vrho^\infty$?
Let $\veta_1^\infty \defeq \mU_1^T \vv_1^\infty$ and $\veta_2^\infty \defeq \mU_2^T \vv_2^\infty$.
Then, $\vv_l^\infty = \mU_l \veta_l^\infty = \mU_l \mU_l^T \vv_l^\infty$ holds because any component orthogonal to the column space of $\mU_l$ stays unchanged while the component in the column space of $\mU_l$ diverges to infinity.
By \eqref{eq:proof-clsf-meta3-15}, $|\veta_1^\infty| = |\veta_2^\infty| = |\vrho^\infty|^{\odot 1/2}$.
By \eqref{eq:proof-clsf-meta3-14}, we have $\sign(\veta_1^\infty) = \sign(y)\odot\sign(\veta_2^\infty)$.

\subsection{Proof of Corollary~\ref{cor:clsf-conv-small}}
The proof of Corollary~\ref{cor:clsf-conv-small} boils down to characterizing the SVD of $\tMconv(\vx)$.

\subsubsection{The $k_1 = 1$ case}
First, it is straightforward to check that for $L=2$ and $k_1 = 1$, we have
\begin{equation*}
    \vbetaconv(\bThetaconv) = \vv_1 \vv_2.
\end{equation*}
Note that the parameter $\vv_1$ is now a scalar although we use a boldface letter.
For $k_1 = 1$, the data tensor is simply $\tMconv(\vx) = \vx^T$. Thus, we have $\mU_1 = 1$, $\mU_2 = \frac{\vx}{\ltwo{\vx}}$, and $\vs = \ltwos{\vx}$.
Substituting $\mU_1$ and $\mU_2$ to the theorem gives the condition on initial directions in Corollary~\ref{cor:clsf-conv-small}. Also, the theorem implies us that the limit direction $\vv_2^\infty$ of $\vv_2$ satisfies $\vv_2^\infty \propto y \vv_1^\infty \vx$. Using this, it is easy to check that
\begin{equation*}
    \vbetaconv(\bThetaconv^\infty) \propto \vv_1^\infty \vv_2^\infty \propto y\vx.
\end{equation*}

\subsubsection{The $k_1 = 2$ case}
First, it is straightforward to check that for $L=2$ and $k_1 = 2$, we have
\begin{equation}
\label{eq:proof-cor-clsf-conv-small-1}
    \vbetaconv(\bThetaconv)
    = \begin{bmatrix}
    [\vv_1]_1 & 0 & 0 & \cdots & 0 & [\vv_1]_2\\
    [\vv_1]_2 & [\vv_1]_1 & 0 & \cdots & 0 & 0\\
    0 & [\vv_1]_2 & [\vv_1]_1 & \cdots & 0 & 0 \\
    \vdots & \vdots & \vdots & \ddots & \vdots & \vdots\\
    0 & 0 & 0 & \cdots & [\vv_1]_1 & 0\\
    0 & 0 & 0 & \cdots & [\vv_1]_2 & [\vv_1]_1
    \end{bmatrix}
    \vv_2.
\end{equation}

For $k_1 = 2$, by definition, the data tensor is
\begin{align*}
    \tMconv(\vx) = \begin{bmatrix}
    \vx^T \\ \overleftarrow\vx^T
    \end{bmatrix},
\end{align*}
and it is straightforward to check that the SVD of this matrix is
\begin{align*}
    \tMconv(\vx) = \begin{bmatrix}
    \vx^T \\ \overleftarrow\vx^T
    \end{bmatrix}
    =
    \begin{bmatrix}
    \nicefrac{1}{\sqrt 2} & \nicefrac{1}{\sqrt 2}\\
    \nicefrac{1}{\sqrt 2} & -\nicefrac{1}{\sqrt 2}
    \end{bmatrix}
    \begin{bmatrix}
    \sqrt{\ltwo{\vx}^2 + \vx^T \overleftarrow \vx} & 0\\
    0 & \sqrt{\ltwo{\vx}^2 - \vx^T \overleftarrow \vx}
    \end{bmatrix}
    \begin{bmatrix}
    \frac{\vx^T + \overleftarrow\vx^T}{\sqrt 2 \sqrt{\ltwo{\vx}^2 + \vx^T \overleftarrow \vx}}\\
    \frac{\vx^T - \overleftarrow\vx^T}{\sqrt 2 \sqrt{\ltwo{\vx}^2 - \vx^T \overleftarrow \vx}}
    \end{bmatrix},
\end{align*}
so
\begin{equation*}
    \mU_1 = \begin{bmatrix}
    \nicefrac{1}{\sqrt 2} & \nicefrac{1}{\sqrt 2}\\
    \nicefrac{1}{\sqrt 2} & -\nicefrac{1}{\sqrt 2}
    \end{bmatrix},
    \mU_2 = \begin{bmatrix}
    \frac{\vx + \overleftarrow\vx}{\sqrt 2 \sqrt{\ltwo{\vx}^2 + \vx^T \overleftarrow \vx}} &
    \frac{\vx - \overleftarrow\vx}{\sqrt 2 \sqrt{\ltwo{\vx}^2 - \vx^T \overleftarrow \vx}}
    \end{bmatrix},
    \vs = \begin{bmatrix}
    \sqrt{\ltwo{\vx}^2 + \vx^T \overleftarrow \vx}\\
    \sqrt{\ltwo{\vx}^2 - \vx^T \overleftarrow \vx}
    \end{bmatrix}.
\end{equation*}
Substituting $\mU_1$ and $\mU_2$ to the theorem gives the conditions on initial directions.
Also, note that the maximum singular value depends on the sign of $\vx^T \overleftarrow \vx$.
Consider the optimization problem in the theorem statement:
\begin{equation*}
    \minimize\nolimits_{\vrho \in \R^m}\quad \norm{\vrho}_{1}
    \quad
    \subjectto \quad y \vs^T \vrho \geq 1.
\end{equation*}
If $\vx^T \overleftarrow \vx > 0$, then the solution $\vrho^\infty$ to this problem is in the direction of $[y~~0]$. Therefore, the limit directions $\vv_1^\infty$ and $\vv_2^\infty$ will be of the form
\begin{equation*}
    \vv_1^\infty \propto c_1 \begin{bmatrix}1\\1\end{bmatrix},~~
    \vv_2^\infty \propto c_2 (\vx+\overleftarrow \vx),
\end{equation*}
where $\sign(c_1)\sign(c_2) = \sign(y)$. Using \eqref{eq:proof-cor-clsf-conv-small-1}, it is straightforward to check that
\begin{equation*}
    \vbetaconv(\bThetaconv^\infty)
    \propto
    y
    \begin{bmatrix}
    1 & 0 & 0 & \cdots & 0 & 1\\
    1 & 1 & 0 & \cdots & 0 & 0\\
    0 & 1 & 1 & \cdots & 0 & 0 \\
    \vdots & \vdots & \vdots & \ddots & \vdots & \vdots\\
    0 & 0 & 0 & \cdots & 1 & 0\\
    0 & 0 & 0 & \cdots & 1 & 1
    \end{bmatrix}
    (\vx + \overleftarrow \vx)
    = y(2\vx + \overleftarrow \vx + \overrightarrow \vx).
\end{equation*}
Similarly, if $\vx^T \overleftarrow \vx < 0$, then the solution $\rho^\infty$ is in the direction of $[0~~y]$. Using \eqref{eq:proof-cor-clsf-conv-small-1}, we have
\begin{equation*}
    \vbetaconv(\bThetaconv^\infty)
    \propto
    y
    \begin{bmatrix}
    1 & 0 & 0 & \cdots & 0 & -1\\
    -1 & 1 & 0 & \cdots & 0 & 0\\
    0 & -1 & 1 & \cdots & 0 & 0 \\
    \vdots & \vdots & \vdots & \ddots & \vdots & \vdots\\
    0 & 0 & 0 & \cdots & 1 & 0\\
    0 & 0 & 0 & \cdots & -1 & 1
    \end{bmatrix}
    (\vx - \overleftarrow \vx)
    = y(2\vx - \overleftarrow \vx - \overrightarrow \vx).
\end{equation*}
\section{Proofs of Theorem~\ref{thm:meta-thm-reg-1}, Corollaries~\ref{cor:reg-diag}, \ref{cor:reg-conv-full} \& \ref{cor:reg-matrix-fac}, and Lemma~\ref{lem:ode-real}}
\label{sec:proof-meta-thm-reg-1}
\subsection{Proof of Lemma~\ref{lem:ode-real}}
In this subsection, we restate Lemma~\ref{lem:ode-real} and prove it.
\lemodereal*

\begin{proof}
For the proof, we omit the subscript $L$ for simplicity.
First, continuity (and also continuous differentiability) of $p(t)$ and $q(t)$ is straightforward because the RHSs of the ODEs are differentiable in $p$ and $q$. Next, define $\tilde p(t) = p(-t)$ and $\tilde q(t) = -q(-t)$. Then, one can show that $\tilde p$ and $\tilde q$ are also the solution of the ODE because
\begin{align*}
    \frac{d}{dt} \tilde p(t) &= \frac{d}{dt} p(-t) = -\dot p(-t) = -p(-t)^{L-2}q(-t) = \tilde p(t)^{L-2} \tilde q(t),\\
    \frac{d}{dt} \tilde q(t) &= - \frac{d}{dt} q(-t) = \dot q(-t) = p(-t)^{L-1} = \tilde p(t)^{L-1}.
\end{align*}
However, by the Picard-Lindel\"{o}f theorem, the solution has to be unique; this means that $p(t) = \tilde p(t) = p(-t)$ and $q(t) = \tilde q(t) = -q(-t)$, which proves that $p$ is even and $q$ is odd and also implies that the domain of $p$ and $q$ has to be of the form $(-c,c)$ (i.e. symmetric around the origin) and $h = p^{L-1}q$ is odd.

To show that $h$ is strictly increasing, it suffices to show that $p$ and $q$ are both strictly increasing on $[0,c)$. To this end, we show that $p(t) \geq 1$ for all $t \in [0,c)$.
First, due to the initial condition $p(0) = 1$ and continuity of $p$, there exists $\epsilon_1 > 0$ such that $p(t) > 0$ for all $t \in [0, \epsilon_1) \eqdef I_1$. 
This implies that $\dot q(t) = p(t)^{L-1} > 0$ for $t \in I_1 \setminus \{0 \}$, so $q$ is strictly increasing on $I_1$. Since $q(0)=0$, we have $q(t) > 0$ for $t \in I_1 \setminus \{0 \}$, which then implies that $\dot p(t) = p(t)^{L-2} q(t) > 0$. Therefore, $p$ is also strictly increasing on $I_1$; this then means $p(t) \geq 1$ for $t \in [0,\epsilon_1]$ because $p(0) = 1$. Now, due to $p(\epsilon_1) \geq 1$ and continuity of $p$, there exists $\epsilon_2 > \epsilon_1$ such that $p(t) > 0$ for all $t \in [\epsilon_1, \epsilon_2) \eqdef I_2$. Using the argument above for $I_2$ results in $p(t) \geq 1$ for $t \in [0, \epsilon_2]$. Repeating this until the end of the domain, we can show that $p(t) \geq 1$ holds for all $t \in [0, c)$.
By $p \geq 1$, we have $\dot q = p^{L-1} \geq 1$ on $[0,c)$, so $q$ is strictly increasing on $[0,c)$. Also, $q(t) > 0$ on $(0,c)$, so $\dot p = p^{L-2} q > 0$ on $(0,c)$ and $p$ is also strictly increasing on $[0,c)$. This proves that $h$ is strictly increasing on $[0,c)$, and also on $(-c,c)$ by oddity of $h$. 

Finally, it is left to show $\lim_{t \uparrow c} h(t) = \infty$ and $\lim_{t \downarrow -c} h(t)= -\infty$. If $c < \infty$, then this together with monotonicity implies that the limits hold. To see why, suppose $c<\infty$ and $\lim_{t \uparrow c} h(t) < \infty$. Then, $p$ and $q$ can be extended beyond $t\geq c$, which contradicts the fact that $(-c,c)$ is the maximal interval of existence of the solution. 
Next, consider the case $c = \infty$. From $p(t) \geq 1$, we have $\dot q(t) \geq 1$ for $t \geq 0$. This implies that $q(t) \geq t$ for $t \geq 0$. Now, $\dot p(t) \geq p(t)^{L-2} q(t) \geq t$, which gives $p(t) \geq \frac{t^2}{2}+1$ for $t \geq 0$. Therefore, we have 
\begin{align*}
\lim_{t \to \infty} h(t) = \lim_{t \to\infty} p(t)^{L-1}q(t) \geq \lim_{t \to \infty} \left ( \frac{t^2}{2}+1 \right )^{L-1} t = \infty,
\end{align*}
hence finishing the proof.
\end{proof}

\subsection{Proof of Theorem~\ref{thm:meta-thm-reg-1}}
\subsubsection{Convergence of loss to zero}
\label{sec:proof-meta-thm-reg-1-conv}
We first show that given the conditions on initialization, the training loss $\mc L(\bTheta(t))$ converges to zero. 
Recall from Section~\ref{sec:tensorform} that
\begin{align*}
    \dot \vv_l 
    = -\nabla_{\vv_l} \mc L(\bTheta)
    = \tM(-\mX^T \vr) \circ (\vv_1, \dots, \vv_{l-1}, \mI_{k_l}, \vv_{l+1}, \dots, \vv_L).
\end{align*}
Applying the structure~\eqref{eq:tensor-orth-decomp} in Assumption~\ref{assm:complex-diag-decomp}, we get
\begin{align*}
    \dot \vv_l 
    &= \tM(-\mX^T \vr) \circ (\vv_1, \dots, \vv_{l-1}, \mI_{k_l}, \vv_{l+1}, \dots, \vv_L)\\
    &= -\!\sum_{j=1}^m [\mS\mX^T\vr]_j (\vv_1^T[\mU_1]_{\cdot,j} \!\otimes \cdots \otimes \vv_{l-1}^T[\mU_{l-1}]_{\cdot,j} \!\otimes [\mU_l]_{\cdot, j} \!\otimes \vv_{l+1}^T[\mU_{l+1}]_{\cdot, j} \!\otimes \cdots \otimes \vv_{L}^T[\mU_{L}]_{\cdot, j})\\
    &= -\sum_{j=1}^m [\mS\mX^T\vr]_j \bigg(\prod_{k\neq l} [\mU_k^T \vv_k]_{j} \bigg) [\mU_l]_{\cdot, j}.
\end{align*}
Left-multiplying $\mU_l^T$ to both sides, we get
\begin{equation}
\label{eq:proof-reg-meta1-1}
    \mU_l^T \dot \vv_l = - \mS \mX^T \vr \odot \prod\nolimits_{k \neq l}^{\odot} \mU_k^T \vv_k,
\end{equation}
where $\prod^\odot$ denotes the product using entry-wise multiplication $\odot$.

Now consider the rate of growth for the second power of the $j$-th component of $\mU_l^T \vv_l$:
\begin{align*}
    \frac{d}{dt} [\mU_l^T \vv_l]_j^2
    = 2 [\mU_l^T \dot \vv_l]_j [\mU_l^T \vv_l]_j
    = -2 [\mS\mX^T \vr]_j \prod\nolimits_{k=1}^L [\mU_k^T \vv_k]_j
    = \frac{d}{dt} [\mU_{l'}^T \vv_{l'}]_j^2
\end{align*}
for any $l' \in [L]$. Thus, for any $j\in [m]$, the second power of the $j$-th components in $\mU_l^T \vv_l$ grow at the same rate for each layer $l \in [L]$. This means that the gap between any two different layers stays constant for all $t \geq 0$. Combining this with our conditions on initial directions, we have
\begin{equation*}
%\label{eq:proof-clsf-meta2-11}
%\begin{aligned}
    [\mU_l^T \vv_l(t)]_j^2 - [\mU_L^T \vv_L(t)]_j^2
    = [\mU_l^T \vv_l(0)]_j^2 - [\mU_L^T \vv_L(0)]_j^2
    = \alpha^2 [\bar \veta]_j^2 \geq \alpha^2 \lambda,
%\end{aligned}
\end{equation*}
for any $j \in [m]$, $l \in [L-1]$, and $t \geq 0$. This inequality also implies
\begin{equation}
\label{eq:proof-reg-meta1-2}
    [\mU_l^T \vv_l(t)]_j^2 
    \geq 
    [\mU_L^T \vv_L(t)]_j^2 + \alpha^2 \lambda
    \geq
    \alpha^2 \lambda.
\end{equation}

Let us now consider the time derivative of $\mc L(\bTheta(t))$.
We have the following chain of upper bounds on the time derivative:
\begin{align}
    \frac{d}{dt} \mc L(\bTheta(t))
    &= \nabla_{\bTheta} \mc L(\bTheta(t))^T \dot \bTheta(t)
    = -\ltwos{\nabla_{\bTheta} \mc L(\bTheta(t))}^2\notag\\
    &\leq -\ltwos{\nabla_{\vv_L} \mc L(\bTheta(t))}^2
    = -\ltwos{\dot \vv_L(t)}^2\notag\\
    &\leqwtxt{(a)} -\ltwos{\mU_L^T \dot \vv_L(t)}^2
    \eqwtxt{(b)} -\ltwo{\mS \mX^T \vr(t) \odot \prod\nolimits_{k \neq L}^{\odot} \mU_k^T \vv_k(t)}^2\notag\\
    &=-\sum\nolimits_{j=1}^m [\mS \mX^T \vr(t)]_j^2 \prod\nolimits_{k \neq L} [\mU_k^T \vv_k(t)]_j^2\notag\\
    &\leqwtxt{(c)} -\alpha^{2L-2} \lambda^{L-1} \sum\nolimits_{j=1}^m [\mS \mX^T \vr(t)]_j^2 \notag\\
    &= -\alpha^{2L-2} \lambda^{L-1} \ltwos{\mS \mX^T \vr(t)}^2\notag\\
    &\leqwtxt{(d)} -\alpha^{2L-2} \lambda^{L-1} s_{\min}(\mS)^2 s_{\min}(\mX)^2
    \ltwos{\vr(t)}^2,\notag\\%\label{eq:proof-clsf-meta2-3}
    &= -2 \alpha^{2L-2} \lambda^{L-1} s_{\min}(\mS)^2 s_{\min}(\mX)^2 \mc L(\bTheta(t)),\label{eq:proof-reg-meta1-3}
\end{align}
where (a) used the fact that $\ltwos{\dot \vv_L(t)}^2 \geq \ltwos{\mU_L \mU_L^T \dot \vv_L(t)}^2$ because it is a projection onto a subspace, and $\ltwos{\mU_L \mU_L^T \dot \vv_L(t)}^2 = \ltwos{\mU_L^T \dot \vv_L(t)}^2$ because $\mU_L^T \mU_L = \mI_{k_L}$; (b) is due to~\eqref{eq:proof-reg-meta1-1}; (c) is due to~\eqref{eq:proof-reg-meta1-2}; and (d) used the fact that $\mS \in \R^{m \times d}$ and $\mX^T \in \R^{d \times n}$ are matrices that have full column rank, so for any $\vz \in \C^n$, we can use $\ltwos{\mS \mX^T \vz} \geq s_{\min}(\mS) s_{\min}(\mX) \ltwos{\vz}$ where $s_{\min}(\cdot)$ denotes the minimum singular value of a matrix.

From~\eqref{eq:proof-reg-meta1-3}, we get
\begin{equation}
\label{eq:proof-reg-meta1-7}
    \mc L(\bTheta(t)) \leq \mc L(\bTheta(0)) \exp(-2 \alpha^{2L-2} \lambda^{L-1} s_{\min}(\mS)^2 s_{\min}(\mX)^2 t),
\end{equation}
so that $\mc L(\bTheta(t)) \to 0$ as $t \to \infty$.


\subsubsection{Characterizing the limit point}
Now, we move on to characterize the limit points of the gradient flow. First, by defining a ``transformed'' version of the parameters $\veta_l(t) \defeq \mU_l^T \vv_l(t)$ and using \eqref{eq:proof-reg-meta1-1}, one can define an equivalent system of ODEs:
\begin{equation}
\label{eq:proof-reg-meta1-4}
\begin{aligned}
    \dot \veta_l &= -\mS \mX^T \vr \odot \prod\nolimits_{k \neq l}^{\odot} \veta_k \text{ for } l \in [L],\\
    \veta_l(0) &= \alpha \bar \veta \text { for } l \in [L-1],~~\veta_L(0) = \zeros.
\end{aligned}
\end{equation}

Using Lemma~\ref{lem:ode-real}, it is straightforward to verify that the solution to \eqref{eq:proof-reg-meta1-4} has the following form. For odd $L$, we have
\begin{equation}
\label{eq:proof-reg-meta1-5}
\begin{aligned}
    \veta_l(t) &= \alpha \bar \veta \odot p_L \left (-\alpha^{L-2} |\bar \veta|^{\odot L-2} \odot \mS \mX^T \int_0^t \vr(\tau)d\tau \right) \text{ for } l \in [L-1],\\
    \veta_L(t) &= \alpha |\bar \veta| \odot q_L\left (- \alpha^{L-2} |\bar \veta|^{\odot L-2} \odot \mS \mX^T \int_0^t \vr(\tau)d\tau \right).
\end{aligned}
\end{equation}

Similarly, for even $L$, the solution for \eqref{eq:proof-reg-meta1-4} satisfies
\begin{equation}
\label{eq:proof-reg-meta1-6}
\begin{aligned}
    \veta_l(t) &= \alpha \bar \veta \odot p_L\left (-\alpha^{L-2} \bar \veta^{\odot L-2} \odot \mS \mX^T \int_0^t \vr(\tau)d\tau \right) \text{ for } l \in [L-1],\\
    \veta_L(t) &= \alpha \bar \veta \odot q_L\left (-\alpha^{L-2} \bar \veta^{\odot L-2} \odot \mS \mX^T \int_0^t \vr(\tau)d\tau \right).
\end{aligned}
\end{equation}
Now that we know how the solutions $\veta_l$ look like, let us see how these relate to the linear coefficients of the network. 
By Assumption~\ref{assm:complex-diag-decomp}, we have
\begin{align*}
    f(\vx; \bTheta) 
    &= \tM(\vx) \circ (\vv_1, \dots, \vv_L)
    = \sum_{j=1}^m [\mS \vx]_j \prod_{l=1}^L [\mU_l^T \vv_l]_j\\
    &= \bigg[ \sum_{j=1}^m \bigg(\prod_{l=1}^L [\veta_l]_j \bigg) [\mS]_{j,\cdot} \bigg ]\vx
    = \vx^T \mS^T \bigg(\prod\nolimits_{l \in [L]}^{\odot} \veta_l \bigg)
    = \vx^T \mS^T \vrho.
\end{align*}
Here, we defined $\vrho \defeq \prod\nolimits_{l \in [L]}^{\odot} \veta_l \in \R^m$.
Therefore, the linear coefficients of the network can be written as $\vbeta(\bTheta(t)) = \mS^T \vrho(t)$.
From the solutions~\eqref{eq:proof-reg-meta1-5} and \eqref{eq:proof-reg-meta1-6}, we can write
\begin{equation*}
    \vrho(t) = \prod_{i=1}^L \veta_l(t)
    = \alpha^{L} |\bar \veta|^{\odot L} \odot h_L \left (
    -\alpha^{L-2} |\bar \veta|^{\odot L-2} \odot \mS \mX^T \int_0^t \vr(\tau)d\tau
    \right ),
\end{equation*}
where $h_L \defeq p_L^{L-1} q_L$, defined in Lemma~\ref{lem:ode-real}.
By the convergence of the loss to zero~\eqref{eq:proof-reg-meta1-7}, we have $\lim_{t \to \infty} \mX \vbeta(\bTheta(t)) = \vy$. Therefore,
\begin{align}
\label{eq:proof-reg-meta1-9}
    \mX \mS^T \underbrace{\left ( 
    \alpha^L |\bar \veta|^{\odot L} \odot h_L \left (
    -\alpha^{L-2} |\bar \veta|^{\odot L-2} \odot \mS \mX^T \int_0^\infty \vr(\tau)d\tau
    \right )
    \right )}_{\eqdef \vrho^\infty}
    = \vy.
\end{align}
Next, we will show that $\vrho^\infty$ is in fact the solution of the following optimization problem
\begin{align}
\label{eq:proof-reg-meta1-8}
    \minimize_{\vrho \in \R^m}\quad Q_{L, \alpha, \bar \veta}(\vrho)
    \quad
    \subjectto \quad \mX \mS^T \vrho= \vy,
\end{align}
where $Q_{L, \alpha, \bar \veta} : \R^m \to \R$ is a norm-like function  defined using $H_L (t) \defeq \int_0^t h_L^{-1} (\tau) d\tau$:
\begin{align*}
    Q_{L, \alpha, \bar \veta}(\vrho) = \alpha^2  \sum_{j=1}^m [\bar \veta]_j^2 H_L \left ( 
    \frac{[\vrho]_j}{\alpha^L |[\bar \veta]_j|^L}
    \right ).
\end{align*}
Note that the KKT conditions for \eqref{eq:proof-reg-meta1-8} are
\begin{align*}
    \mX \mS^T \vrho = \vy, \quad \nabla_\vrho Q_{L, \alpha, \bar \veta}(\vrho) = \mS \mX^T \vnu,
\end{align*}
for some $\vnu \in \R^n$. It is clear from \eqref{eq:proof-reg-meta1-9} that $\vrho^\infty$ satisfies the first condition (primal feasibility), so let us check the other one.
Through a straightforward calculation, we get
\begin{align*}
    \nabla_\vrho Q_{L, \alpha, \bar \veta}(\vrho)
    = \alpha^{2-L} |\bar \veta|^{\odot 2-L} \odot
    h_L^{-1} \left ( \alpha^{-L} |\bar \veta|^{\odot (-L)} \odot \vrho \right ).
\end{align*}
Equating this with $\mS \mX^T \vnu$ gives
\begin{align*}
    &\alpha^{2-L} |\bar \veta|^{\odot 2-L} \odot
    h_L^{-1} \left ( \alpha^{-L} |\bar \veta|^{\odot(-L)} \odot \vrho \right ) = \mS \mX^T \vnu\\
    \iff~&
    h_L^{-1} \left ( \alpha^{-L} |\bar \veta|^{\odot(-L)} \odot \vrho \right )
    = \alpha^{L-2} |\bar \veta|^{\odot L-2} \odot \mS \mX^T \vnu\\
    \iff~&\vrho = \alpha^L |\bar \veta|^{\odot L} \odot
    h_L \left ( \alpha^{L-2} |\bar \veta|^{\odot L-2} \odot \mS \mX^T \vnu \right ).
\end{align*}
Hence, by setting $\vnu = - \int_0^\infty \vr(\tau) d\tau$, $\vrho^\infty$ satisfies the second KKT condition as well.
Also, if $\mS$ is invertible, we can substitute $\vrho = \mS^{-T} \vz$ to \eqref{eq:proof-reg-meta1-8} to get the last statement of the theorem.
This finishes the proof.


\subsection{Proof of Corollary~\ref{cor:reg-diag}}
The proof is a direct consequence of the fact that Assumption~\ref{assm:complex-diag-decomp} holds with $\mS = \mU_1 = \dots = \mU_L = \mI_d$ for linear diagonal networks. Hence, the proof is the same as Corollary~\ref{cor:clsf-diag}, proved in Appendix~\ref{sec:proof-cor-clsf-diag}.


\subsection{Proof of Corollary~\ref{cor:reg-conv-full}}
\label{sec:proof-reg-conv-full}
We start by showing the DFT of a real and even vector is also real and even. Suppose that $\vx \in \reals^d$ is real and even. First,
\begin{align*}
    [\mF \vx]_j &= \frac{1}{\sqrt{d}} \sum_{k=1}^d [\vx]_k \exp \left (-\frac{\sqrt{-1} \cdot 2 \pi (j-1)(k-1)}{d} \right)\\
    &=
    \frac{1}{\sqrt{d}} \sum_{k=1}^d [\vx]_k \cos \left (-\frac{2 \pi (j-1)(k-1)}{d} \right)
    + 
    \frac{\sqrt{-1}}{\sqrt{d}} \sum_{k=1}^d [\vx]_k \sin \left (-\frac{2 \pi (j-1)(k-1)}{d} \right)\\
    &=
    \frac{1}{\sqrt{d}} \sum_{k=1}^d [\vx]_k \cos \left (-\frac{2 \pi (j-1)(k-1)}{d} \right) \in \reals,
\end{align*}
for all $j \in [d]$.
To prove that $\mF\vx$ is even, for $j = 0, \dots, \lfloor \frac{d-3}{2} \rfloor$, we have
\begin{align*}
    [\mF \vx]_{j+2} &= \frac{1}{\sqrt{d}} \sum_{k=1}^d [\vx]_k \cos \left (-\frac{2 \pi (j+1)(k-1)}{d} \right)\\
    &=
    \frac{1}{\sqrt{d}} \sum_{k=1}^d [\vx]_k \cos \left (2\pi(k-1)-\frac{2 \pi (j+1)(k-1)}{d} \right)\\
    &=
    \frac{1}{\sqrt{d}} \sum_{k=1}^d [\vx]_k \cos \left (\frac{2 \pi (d-j-1)(k-1)}{d} \right)\\
    &=
    \frac{1}{\sqrt{d}} \sum_{k=1}^d [\vx]_k \cos \left (-\frac{2 \pi (d-j-1)(k-1)}{d} \right)\\
    &= [\mF \vx]_{d-j}.
\end{align*}

It is proved in Appendix~\ref{sec:proof-cor-clsf-conv-full} that linear full-length convolutional networks ($k_1 = \dots = k_L = d$) satisfy Assumption~\ref{assm:complex-diag-decomp} with
$\mS = d^{\frac{L-1}{2}} \mF$ and $\mU_1 = \dots = \mU_L = \mF^*$, where $\mF \in \C^{d \times d}$ is the matrix of discrete Fourier transform basis $[\mF]_{j,k} = \frac{1}{\sqrt{d}} \exp(-\frac{\sqrt{-1} \cdot 2 \pi (j-1)(k-1)}{d})$ and $\mF^*$ is the complex conjugate of $\mF$.

The proof of convergence of loss to zero in Appendix~\ref{sec:proof-meta-thm-reg-1-conv} is written for real matrices $\mS, \mU_1, \dots, \mU_L$, but we can actually apply the same argument as in Appendix~\ref{sec:proof-meta-thm-clsf-2-conv} and prove that the loss converges to zero, even in the case where $\mS, \mU_1, \dots, \mU_L$ are complex.

Next, since $\mU_l$'s are complex, we can write the system of ODE as (see \eqref{eq:proof-clsf-meta2-1} for its derivation)
\begin{equation}
\label{eq:proof-cor-reg-conv-full-1}
    \mF \dot \vw_l = - d^{\frac{L-1}{2}} \mF \mX^T \vr \odot \prod\nolimits_{k \neq l}^{\odot} \mF^* \vw_k,
\end{equation}
Since all data points $\vx_i$ and initialization $\vw_l(0)$ are real and even, we have that $\mF \mX^T \vr$ is real and even, and $\mF^* \vw_l(0) = \mF \vw_l(0)$'s are real and even. By \eqref{eq:proof-cor-reg-conv-full-1}, we see that the time derivatives of $\mF \vw_l$ are also real and even. Thus, the parameters $\vw_l(t)$ are all real and even for all $t \geq 0$.
From this observation, we can define
$\veta_l(t) \defeq \mF \vw_l(t)$, $\bar \veta \defeq \mF \bar \vw$, and $\mS \defeq d^{\frac{L-1}{2}} \Re(\mF)$, which are all real by the even symmetry.
Then, starting from \eqref{eq:proof-reg-meta1-4}, the proof goes through.


\subsection{Proof of Corollary~\ref{cor:reg-matrix-fac}}
Since the sensor matrices $\mA_1, \dots, \mA_n$ commute, they are simultaneously diagonalizable with a real unitary matrix $\mU \in \reals^{d \times d}$, i.e., $\mU^T \mA_i \mU$'s are diagonal matrices for all $i \in [n]$.
From the deep matrix sensing problem~\eqref{eq:matrix-sensing}, we can compute $\nabla_{\mW_l} \mc L_{\rm ms}$, which gives the gradient flow dynamics of $\mW_l$.
\begin{equation*}
    \dot \mW_l = -\nabla_{\mW_l} \mc L_{\rm ms}
    = - \mW_{l-1}^T \cdots \mW_1^T (\sum\nolimits_{i=1}^n r_i \mA_i) \mW_L^T \cdots \mW_{l+1}^T,
\end{equation*}
where $r_i = \< \mA_i, \mW_1\cdots \mW_L \> - y_i$ is the residual for the $i$-th sensor matrix.
If we left-multiply $\mU^T$ and right-multiply $\mU$ to both sides, we get
\begin{equation}
\label{eq:proof-cor-reg-matrix-fac-1}
    \mU^T \dot \mW_l \mU 
    = - \mU^T \mW_{l-1}^T \mU \cdots \mU^T \mW_1^T \mU (\sum\nolimits_{i=1}^n r_i \mU^T \mA_i \mU) \mU^T \mW_L^T \mU \cdots \mU^T \mW_{l+1}^T \mU.
\end{equation}
If $\mU^T \mW_{k}^T \mU$ is a diagonal matrix for all $k \neq l$, then $\mU^T \dot \mW_l \mU$ is also a diagonal matrix.
Note also that, since $\mW_l(0) = \alpha \mI_d = \alpha \mU \mU^T$ for $l \in [L-1]$, the product $\mU^T\mW_l\mU$ is a diagonal matrix at initialization. 
These observations imply that $\mW_l(t)$'s are all diagonalizable with $\mU$ for all $t \geq 0$.

Now, define $\vv_l(t) = \eig(\mW_l(t))$, i.e., $\mU^T \mW_l \mU = \diag(\vv_l)$. Also, let $\vx_i = \eig(\mA_i)$. Then, \eqref{eq:proof-cor-reg-matrix-fac-1} can be written as
\begin{equation*}
    \dot \vv_l = -(\sum\nolimits_{i=1}^n r_i \vx_i) \odot \prod\nolimits_{k\neq l}^{\odot} \vv_k.
\end{equation*}
Therefore, this is equivalent to the regression problem with linear diagonal networks, initialized at $\vv_l(0) = \alpha \ones$ for $l \in [L-1]$ and $\vv_L(0) = \zeros$.
Given this equivalence, Corollary~\ref{cor:reg-matrix-fac} can be implied from Corollary~\ref{cor:reg-diag}.



\section{Proof of Theorem~\ref{thm:meta-thm-reg-2}}
\label{sec:proof-meta-thm-reg-2}
\subsection{Convergence of loss to zero}
We first show that given the conditions on initialization, the training loss $\mc L(\bTheta(t))$ converges to zero.
Since $L=2$ and $\tM(\vx) = \mU_1\diag(\vs)\mU_2^T$, we can write the gradient flow dynamics from Section~\ref{sec:tensorform} as
\begin{equation}
\label{eq:proof-reg-meta2-1}
\begin{aligned}
    \dot \vv_1
    &= -\tM(\mX^T \vr) \circ (\mI_{k_1}, \vv_2)
    = -r \mU_1 \diag(\vs) \mU_2^T \vv_2,\\
    \dot \vv_2
    &= -\tM(\mX^T \vr) \circ (\vv_1, \mI_{k_2})
    = -r \mU_2 \diag(\vs) \mU_1^T \vv_1,
\end{aligned}
\end{equation}
where $r(t) = f(\vx;\bTheta(t))-y$ is the residual of the data point $(\vx,y)$.
From \eqref{eq:proof-reg-meta2-1} we get
\begin{equation}
\label{eq:proof-reg-meta2-2}
    \mU_1^T \dot \vv_1 = -r \vs \odot \mU_2^T \vv_2,~~
    \mU_2^T \dot \vv_2 = -r \vs \odot \mU_1^T \vv_1.
\end{equation}

Now consider the rate of growth for the $j$-th component of $\mU_1^T \vv_1$ squared:
\begin{equation*}
%\label{eq:proof-clsf-meta3-15}
    \frac{d}{dt} [\mU_1^T \vv_1]_j^2
    = 2 [\mU_1^T \vv_1]_j [\mU_1^T \dot \vv_1]_j
    = -2 r [\vs]_j [\mU_1^T \vv_1]_j [\mU_2^T \vv_2]_j
    = \frac{d}{dt} [\mU_{2}^T \vv_{2}]_j^2.
\end{equation*}
So for any $j\in [m]$, $[\mU_1^T \vv_1]_j^2$ and $[\mU_2^T \vv_2]_j^2$ grow at the same rate. This means that the gap between the two layers stays constant for all $t \geq 0$. Combining this with our conditions on initial directions,
\begin{equation*}
%\label{eq:proof-reg-meta2-3}
\begin{aligned}
    [\mU_1^T \vv_1(t)]_j^2 - [\mU_2^T \vv_2(t)]_j^2
    &= [\mU_1^T \vv_1(0)]_j^2 - [\mU_2^T \vv_2(0)]_j^2\\
    &= \alpha^2 [\mU_1^T \bar \vv_1]_j^2 - \alpha^2 [\mU_2^T \bar \vv_2]_j^2 \geq \alpha^2 \lambda,
\end{aligned}
\end{equation*}
for any $j \in [m]$ and $t \geq 0$. This inequality implies
\begin{equation}
\label{eq:proof-reg-meta2-3}
    [\mU_1^T \vv_1(t)]_j^2 
    \geq 
    [\mU_2^T \vv_2(t)]_j^2 + \alpha^2 \lambda
    \geq
    \alpha^2 \lambda.
\end{equation}

Let us now consider the time derivative of $\mc L(\bTheta(t))$.
We have the following chain of upper bounds on the time derivative:
\begin{align}
    \frac{d}{dt} \mc L(\bTheta(t))
    &= \nabla_{\bTheta} \mc L(\bTheta(t))^T \dot \bTheta(t)
    = -\ltwos{\nabla_{\bTheta} \mc L(\bTheta(t))}^2\notag\\
    &\leq -\ltwos{\nabla_{\vv_2} \mc L(\bTheta(t))}^2
    = -\ltwos{\dot \vv_2(t)}^2\notag\\
    &\leqwtxt{(a)} -\ltwos{\mU_2^T \dot \vv_2(t)}^2
    \eqwtxt{(b)} -r(t)^2 \ltwo{\vs \odot \mU_1^T \vv_1(t)}^2\notag\\
    &=-r(t)^2 \sum\nolimits_{j=1}^m [\vs]_j^2 [\mU_1^T \vv_1(t)]_j^2\notag\\
    &\leqwtxt{(c)} -\alpha^{2} \lambda r(t)^2 \sum\nolimits_{j=1}^m [\vs]_j^2 \notag\\
    &= -2 \alpha^{2} \lambda \ltwos{\vs}^2 \mc L(\bTheta(t))\notag,
\end{align}
where (a) used the fact that $\ltwos{\dot \vv_2(t)}^2 \geq \ltwos{\mU_2 \mU_2^T \dot \vv_2(t)}^2$ because it is a projection onto a subspace, and $\ltwos{\mU_2 \mU_2^T \dot \vv_L(t)}^2 = \ltwos{\mU_2^T \dot \vv_2(t)}^2$ because $\mU_2^T \mU_2 = \mI_{k_2}$; (b) is due to~\eqref{eq:proof-reg-meta2-2}; (c) is due to~\eqref{eq:proof-reg-meta2-3}.
From this, we get
\begin{equation}
\label{eq:proof-reg-meta2-9}
    \mc L(\bTheta(t)) \leq 
    \mc L(\bTheta(0)) \exp(-2\alpha^{2} \lambda \ltwos{\vs}^2 t).
\end{equation}
Therefore, $\mc L(\bTheta(t)) \to 0$ as $t \to \infty$.

\subsection{Characterizing the limit point}
Now, we move on to characterize the limit points of the gradient flow. First, note that any changes made in $\vv_l$ over time are in the subspace spanned by the columns of $\mU_l$. Therefore, any component in the initialization $\vv_l(0) = \alpha \bar \vv_l$ that is orthogonal to the column space of $\mU_l$ stays constant.

So, we can focus on the evolution of $\vv_l$ in the column space of $\mU_l$; this can be done by defining a ``transformed'' version of the parameters $\veta_l(t) \defeq \mU_l^T \vv_l(t)$ and using \eqref{eq:proof-reg-meta2-2}, one can define an equivalent system of ODEs:
\begin{equation}
\label{eq:proof-reg-meta2-4}
\begin{aligned}
    \dot \veta_1 = -r \vs \odot \veta_2,~&~
    \dot \veta_2 = -r \vs \odot \veta_1,\\
    \veta_1(0) = \alpha \bar \veta_1,~&~
    \veta_2(0) = \alpha \bar \veta_2,
\end{aligned}
\end{equation}
where $\bar \veta_1 \defeq \mU_1^T \bar \vv_1$, $\bar \veta_2 \defeq \mU_2^T \bar \vv_2$.
It is straightforward to verify that the solution to \eqref{eq:proof-reg-meta2-4} has the following form.
\begin{equation}
\label{eq:proof-reg-meta2-7}
\begin{aligned}
    \veta_1(t) &= 
    \alpha \bar \veta_1 \odot \cosh \left (-\vs \int_0^t r(\tau)d\tau \right) 
    + \alpha \bar \veta_2 \odot \sinh \left (-\vs \int_0^t r(\tau)d\tau \right)
    ,\\
    \veta_2(t) &= 
    \alpha \bar \veta_1 \odot \sinh \left (-\vs \int_0^t r(\tau)d\tau \right) 
    + \alpha \bar \veta_2 \odot \cosh \left (-\vs \int_0^t r(\tau)d\tau \right).
\end{aligned}
\end{equation}
By the convergence of the loss to zero~\eqref{eq:proof-reg-meta2-9}, we have $\lim_{t \to \infty} f(\vx; \bTheta(t)) = y$. Note that $f(\vx; \bTheta(t))$ can be written as
\begin{align*}
    f(\vx; \bTheta(t))
    &= \tM(\vx) \circ (\vv_1(t), \vv_2(t))
    = \vv_1(t)^T \tM(\vx) \vv_2(t)\\
    &= \vv_1(t)^T \mU_1 \diag(\vs) \mU_2^T \vv_2(t)
    = \vs^T (\veta_1(t) \odot \veta_2(t)).
\end{align*}
Therefore,
\begin{align}
    &~\lim_{t \to \infty} f(\vx; \bTheta(t)) = \lim_{t \to \infty} \vs^T (\veta_1(t) \odot \veta_2(t))\notag \\
    =&~ \alpha^2 \vs^T \bigg[ 
    (\bar \veta_1^{\odot 2} + \bar \veta_2^{\odot 2}) 
    \odot \cosh \left (-\vs \int_0^\infty r(\tau)d\tau \right) 
    \odot \sinh \left (-\vs \int_0^\infty r(\tau)d\tau \right)\notag\\
    &\quad\quad\quad+(\bar \veta_1 \odot \bar \veta_2) \odot \left ( \cosh^{\odot 2} \left (-\vs \int_0^\infty r(\tau)d\tau \right) + \sinh^{\odot 2} \left (-\vs \int_0^\infty r(\tau)d\tau \right)\right) \bigg ]\notag\\
    =&~ \alpha^2 \vs^T \bigg [
    \frac{\bar \veta_1^{\odot 2} + \bar \veta_2^{\odot 2}}{2}
    \odot \sinh \left (-2 \vs \int_0^\infty r(\tau)d\tau \right) 
    +(\bar \veta_1 \odot \bar \veta_2) 
    \odot \cosh \left (-2 \vs \int_0^\infty r(\tau)d\tau \right) 
    \bigg ]\notag\\
    =&~\alpha^2 \sum_{j=1}^m 
    [\vs]_j \left (\frac{[\bar \veta_1]_j^2+[\bar \veta_2]_j^2}{2} \sinh\left (2[\vs]_j \nu \right ) 
    + [\bar \veta_1]_j [\bar \veta_2]_j \cosh\left (2[\vs]_j \nu \right) \right)\notag\\
    =&~y,\label{eq:proof-reg-meta2-6}
\end{align}
where we defined $\nu \defeq -\int_0^\infty r(\tau) d \tau$.
Consider the function $\nu \mapsto a \sinh(\nu) + b\cosh(\nu)$. This is a strictly increasing function if $a>|b|$. Note also that 
\begin{equation}
\label{eq:proof-reg-meta2-5}
    \frac{[\bar \veta_1]_j^2 + [\bar \veta_2]_j^2}{2} \geq |[\bar \veta_1]_j [\bar \veta_2]_j|,
\end{equation}
which holds with equality if and only if $|[\bar \veta_1]_j| = |[\bar \veta_2]_j|$. However, recall from our assumptions on initialization that
$[\bar \veta_1]_j^2-[\bar \veta_2]_j^2 \geq \lambda > 0$, so \eqref{eq:proof-reg-meta2-5} can only hold with strict inequality. Therefore,
\begin{equation*}
    g(\nu) \defeq \sum_{j=1}^m [\vs]_j \left (\frac{[\bar \veta_1]_j^2+[\bar \veta_2]_j^2}{2} \sinh(2[\vs]_j \nu) + [\bar \veta_1]_j [\bar \veta_2]_j \cosh(2[\vs]_j \nu) \right)
\end{equation*}
is a strictly increasing (hence invertible) function because it is a sum of $m$ strictly increasing function. Using this $g(\nu)$, \eqref{eq:proof-reg-meta2-6} can be written as $\alpha^2 g(\nu) = y$, and by using the inverse of $g$, we have \begin{equation}
\label{eq:proof-reg-meta2-8}
    \nu = - \int_0^\infty r(\tau) d\tau = g^{-1} \left( \frac{y}{\alpha^2} \right ).
\end{equation}
Plugging \eqref{eq:proof-reg-meta2-8} into \eqref{eq:proof-reg-meta2-7}, we get
\begin{align*}
    &~\lim_{t\to\infty} \vv_1(t) \\
    =&~ \mU_1 \lim_{t \to \infty} \veta_1(t) + \alpha (\mI_{k_1} - \mU_1 \mU_1^T ) \bar \vv_1\\
    =&~ \alpha \mU_1 \left ( 
    \bar \veta_1 \odot \cosh \left ( g^{-1} \left( \frac{y}{\alpha^2} \right ) \vs \right) 
    +  \bar \veta_2 \odot \sinh \left ( g^{-1} \left( \frac{y}{\alpha^2} \right ) \vs \right ) 
    \right )
    + \alpha (\mI_{k_1} - \mU_1 \mU_1^T )\bar \vv_1,\\
    &~\lim_{t\to\infty} \vv_2(t) \\
    =&~ \mU_2 \lim_{t \to \infty} \veta_2(t) + \alpha (\mI_{k_2} - \mU_2 \mU_2^T ) \bar \vv_2\\
    =&~ \alpha \mU_2 \left ( 
    \bar \veta_1 \odot \sinh \left ( g^{-1} \left( \frac{y}{\alpha^2} \right ) \vs \right) 
    +  \bar \veta_2 \odot \cosh \left ( g^{-1} \left( \frac{y}{\alpha^2} \right ) \vs \right ) 
    \right )
    + \alpha (\mI_{k_2} - \mU_2 \mU_2^T )\bar \vv_2.
\end{align*}
This finishes the proof.

\section{Proof of Theorem~\ref{thm:fc-reg}}
\label{sec:proof-thm-fc-reg}
\subsection{Proof of Theorem~\ref{thm:fc-reg}.\ref{item:thm-fc-reg-1}: convergence of loss to zero}
We first show that given the conditions on initialization, the training loss $\mc L(\bTheta(t))$ converges to zero.
Recall from \eqref{eq:fcnet-linear} that the linear fully-connected network can be written as
\begin{equation*}
    \ffc(\vx; \bThetafc) = \vx^T \mW_1 \mW_2 \cdots \mW_{L-1} \vw_L.
\end{equation*}
From the definition of the training loss $\mc L$, it is straightforward to check that the gradient flow dynamics read
\begin{equation}
\label{eq:proof-fc-reg-1}
\begin{aligned}
    \dot \mW_l &= - \nabla_{\mW_l} \mc L(\bThetafc) = - \mW_{l-1}^T \cdots \mW_1^T \mX^T \vr \vw_L^T \mW_{L-1}^T \cdots \mW_{l+1}^T \text { for } l \in [L-1],\\
    \dot \vw_L &= - \nabla_{\vw_L} \mc L(\bThetafc) = - \mW_{L-1}^T \cdots \mW_1^T \mX^T \vr,\\
    \mW_l(0) &= \alpha \bar \mW_l \text { for } l \in [L-1],\\
    \vw_L(0) &= \alpha \bar \vw_L,
\end{aligned}    
\end{equation}
where $\vr \in \R^n$ is the residual vector satisfying $[\vr]_i = \ffc(\vx_i;\bThetafc) - y_i$, as defined in Section~\ref{sec:tensorform}.
From~\eqref{eq:proof-fc-reg-1}, we have
\begin{align*}
    &\mW_l^T \dot \mW_l = \dot \mW_{l+1} \mW_{l+1}^T = - \mW_l^T \cdots \mW_1^T \mX^T \vr \vw_L^T \mW_{L-1}^T \cdots \mW_{l+1}^T,\\
    &\dot \mW_l^T \mW_l = \mW_{l+1} \dot \mW_{l+1}^T = - \mW_{l+1} \cdots \mW_{L-1} \vw_L \vr^T \mX \mW_1 \cdots \mW_l,
\end{align*}
for any $l \in [L-2]$. From this, we have
\begin{align*}
    \frac{d}{dt} \mW_l^T \mW_l = \frac{d}{dt} \mW_{l+1} \mW_{l+1}^T,
\end{align*}
and thus
\begin{equation}
\label{eq:proof-fc-reg-3}
\begin{aligned}
    \mW_l(t)^T \mW_l(t) - \mW_{l+1}(t) \mW_{l+1}(t)^T 
    &= \mW_l(0)^T \mW_l(0) - \mW_{l+1}(0) \mW_{l+1}(0)^T \\
    &= \alpha^2 \bar \mW_l^T \bar \mW_l - \alpha^2 \bar \mW_{l+1} \bar \mW_{l+1}^T,
\end{aligned}    
\end{equation}
for any $l \in [L-2]$.
Similarly, we have
\begin{equation}
\label{eq:proof-fc-reg-4}
\begin{aligned}
    \mW_{L-1}(t)^T \mW_{L-1}(t) - \vw_{L}(t) \vw_{L}(t)^T 
    &= \mW_{L-1}(0)^T \mW_{L-1}(0) - \vw_{L}(0) \vw_{L}(0)^T\\ 
    &= \alpha^2 \bar \mW_{L-1}^T \bar \mW_{L-1} - \alpha^2 \bar \vw_{L} \bar \vw_{L}^T.
\end{aligned}
\end{equation}
Let us now consider the time derivative of $\mc L(\bThetafc(t))$.
We have the following chain of upper bounds on the time derivative:
\begin{align}
    \frac{d}{dt} \mc L(\bThetafc(t))
    &= \nabla_{\bThetafc} \mc L(\bThetafc(t))^T \dot \bTheta_{\rm fc}(t)
    = -\ltwos{\nabla_{\bThetafc} \mc L(\bThetafc(t))}^2\notag\\
    &\leq -\ltwos{\nabla_{\vw_L} \mc L(\bThetafc(t))}^2
    = -\ltwos{\dot \vw_L(t)}^2\notag\\
    &= -\ltwos{\mW_{L-1}^T \cdots \mW_1^T \mX^T \vr}^2.\label{eq:proof-fc-reg-2}
    %&= -\vr^T \mX \mW_1 \cdots \mW_{L-1} \mW_{L-1}^T \cdots \mW_1^T \mX^T \vr.
\end{align}
Note from \eqref{eq:proof-fc-reg-2} that if $\mW_{L-1}^T \cdots \mW_1^T$ is full-rank, its minimum singular value is positive, and one can bound
\begin{equation}
\label{eq:proof-fc-reg-6}
    \ltwos{\mW_{L-1}^T \cdots \mW_1^T \mX^T \vr} \geq \sigma_{\min}(\mW_{L-1}^T \cdots \mW_1^T) \ltwos{\mX^T \vr}.
\end{equation}
We now prove that the matrix $\mW_{L-1}^T \cdots \mW_1^T$ is full-rank, and its minimum singular value is bounded from below by $\alpha^{L-1} \lambda^{(L-1)/2}$ for any $t \geq 0$. To show this, it suffices to show that
\begin{align}
\label{eq:proof-fc-reg-5}
    \mW_{L-1}^T \cdots \mW_1^T \mW_1 \cdots \mW_{L-1} \succeq \alpha^{2L-2} \lambda^{L-1} \mI_d.
\end{align}
Now,
\begin{align*}
    &~\mW_{L-1}^T \cdots \mW_2^T \mW_1^T \mW_1 \mW_2 \cdots \mW_{L-1}\\
    \eqwtxt{(a)}&~\mW_{L-1}^T \cdots \mW_2^T (\mW_2 \mW_2^T + \alpha^2 \bar \mW_1^T \bar \mW_1 - \alpha^2 \bar \mW_2 \bar \mW_2^T) \mW_2 \cdots \mW_{L-1}\\
    \succeqwtxt{(b)}&~  
    \mW_{L-1}^T \cdots \mW_3^T \mW_2^T \mW_2 \mW_2^T \mW_2 \mW_3 \cdots \mW_{L-1}\\
    \eqwtxt{(a)}&~ \mW_{L-1}^T \cdots \mW_3^T (\mW_3 \mW_3^T + \alpha^2 \bar \mW_2^T \bar \mW_2 - \alpha^2 \bar \mW_3 \bar \mW_3^T )^2 \mW_3 \cdots \mW_{L-1}\\
    \succeqwtxt{(b)}&~
    \mW_{L-1}^T \cdots \mW_3^T (\mW_3 \mW_3^T)^2 \mW_3 \cdots \mW_{L-1}\\
    =&~ \cdots \succeq (\mW_{L-1}^T \mW_{L-1})^{L-1},
\end{align*}
where equalities marked in (a) used \eqref{eq:proof-fc-reg-3}, and inequalities marked in (b) used the initialization conditions $\bar \mW_l^T \bar \mW_l \succeq \bar \mW_{l+1} \bar \mW_{l+1}^T$.
Next, it follows from \eqref{eq:proof-fc-reg-4} that
\begin{align*}
    (\mW_{L-1}^T \mW_{L-1})^{L-1} 
    &= (\vw_L \vw_L^T + \alpha^2 \bar \mW_{L-1}^T \bar \mW_{L-1} - \alpha^2 \bar \vw_L \bar \vw_L^T)^{L-1}\\
    &\succeq \alpha^{2L-2} (\bar \mW_{L-1}^T \bar \mW_{L-1} - \bar \vw_L \bar \vw_L^T)^{L-1}\\
    &\succeqwtxt{(c)} \alpha^{2L-2} \lambda^{L-1} \mI_d.
\end{align*}
where (c) used the assumption that $\bar \mW_{L-1}^T \bar \mW_{L-1} - \bar \vw_L \bar \vw_L^T \succeq \lambda \mI_d$. This proves \eqref{eq:proof-fc-reg-5}. Applying \eqref{eq:proof-fc-reg-5} to \eqref{eq:proof-fc-reg-2} then gives
\begin{align*}
    \frac{d}{dt} \mc L(\bThetafc(t)) 
    &\leq -\ltwos{\mW_{L-1}^T \cdots \mW_1^T \mX^T \vr}^2\\
    &\leq -\sigma_{\min}(\mW_{L-1}^T \cdots \mW_1^T)^2 \ltwos{\mX^T \vr}^2\\
    &\leq -\alpha^{2L-2} \lambda^{L-1} \ltwos{\mX^T \vr}^2\\
    &\leqwtxt{(d)} -\alpha^{2L-2} \lambda^{L-1} \sigma_{\min}(\mX)^2 \ltwos{\vr}^2\\
    &=-\alpha^{2L-2} \lambda^{L-1} \sigma_{\min}(\mX)^2 \mc L(\bThetafc(t)),
\end{align*}
where (d) used the fact that $\mX^T$ is a full column rank matrix to apply a bound similar to \eqref{eq:proof-fc-reg-6}. From this, we get
\begin{equation*}
    \mc L(\bThetafc(t)) \leq \mc L(\bThetafc(0)) \exp(-\alpha^{2L-2} \lambda^{L-1} \sigma_{\min}(\mX)^2 t),
\end{equation*}
hence proving $\mc L(\bThetafc(t)) \to 0$ as $t \to \infty$.


\subsection{Proof of Theorem~\ref{thm:fc-reg}.\ref{item:thm-fc-reg-2}: characterizing the limit point when $\alpha \to 0$}
Now, under the assumption that the initial directions $\bar \mW_1, \dots, \bar \mW_{L-1}, \bar \vw_L$ lead to $\mc L(\bTheta(t)) \to 0$ (which we proved in the previous subsection for certain sufficient conditions), we move on to characterize the limit points of the gradient flow, for the ``active regime'' case $\alpha \to 0$. This part of the proof is motivated from the analysis in \citet{ji2019gradient}.

Let $\vu_l$ and $\vv_l$ be the top left and right singular vectors of $\mW_l$, for $l \in [L-1]$. Note that since $\mW_l$ varies over time, the singular vectors and singular value also vary over time. Similarly, let $s_l$ be the largest singular value of $\mW_l$. We will show that the linear coefficients $\vbetafc(\bThetafc) = \mW_1\cdots\mW_{L-1}\vw_L$ align with $\vu_1$ as $\alpha \to 0$, and $\vu_1$ is in the subspace of $\rowsp(\mX)$ in the limit $\alpha \to 0$, hence proving that $\vbetafc(\bThetafc)$ is the minimum $\ell_2$ norm solution in the limit $\alpha \to 0$.

First, note from \eqref{eq:proof-fc-reg-3} and \eqref{eq:proof-fc-reg-4} that if we take trace of both sides, we get
\begin{align*}
    \nfro{\mW_l}^2 - \nfro{\mW_{l+1}}^2 &= \alpha^2 (\nfro{\bar \mW_l}^2 - \nfro{\bar \mW_{l+1}}^2)~~\text{ for } l \in [L-2],\\
    \nfro{\mW_{L-1}}^2 - \ltwo{\vw_L}^2 &= \alpha^2 (\nfro{\bar \mW_{L-1}}^2 - \ltwo{\bar \vw_L}^2).
\end{align*}
Summing the equations above for $l, l+1, \dots, L-1$, we get
\begin{equation}
\label{eq:proof-fc-reg-7}
    \nfro{\mW_l}^2 - \ltwo{\vw_L}^2 = \alpha^2(\nfro{\bar \mW_l}^2 - \ltwo{\bar \vw_L}^2).
\end{equation}
Next, consider the operator norms (i.e., the maximum singular values), denoted as $\ltwo{\cdot}$, of the matrices.
\begin{align*}
    \ltwo{\mW_l}^2 
    &\geq \vu_{l+1}^T \mW_l^T \mW_l \vu_{l+1}\\
    &\eqwtxt{(e)} \vu_{l+1}^T \mW_{l+1} \mW_{l+1}^T \vu_{l+1}
    + \alpha^2 \vu_{l+1}^T (\bar \mW_{l}^T \bar \mW_{l}- \bar \mW_{l+1} \bar \mW_{l+1}^T) \vu_{l+1}\\
    &= \ltwo{\mW_{l+1}}^2 + \alpha^2 \vu_{l+1}^T (\bar \mW_{l}^T \bar \mW_{l}- \bar \mW_{l+1} \bar \mW_{l+1}^T) \vu_{l+1}\\
    &\geq \ltwo{\mW_{l+1}}^2 - \alpha^2 \ltwos{\bar \mW_{l}^T \bar \mW_{l}- \bar \mW_{l+1} \bar \mW_{l+1}^T}~~\text{ for } l \in [L-2],\\
    \ltwo{\mW_{L-1}}^2
    &\geq \frac{\vw_L}{\ltwo{\vw_L}} \mW_{L-1}^T \mW_{L-1} \frac{\vw_L}{\ltwo{\vw_L}}\\
    &\eqwtxt{(f)} \frac{\vw_L}{\ltwo{\vw_L}} \vw_L \vw_L^T \frac{\vw_L}{\ltwo{\vw_L}} + \alpha^2 \frac{\vw_L}{\ltwo{\vw_L}} (\bar \mW_{L-1}^T \bar \mW_{L-1} - \bar \vw_L \bar \vw_L^T) \frac{\vw_L}{\ltwo{\vw_L}}\\
    &\geq \ltwo{\vw_L}^2 - \alpha^2 \ltwos{\bar \mW_{L-1}^T \bar \mW_{L-1} - \bar \vw_L \bar \vw_L^T}.
\end{align*}
where (e) used \eqref{eq:proof-fc-reg-3} and (f) used \eqref{eq:proof-fc-reg-4}.
Summing the inequalities for $l, l+1, \dots, L-1$ gives
\begin{equation}
	\label{eq:proof-fc-reg-8}
	\begin{aligned}
		\ltwo{\mW_l}^2 \geq \ltwo{\vw_L}^2 
		&- \alpha^2 \sum_{k=l}^{L-2} \ltwos{\bar \mW_{k}^T \bar \mW_{k}- \bar \mW_{k+1} \bar \mW_{k+1}^T}\\
		&- \alpha^2 \ltwos{\bar \mW_{L-1}^T \bar \mW_{L-1}- \bar \vw_{L} \bar \vw_{L}^T}
	\end{aligned}
\end{equation}
From \eqref{eq:proof-fc-reg-7} and \eqref{eq:proof-fc-reg-8}, we get a bound on the gap between the second powers of the Frobenius norm (or the $\ell_2$ norm of singular values) and operator norm (or the maximum singular value $s_l$) of $\mW_l$:
\begin{align}
	\nfro{\mW_l(t)}^2 - \ltwo{\mW_l(t)}^2 
	&\leq 
	\alpha^2(\nfro{\bar \mW_l}^2 - \ltwo{\bar \vw_L}^2) 
	+ \alpha^2 \sum_{k=l}^{L-2} \ltwos{\bar \mW_{k}^T \bar \mW_{k}- \bar \mW_{k+1} \bar \mW_{k+1}^T}\notag\\
	&\quad+ \alpha^2 \ltwos{\bar \mW_{L-1}^T \bar \mW_{L-1}- \bar \vw_{L} \bar \vw_{L}^T},\label{eq:proof-fc-reg-9}
\end{align}
which holds for any $t \geq 0$.
The gap~\eqref{eq:proof-fc-reg-9} implies that each $\mW_l$, for $l \in [L-1]$, can be written as
\begin{equation}
\label{eq:proof-fc-reg-16}
    \mW_l(t) = s_l(t) \vu_l(t) \vv_l(t)^T + O(\alpha^2).
\end{equation}
Next, we show that the ``adjacent'' singular vectors $\vv_l$ and $\vu_{l+1}$ align with each other as $\alpha \to 0$. 
To this end, we will get lower and upper bounds for a quantity $\vv_l^T \mW_{l+1} \mW_{l+1}^T \vv_l$, where $l \in [L-2]$.
\begin{align}
    \vv_l^T \mW_{l+1} \mW_{l+1}^T \vv_l
    &= \vv_l^T \mW_l^T \mW_l \vv_l - \alpha^2 \vv_l^T \bar \mW_l^T \bar \mW_l \vv_l + \alpha^2 \vv_l^T \bar \mW_{l+1} \bar \mW_{l+1}^T \vv_l\notag\\
    &\geq \ltwo{\mW_l}^2 - \alpha^2 \ltwo{\bar \mW_l^T \bar \mW_l -  \bar \mW_{l+1} \bar \mW_{l+1}^T}\notag\\
    &= s_l^2 - \alpha^2 \ltwo{\bar \mW_l^T \bar \mW_l -  \bar \mW_{l+1} \bar \mW_{l+1}^T},\label{eq:proof-fc-reg-10}\\
    \vv_l^T \mW_{l+1} \mW_{l+1}^T \vv_l
    &= \vv_l^T (s_{l+1}^2 \vu_{l+1} \vu_{l+1}^T + \mW_{l+1} \mW_{l+1}^T - s_{l+1}^2 \vu_{l+1} \vu_{l+1}^T) \vv_l\notag\\
    &= s_{l+1}^2 (\vv_l^T \vu_{l+1})^2 + \vv_l^T (\mW_{l+1} \mW_{l+1}^T - s_{l+1}^2 \vu_{l+1} \vu_{l+1}^T) \vv_l\notag\\
    &\leq s_{l+1}^2 (\vv_l^T \vu_{l+1})^2 + \nfro{\mW_{l+1}}^2 - \ltwo{\mW_{l+1}}^2.\label{eq:proof-fc-reg-11}
\end{align}
Combining \eqref{eq:proof-fc-reg-10}, \eqref{eq:proof-fc-reg-11}, and \eqref{eq:proof-fc-reg-9}, we get
\begin{align}
    s_l^2 
    &\leq s_{l+1}^2(\vv_l^T \vu_{l+1})^2 
    + \alpha^2 \ltwo{\bar \mW_l^T \bar \mW_l -  \bar \mW_{l+1} \bar \mW_{l+1}^T}
    + \nfro{\mW_{l+1}}^2 - \ltwo{\mW_{l+1}}^2\notag\\
    &\leq s_{l+1}^2(\vv_l^T \vu_{l+1})^2
    + \alpha^2 (\nfro{\bar \mW_{l+1}}^2 - \ltwo{\bar \vw_L}^2)
    + \alpha^2 \sum_{k=l}^{L-2} \ltwos{\bar \mW_{k}^T \bar \mW_{k}- \bar \mW_{k+1} \bar \mW_{k+1}^T}\notag\\
    &
    \quad+ \alpha^2 \ltwos{\bar \mW_{L-1}^T \bar \mW_{L-1}- \bar \vw_{L} \bar \vw_{L}^T}.
    \label{eq:proof-fc-reg-12}
\end{align}
Next, by a similar reasoning as \eqref{eq:proof-fc-reg-10}, we have
\begin{align}
    s_l^2 \geq \vu_{l+1}^T \mW_{l}^T \mW_{l} \vu_{l+1} \geq s_{l+1}^2 - \alpha^2 \ltwo{\bar \mW_l^T \bar \mW_l -  \bar \mW_{l+1} \bar \mW_{l+1}^T}.\label{eq:proof-fc-reg-13}
\end{align}
Combining \eqref{eq:proof-fc-reg-12} and \eqref{eq:proof-fc-reg-13} and dividing both sides by $s_{l+1}^2$, we get
\begin{equation}
\label{eq:proof-fc-reg-14}
    (\vv_l(t)^T \vu_{l+1}(t))^2 \geq 1 - \alpha^2 
    \frac{G_l}{s_{l+1}(t)^2}
\end{equation}
for $t \geq 0$, where
\begin{align*}
	G_l &\defeq \ltwo{\bar \mW_l^T \bar \mW_l -  \bar \mW_{l+1} \bar \mW_{l+1}^T}
	+ \nfro{\bar \mW_{l+1}}^2 - \ltwo{\bar \vw_L}^2\\
	&\quad + \sum_{k=l}^{L-2} \ltwos{\bar \mW_{k}^T \bar \mW_{k}- \bar \mW_{k+1} \bar \mW_{k+1}^T}
	+\ltwos{\bar \mW_{L-1}^T \bar \mW_{L-1}- \bar \vw_{L} \bar \vw_{L}^T}.
\end{align*}
Applying a similar argument to $l = L-1$, we can also get
\begin{equation}
	\label{eq:proof-fc-reg-15}
	\frac{(\vv_{L-1}(t)^T \vw_L(t))^2}{\ltwo{\vw_L(t)}^2} \geq 1 - \alpha^2 \frac{G_{L-1}}{\ltwo{\vw_L(t)}^2},
\end{equation}
where
%\begin{align*}
$G_{L-1} \defeq 2 \ltwo{\bar \mW_{L-1}^T \bar \mW_{L-1} -  \bar \vw_{L} \bar \vw_{L}^T}$.
%\end{align*}

From \eqref{eq:proof-fc-reg-14} and \eqref{eq:proof-fc-reg-15}, we can note that as $\alpha \to 0$, the inner product between the adjacent singular vectors converges to $\pm 1$, unless $s_{2}, \dots, s_{L-1}, \ltwo{\vw_L}$ also diminish to zero. So it is left to show that the singular values do not diminish to zero as $\alpha \to 0$.
To this end, recall that we assume for this part that
\begin{equation*}
 \lim_{t \to \infty} \mX \mW_1(t) \cdots \mW_{L-1}(t) \vw_L(t) = \vy.   
\end{equation*}
A necessary condition for this to hold is that 
\begin{equation*}
    \frac{\ltwo{\vy}}{\ltwo{\mX}} \leq \lim_{t \to \infty} \ltwo{\mW_1(t) \cdots \mW_{L-1}(t)\vw_L(t)} \leq \lim_{t \to \infty} \prod_{l=1}^{L-1} s_l(t) \ltwo{\vw_L(t)}.
\end{equation*}
This means that after converging to the global minimum solution of the problem (i.e., $t \to \infty$), the product of the singular values must be at least greater than some constant independent of $\alpha$. Moreover, we can see from \eqref{eq:proof-fc-reg-10} and \eqref{eq:proof-fc-reg-13} that the gap between singular values squared of adjacent layers is bounded by $O(\alpha^2)$, for all $t \geq 0$; so the maximum singular values become closer and closer to each other as $\alpha$ diminishes. 
This implies that 
\begin{equation*}
    \lim_{\alpha \to 0} \lim_{t \to \infty} s_l(t) \geq \frac{\ltwo{\vy}^{1/L}}{\ltwo{\mX}^{1/L}}~~\text{ for } l \in [L-1],~~
    \lim_{\alpha \to 0} \lim_{t \to \infty} \ltwo{\vw_L(t)} \geq \frac{\ltwo{\vy}^{1/L}}{\ltwo{\mX}^{1/L}}.
\end{equation*}
Therefore, we have the alignment of singular vectors at convergence as $\alpha \to 0$:
\begin{equation}
\label{eq:proof-fc-reg-17}
    \lim_{\alpha \to 0} \lim_{t \to \infty} (\vv_l(t)^T \vu_{l+1}(t))^2 = 1,~~\text{ for } l \in [L-2],~~
    \lim_{\alpha \to 0} \lim_{t \to \infty} \frac{(\vv_{L-1}(t)^T \vw_L(t))^2}{\ltwo{\vw_L(t)}^2} = 1.
\end{equation}

So far, we saw from \eqref{eq:proof-fc-reg-16} that $\mW_l(t)$'s become rank-1 matrices as $\alpha \to 0$, and from \eqref{eq:proof-fc-reg-17} that the top singular vectors align with each other as $t \to \infty$ and $\alpha \to 0$. These imply that, as $t \to \infty$ and $\alpha \to 0$, $\vbetafc(\bThetafc)$ is a scalar multiple of $\vu_1$, the top left singular vector of $\mW_1$:
\begin{equation}
\label{eq:proof-fc-reg-18}
    \lim_{\alpha \to 0} \lim_{t \to \infty} \vbetafc(\bThetafc(t))
    = c \cdot \lim_{\alpha \to 0} \lim_{t \to \infty} \vu_1(t),
\end{equation}
for some $c \in \reals$.

In light of \eqref{eq:proof-fc-reg-18}, it remains to take a close look at $\vu_1(t)$.
Note from the gradient flow dynamics of $\mW_1$ that $\dot \mW_1$ is always a rank-1 matrix whose columns are in the row space of $\mX$, since $\mX ^T \vr \in \rowsp(\mX)$. This implies that, if we decompose $\mW_1$ into two orthogonal components $\mW_1^\perp$ and $\mW_1^\parallel$ so that the columns in $\mW_1^\parallel$ are in $\rowsp(\mX)$ and the columns in $\mW_1^\perp$ are in the orthogonal subspace $\rowsp(\mX)^\perp$, we have
\begin{equation*}
    \dot \mW_1^\perp = \zeros,~~
    \dot \mW_1^\parallel = \dot \mW_1.
\end{equation*}
That is, any component $\mW_1^\perp(0)$ orthogonal to $\rowsp(\mX)$ remains unchanged for all $t \geq 0$, while the component $\mW_1^\parallel$ changes by the gradient flow.
Since we have 
\begin{equation*}
    \nfro{\mW_1^\perp(t)} = \nfro{\mW_1^\perp(0)} \leq \alpha \nfro{\bar \mW_l},
\end{equation*} 
the component in $\mW_1$ that is orthogonal to $\rowsp(\mX)$ diminishes to zero as $\alpha \to 0$. This means that in the limit $\alpha \to 0$, the columns of $\mW_1$ are entirely from $\rowsp(\mX)$, which also means that
\begin{equation*}
    \lim_{\alpha \to 0} \lim_{t \to \infty} \vbetafc(\bThetafc(t)) \in \rowsp(\mX).
\end{equation*}
However, recall that there is only one unique global minimum $\vz$ satisfying $\mX \vz = \vy$ that is in $\rowsp(\mX)$: namely, $\vz = \mX^T(\mX \mX^T)^{-1} \vy$, the minimum $\ell_2$ norm solution. This finishes the proof.


\end{document}
